\begin{register}{H}{MCPWM\_CLK\_CFG\_REG}{0x0000}\label{regdesc:MCPWMCLKCFGREG}
\regfield{(reserved)}{24}{8}{000000000000000000000000}%
\regfield{MCPWM\_CLK\_PRESCALE}{8}{0}{{0x0}}%
\reglabel{Reset}\regnewline%
\begin{regdesc}\begin{reglist}
\label{fielddesc:MCPWMCLKPRESCALE}\item [MCPWM\_CLK\_PRESCALE] 配置时钟预分频系数，使得 PWM\_CLK 的周期为 6.25ns * (PWM\_CLK\_PRESCALE + 1)。 (R/W)
\end{reglist}\end{regdesc}
\end{register}


\begin{register}{H}{MCPWM\_TIMER0\_CFG0\_REG}{0x0004}\label{regdesc:MCPWMTIMER0CFG0REG}
\regfield{(reserved)}{6}{26}{000000}%
\regfield{MCPWM\_TIMER0\_PERIOD\_UPMETHOD}{2}{24}{{0}}%
\regfield{MCPWM\_TIMER0\_PERIOD}{16}{8}{{0xff}}%
\regfield{MCPWM\_TIMER0\_PRESCALE}{8}{0}{{0x0}}%
\reglabel{Reset}\regnewline%
\begin{regdesc}\begin{reglist}
\label{fielddesc:MCPWMTIMER0PRESCALE}\item [MCPWM\_TIMER0\_PRESCALE] 配置定时器 0 的预分频系数，使得 PT0\_CLK 的周期 = PWM\_CLK 的周期 * (PWM\_TIMER0\_PRESCALE + 1)。 (R/W)
\label{fielddesc:MCPWMTIMER0PERIOD}\item [MCPWM\_TIMER0\_PERIOD] 配置定时器 0 的影子周期。(R/W)
\label{fielddesc:MCPWMTIMER0PERIODUPMETHOD}\item [MCPWM\_TIMER0\_PERIOD\_UPMETHOD] 配置 PWM 定时器 0 周期有效寄存器的更新方式。\\
0：立即更新\\
1：发生 TEZ 事件时更新\\
2：同步时更新\\
3：发生 TEZ 事件或同步时更新\\
本文档中，TEZ 指定时器为 0 时的事件。\\ (R/W)
\end{reglist}\end{regdesc}
\end{register}


\begin{register}{H}{MCPWM\_TIMER0\_CFG1\_REG}{0x0008}\label{regdesc:MCPWMTIMER0CFG1REG}
\regfield{(reserved)}{27}{5}{000000000000000000000000000}%
\regfield{MCPWM\_TIMER0\_MOD}{2}{3}{{0x0}}%
\regfield{MCPWM\_TIMER0\_START}{3}{0}{{0x0}}%
\reglabel{Reset}\regnewline%
\begin{regdesc}\begin{reglist}
\label{fielddesc:MCPWMTIMER0START}\item [MCPWM\_TIMER0\_START] 控制 PWM 定时器的开启与关闭。\\
0：如果开启，在 TEZ 事件发生时停止\\
1：如果开启，在 TEP 事件发生时停止\\
2：开启\\
3：开启，并在下一个 TEZ 事件发生时停止\\
4：开启，并在下一个 TEP 事件发生时停止\\
其他值：无效值，不产生影响\\
本文档中，TEP 指定时器为周期值时发生的事件。\\ (R/W/SC)
\label{fielddesc:MCPWMTIMER0MOD}\item [MCPWM\_TIMER0\_MOD] 配置 PWM 定时器 0 的工作模式。\\
0：暂停\\
1：递增模式\\
2：递减模式\\
3：递增递减循环模式\\
(R/W)
\end{reglist}\end{regdesc}
\end{register}


\begin{register}{H}{MCPWM\_TIMER0\_SYNC\_REG}{0x000C}\label{regdesc:MCPWMTIMER0SYNCREG}
\regfield{(reserved)}{11}{21}{00000000000}%
\regfield{MCPWM\_TIMER0\_PHASE\_DIRECTION}{1}{20}{{0}}%
\regfield{MCPWM\_TIMER0\_PHASE}{16}{4}{{0}}%
\regfield{MCPWM\_TIMER0\_SYNCO\_SEL}{2}{2}{{0}}%
\regfield{MCPWM\_TIMER0\_SYNC\_SW}{1}{1}{{0}}%
\regfield{MCPWM\_TIMER0\_SYNCI\_EN}{1}{0}{{0}}%
\reglabel{Reset}\regnewline%
\begin{regdesc}\begin{reglist}
\label{fielddesc:MCPWMTIMER0SYNCIEN}\item [MCPWM\_TIMER0\_SYNCI\_EN] 置 1 时，使能在同步输入事件发生时的定时器相位重载。 (R/W)
\label{fielddesc:MCPWMTIMER0SYNCSW}\item [MCPWM\_TIMER0\_SYNC\_SW] 此位取反，触发软件同步。 (R/W)
\label{fielddesc:MCPWMTIMER0SYNCOSEL}\item [MCPWM\_TIMER0\_SYNCO\_SEL] 选择 PWM 定时器 0 的同步输出来源。\\
0：同步。取反 \hyperref[fielddesc:MCPWMTIMER0SYNCSW]{MCPWM\_TIMER0\_SYNC\_SW} 位时始终生成同步输出。\\
1：TEZ\\
2：TEP\\
3：无效\\
(R/W)
\label{fielddesc:MCPWMTIMER0PHASE}\item [MCPWM\_TIMER0\_PHASE] 同步事件中定时器重载的相位。 (R/W)
\label{fielddesc:MCPWMTIMER0PHASEDIRECTION}\item [MCPWM\_TIMER0\_PHASE\_DIRECTION] 当定时器 0 为递增-递减循环模式时，配置该定时器方向。\\
0：递增\\
1：递减\\
(R/W)
\end{reglist}\end{regdesc}
\end{register}


\begin{register}{H}{MCPWM\_TIMER0\_STATUS\_REG}{0x0010}\label{regdesc:MCPWMTIMER0STATUSREG}
\regfield{(reserved)}{15}{17}{000000000000000}%
\regfield{MCPWM\_TIMER0\_DIRECTION}{1}{16}{{0}}%
\regfield{MCPWM\_TIMER0\_VALUE}{16}{0}{{0}}%
\reglabel{Reset}\regnewline%
\begin{regdesc}\begin{reglist}
\label{fielddesc:MCPWMTIMER0VALUE}\item [MCPWM\_TIMER0\_VALUE] 表示当前 PWM 定时器 0 计数器的值。 (RO)
\label{fielddesc:MCPWMTIMER0DIRECTION}\item [MCPWM\_TIMER0\_DIRECTION] 配置当前 PWM 定时器 0 的计数器模式。\\
0：递增模式\\
1：递减模式\\
(RO)
\end{reglist}\end{regdesc}
\end{register}


\begin{register}{H}{MCPWM\_TIMER1\_CFG0\_REG}{0x0014}\label{regdesc:MCPWMTIMER1CFG0REG}
\regfield{(reserved)}{6}{26}{000000}%
\regfield{MCPWM\_TIMER1\_PERIOD\_UPMETHOD}{2}{24}{{0}}%
\regfield{MCPWM\_TIMER1\_PERIOD}{16}{8}{{0xff}}%
\regfield{MCPWM\_TIMER1\_PRESCALE}{8}{0}{{0x0}}%
\reglabel{Reset}\regnewline%
\begin{regdesc}\begin{reglist}
\label{fielddesc:MCPWMTIMER1PRESCALE}\item [MCPWM\_TIMER1\_PRESCALE] 配置 timer1 预分频系数，使得 PT0\_CLK 周期 = PWM\_CLK 周期 * (PWM\_timer1\_PRESCALE + 1)。 (R/W)
\label{fielddesc:MCPWMTIMER1PERIOD}\item [MCPWM\_TIMER1\_PERIOD] 定时器 1 的影子周期寄存器。 (R/W)
\label{fielddesc:MCPWMTIMER1PERIODUPMETHOD}\item [MCPWM\_TIMER1\_PERIOD\_UPMETHOD] 配置 PWM 定时器 1 周期有效寄存器的更新方式。\\
0：立即更新\\
1：发生 TEZ 事件时更新\\
2：同步时更新\\
3：发生 TEZ 事件或同步时更新\\
本文档中，TEZ 指定时器为 0 时的事件。\\(R/W)
\end{reglist}\end{regdesc}
\end{register}


\begin{register}{H}{MCPWM\_TIMER1\_CFG1\_REG}{0x0018}\label{regdesc:MCPWMTIMER1CFG1REG}
\regfield{(reserved)}{27}{5}{000000000000000000000000000}%
\regfield{MCPWM\_TIMER1\_MOD}{2}{3}{{0x0}}%
\regfield{MCPWM\_TIMER1\_START}{3}{0}{{0x0}}%
\reglabel{Reset}\regnewline%
\begin{regdesc}\begin{reglist}
\label{fielddesc:MCPWMTIMER1START}\item [MCPWM\_TIMER1\_START] PWM 控制定时器 1 的开启与停止。\\
0：如果开启，在发生 TEZ 事件时停止\\
1：如果开启，在发生 TEP 事件时停止\\
2：开启\\
3：开启并在下一次 TEZ 事件中停止\\
4：开启并在下一次 TEP 事件中停止\\
其他值：无效值，不产生影响。\\
本文档中，TEP 指定时器为周期数时的事件。\\(R/W/SC)
\label{fielddesc:MCPWMTIMER1MOD}\item [MCPWM\_TIMER1\_MOD] PWM 定时器 1 的工作模式。\\
0：暂停\\
1：递增模式\\
2：递减模式\\
3：递增-递减循环模式\\
(R/W)
\end{reglist}\end{regdesc}
\end{register}


\begin{register}{H}{MCPWM\_TIMER1\_SYNC\_REG}{0x001C}\label{regdesc:MCPWMTIMER1SYNCREG}
\regfield{(reserved)}{11}{21}{00000000000}%
\regfield{MCPWM\_TIMER1\_PHASE\_DIRECTION}{1}{20}{{0}}%
\regfield{MCPWM\_TIMER1\_PHASE}{16}{4}{{0}}%
\regfield{MCPWM\_TIMER1\_SYNCO\_SEL}{2}{2}{{0}}%
\regfield{MCPWM\_TIMER1\_SYNC\_SW}{1}{1}{{0}}%
\regfield{MCPWM\_TIMER1\_SYNCI\_EN}{1}{0}{{0}}%
\reglabel{Reset}\regnewline%
\begin{regdesc}\begin{reglist}
\label{fielddesc:MCPWMTIMER1SYNCIEN}\item [MCPWM\_TIMER1\_SYNCI\_EN] 置 1 时，使能在同步输入事件时的定时器相位重载。 (R/W)
\label{fielddesc:MCPWMTIMER1SYNCSW}\item [MCPWM\_TIMER1\_SYNC\_SW] 此位取反，触发软件同步事件。 (R/W)
\label{fielddesc:MCPWMTIMER1SYNCOSEL}\item [MCPWM\_TIMER1\_SYNCO\_SEL] 选择 PWM 定时器 1 的同步输出来源。\\
0：同步\\
1：TEZ\\
2：TEP。取反 reg\_timer1\_sw 位时始终生成同步输出\\
3：无效\\
(R/W)
\label{fielddesc:MCPWMTIMER1PHASE}\item [MCPWM\_TIMER1\_PHASE] 同步时间中定时器重载的相位。(R/W)
\label{fielddesc:MCPWMTIMER1PHASEDIRECTION}\item [MCPWM\_TIMER1\_PHASE\_DIRECTION] 当定时器 1 为递增-递减循环模式时，配置该定时器方向。\\
0：递增\\
1：递减\\
(R/W)
\end{reglist}\end{regdesc}
\end{register}


\begin{register}{H}{MCPWM\_TIMER1\_STATUS\_REG}{0x0020}\label{regdesc:MCPWMTIMER1STATUSREG}
\regfield{(reserved)}{15}{17}{000000000000000}%
\regfield{MCPWM\_TIMER1\_DIRECTION}{1}{16}{{0}}%
\regfield{MCPWM\_TIMER1\_VALUE}{16}{0}{{0}}%
\reglabel{Reset}\regnewline%
\begin{regdesc}\begin{reglist}
\label{fielddesc:MCPWMTIMER1VALUE}\item [MCPWM\_TIMER1\_VALUE] 配置当前 PWM 定时器 1 的计数器值。 (RO)
\label{fielddesc:MCPWMTIMER1DIRECTION}\item [MCPWM\_TIMER1\_DIRECTION] 配置当前 PWM 定时器 1 的计数器模式。\\
0：递增\\
1：递减\\
(RO)
\end{reglist}\end{regdesc}
\end{register}


\begin{register}{H}{MCPWM\_TIMER2\_CFG0\_REG}{0x0024}\label{regdesc:MCPWMTIMER2CFG0REG}
\regfield{(reserved)}{6}{26}{000000}%
\regfield{MCPWM\_TIMER2\_PERIOD\_UPMETHOD}{2}{24}{{0}}%
\regfield{MCPWM\_TIMER2\_PERIOD}{16}{8}{{0xff}}%
\regfield{MCPWM\_TIMER2\_PRESCALE}{8}{0}{{0x0}}%
\reglabel{Reset}\regnewline%
\begin{regdesc}\begin{reglist}
\label{fielddesc:MCPWMTIMER2PRESCALE}\item [MCPWM\_TIMER2\_PRESCALE] 配置定时器 2 的预分频系数，使得 PT0\_CLK 周期 = PWM\_CLK 周期 * (PWM\_timer2\_PRESCALE + 1)。 (R/W)
\label{fielddesc:MCPWMTIMER2PERIOD}\item [MCPWM\_TIMER2\_PERIOD] PWM 定时器 2 的影子周期寄存器。 (R/W)
\label{fielddesc:MCPWMTIMER2PERIODUPMETHOD}\item [MCPWM\_TIMER2\_PERIOD\_UPMETHOD] 配置PWM 定时器 2 周期有效寄存器的更新方式。\\0：立即更新\\1：发生 TEZ 事件时更新\\2：同步时更新\\3：发生 TEZ 事件或同步时更新\\本文档中，TEZ 指定时器为 0 时的事件。\\(R/W)
\end{reglist}\end{regdesc}
\end{register}


\begin{register}{H}{MCPWM\_TIMER2\_CFG1\_REG}{0x0028}\label{regdesc:MCPWMTIMER2CFG1REG}
\regfield{(reserved)}{27}{5}{000000000000000000000000000}%
\regfield{MCPWM\_TIMER2\_MOD}{2}{3}{{0x0}}%
\regfield{MCPWM\_TIMER2\_START}{3}{0}{{0x0}}%
\reglabel{Reset}\regnewline%
\begin{regdesc}\begin{reglist}
\label{fielddesc:MCPWMTIMER2START}\item [MCPWM\_TIMER2\_START] 配置 PWM 控制定时器 2 的开启与停止。\\0：如果开启，在发生 TEZ 事件时停止\\1：如果开启，在发生 TEP 事件时停止\\2：开启\\3：开启并在下一次 TEZ 事件中停止\\4：开启并在下一次 TEP 事件中停止\\其他值：无效值，不产生影响\\
本文档中，TEP 指定时器为周期数时的事件。\\(R/W/SC)
\label{fielddesc:MCPWMTIMER2MOD}\item [MCPWM\_TIMER2\_MOD] 配置 PWM 定时器 1 的工作模式。\\0：暂停\\1：递增模式\\2：递减模式\\3：递增-递减循环模式 \\(R/W)
\end{reglist}\end{regdesc}
\end{register}


\begin{register}{H}{MCPWM\_TIMER2\_SYNC\_REG}{0x002C}\label{regdesc:MCPWMTIMER2SYNCREG}
\regfield{(reserved)}{11}{21}{00000000000}%
\regfield{MCPWM\_TIMER2\_PHASE\_DIRECTION}{1}{20}{{0}}%
\regfield{MCPWM\_TIMER2\_PHASE}{16}{4}{{0}}%
\regfield{MCPWM\_TIMER2\_SYNCO\_SEL}{2}{2}{{0}}%
\regfield{MCPWM\_TIMER2\_SYNC\_SW}{1}{1}{{0}}%
\regfield{MCPWM\_TIMER2\_SYNCI\_EN}{1}{0}{{0}}%
\reglabel{Reset}\regnewline%
\begin{regdesc}\begin{reglist}
\label{fielddesc:MCPWMTIMER2SYNCIEN}\item [MCPWM\_TIMER2\_SYNCI\_EN] 置 1 时使能在同步输入事件时的定时器相位重载。 (R/W)
\label{fielddesc:MCPWMTIMER2SYNCSW}\item [MCPWM\_TIMER2\_SYNC\_SW] 此位取反，触发软件同步事件。 (R/W)
\label{fielddesc:MCPWMTIMER2SYNCOSEL}\item [MCPWM\_TIMER2\_SYNCO\_SEL] 选择 PWM 定时器 2 的同步输出来源。\\0：同步\\1：TEZ\\2：TEP。取反 reg\_timer2\_sw 位时始终生成同步输出\\3：无效\\ (R/W)
\label{fielddesc:MCPWMTIMER2PHASE}\item [MCPWM\_TIMER2\_PHASE] 同步事件中定时器重载相位。 (R/W)
\label{fielddesc:MCPWMTIMER2PHASEDIRECTION}\item [MCPWM\_TIMER2\_PHASE\_DIRECTION] 当定时器 2 为递增-递减循环模式时，配置该定时器方向。\\
0：递增\\
1：递减\\
(R/W)
\end{reglist}\end{regdesc}
\end{register}


\begin{register}{H}{MCPWM\_TIMER2\_STATUS\_REG}{0x0030}\label{regdesc:MCPWMTIMER2STATUSREG}
\regfield{(reserved)}{15}{17}{000000000000000}%
\regfield{MCPWM\_TIMER2\_DIRECTION}{1}{16}{{0}}%
\regfield{MCPWM\_TIMER2\_VALUE}{16}{0}{{0}}%
\reglabel{Reset}\regnewline%
\begin{regdesc}\begin{reglist}
\label{fielddesc:MCPWMTIMER2VALUE}\item [MCPWM\_TIMER2\_VALUE] 表示当前 PWM 定时器 2 计数器的值。 (RO)
\label{fielddesc:MCPWMTIMER2DIRECTION}\item [MCPWM\_TIMER2\_DIRECTION] 表示当前 PWM 定时器 2 的计数器模式。\\
0：递增\\
1：递减\\
(RO)
\end{reglist}\end{regdesc}
\end{register}


\begin{register}{H}{MCPWM\_TIMER\_SYNCI\_CFG\_REG}{0x0034}\label{regdesc:MCPWMTIMERSYNCICFGREG}
\regfield{(reserved)}{20}{12}{00000000000000000000}%
\regfield{MCPWM\_EXTERNAL\_SYNCI2\_INVERT}{1}{11}{{0}}%
\regfield{MCPWM\_EXTERNAL\_SYNCI1\_INVERT}{1}{10}{{0}}%
\regfield{MCPWM\_EXTERNAL\_SYNCI0\_INVERT}{1}{9}{{0}}%
\regfield{MCPWM\_TIMER2\_SYNCISEL}{3}{6}{{0}}%
\regfield{MCPWM\_TIMER1\_SYNCISEL}{3}{3}{{0}}%
\regfield{MCPWM\_TIMER0\_SYNCISEL}{3}{0}{{0}}%
\reglabel{Reset}\regnewline%
\begin{regdesc}\begin{reglist}
\label{fielddesc:MCPWMTIMER0SYNCISEL}\item [MCPWM\_TIMER0\_SYNCISEL] 选择 PWM 定时器 0 的同步输入来源。 \\
1：PWM 定时器 0 同步输出\\
2：PWM 定时器 1 同步输出\\
3：PWM 定时器 2 同步输出\\
4：来自 GPIO 矩阵的 SYNC0\\
5：来自 GPIO 矩阵的 SYNC1\\
6：来自 GPIO 矩阵的 SYNC2\\
其他值：未选择任何同步输入\\
(R/W)
\label{fielddesc:MCPWMTIMER1SYNCISEL}\item [MCPWM\_TIMER1\_SYNCISEL] 选择 PWM 定时器 1 的同步输入来源。 \\
1：PWM 定时器 0 同步输出\\
2：PWM 定时器 1 同步输出\\
3：PWM 定时器 2 同步输出\\
4：来自 GPIO 矩阵的 SYNC0\\
5：来自 GPIO 矩阵的 SYNC1\\
6：来自 GPIO 矩阵的 SYNC2\\
其他值：未选择任何同步输入\\
(R/W)
\label{fielddesc:MCPWMTIMER2SYNCISEL}\item [MCPWM\_TIMER2\_SYNCISEL] 选择 PWM 定时器 2 的同步输入来源。 \\
1：PWM 定时器 0 同步输出\\
2：PWM 定时器 1 同步输出\\
3：PWM 定时器 2 同步输出\\
4：来自 GPIO 矩阵的 SYNC0\\
5：来自 GPIO 矩阵的 SYNC1\\
6：来自 GPIO 矩阵的 SYNC2\\
其他值：未选择任何同步输入\\
(R/W)
\label{fielddesc:MCPWMEXTERNALSYNCI0INVERT}\item [MCPWM\_EXTERNAL\_SYNCI0\_INVERT] 将来自 GPIO 矩阵的 SYNC0 反相。 (R/W)
\label{fielddesc:MCPWMEXTERNALSYNCI1INVERT}\item [MCPWM\_EXTERNAL\_SYNCI1\_INVERT] 将来自 GPIO 矩阵的 SYNC1 反相。 (R/W)
\label{fielddesc:MCPWMEXTERNALSYNCI2INVERT}\item [MCPWM\_EXTERNAL\_SYNCI2\_INVERT] 将来自 GPIO 矩阵的 SYNC2 反相。(R/W)
\end{reglist}\end{regdesc}
\end{register}


\begin{register}{H}{MCPWM\_OPERATOR\_TIMERSEL\_REG}{0x0038}\label{regdesc:MCPWMOPERATORTIMERSELREG}
\regfield{(reserved)}{26}{6}{00000000000000000000000000}%
\regfield{MCPWM\_OPERATOR2\_TIMERSEL}{2}{4}{{0}}%
\regfield{MCPWM\_OPERATOR1\_TIMERSEL}{2}{2}{{0}}%
\regfield{MCPWM\_OPERATOR0\_TIMERSEL}{2}{0}{{0}}%
\reglabel{Reset}\regnewline%
\begin{regdesc}\begin{reglist}
\label{fielddesc:MCPWMOPERATOR0TIMERSEL}\item [MCPWM\_OPERATOR0\_TIMERSEL] 选择 PWM 操作器 0 的定时参考来源。\\0：定时器 0\\1：定时器 1\\2：定时器 2\\3：无效\\ (R/W)
\label{fielddesc:MCPWMOPERATOR1TIMERSEL}\item [MCPWM\_OPERATOR1\_TIMERSEL] 选择 PWM 操作器 1 的定时参考来源。\\0：定时器 0\\1：定时器 1\\2：定时器 2\\3：无效\\ (R/W)
\label{fielddesc:MCPWMOPERATOR2TIMERSEL}\item [MCPWM\_OPERATOR2\_TIMERSEL] 选择 PWM 操作器 2 的定时参考来源。\\0：定时器 0\\1：定时器 1\\2：定时器 2\\3：无效\\ (R/W)
\end{reglist}\end{regdesc}
\end{register}


\begin{register}{H}{MCPWM\_GEN0\_STMP\_CFG\_REG}{0x003C}\label{regdesc:MCPWMGEN0STMPCFGREG}
\regfield{(reserved)}{22}{10}{0000000000000000000000}%
\regfield{MCPWM\_GEN0\_B\_SHDW\_FULL}{1}{9}{{0}}%
\regfield{MCPWM\_GEN0\_A\_SHDW\_FULL}{1}{8}{{0}}%
\regfield{MCPWM\_GEN0\_B\_UPMETHOD}{4}{4}{{0}}%
\regfield{MCPWM\_GEN0\_A\_UPMETHOD}{4}{0}{{0}}%
\reglabel{Reset}\regnewline%
\begin{regdesc}\begin{reglist}
\label{fielddesc:MCPWMGEN0AUPMETHOD}\item [MCPWM\_GEN0\_A\_UPMETHOD] 配置 PWM 生成器 0 时间戳寄存器 A 的有效寄存器的更新方式。\\
所有位值为 0：立即更新\\
bit0 为 1：发生 TEZ 事件时更新\\
bit1 为 1：发生 TEP 事件时更新\\
bit2 为 1：发生同步时间时更新\\
bit3 为 1：关闭更新\\
(R/W)
\label{fielddesc:MCPWMGEN0BUPMETHOD}\item [MCPWM\_GEN0\_B\_UPMETHOD] 配置 PWM 生成器 0 时间戳寄存器 B 有效寄存器的更新方式。\\
所有位值为 0：立即更新\\
bit0 为 1：发生 TEZ 事件时更新\\
bit1 为 1：发生 TEP 事件时更新\\
bit2 为 1：发生同步时间时更新\\
bit3 为 1：关闭更新\\
(R/W)
\label{fielddesc:MCPWMGEN0ASHDWFULL}\item [MCPWM\_GEN0\_A\_SHDW\_FULL] 配置是否在某个特定时间将影子寄存器中的值写入对应的有效寄存器中。由硬件置 1 或复位。\\
0：将影子寄存器中最新的值写入有效寄存器 A 中\\
1：将值写入 PWM 生成器 0 时间戳寄存器 A 的影子寄存器中，等待传输给有效寄存器 A。\\
(R/SC/WTC)
\label{fielddesc:MCPWMGEN0BSHDWFULL}\item [MCPWM\_GEN0\_B\_SHDW\_FULL] 配置是否在某个特定时间将影子寄存器中的值写入对应的有效寄存器中。由硬件置 1 或复位。\\
0：将影子寄存器中最新的值写入有效寄存器 B 中\\
1：将值写入 PWM 生成器 0 时间戳寄存器 B 的影子寄存器中，等待传输给有效寄存器 B。\\
(R/SC/WTC)
\end{reglist}\end{regdesc}
\end{register}


\begin{register}{H}{MCPWM\_GEN0\_TSTMP\_A\_REG}{0x0040}\label{regdesc:MCPWMGEN0TSTMPAREG}
\regfield{(reserved)}{16}{16}{0000000000000000}%
\regfield{MCPWM\_GEN0\_A}{16}{0}{{0}}%
\reglabel{Reset}\regnewline%
\begin{regdesc}\begin{reglist}
\label{fielddesc:MCPWMGEN0A}\item [MCPWM\_GEN0\_A] PWM 生成器 0 时间戳寄存器 A 的影子寄存器。 (R/W)
\end{reglist}\end{regdesc}
\end{register}


\begin{register}{H}{MCPWM\_GEN0\_TSTMP\_B\_REG}{0x0044}\label{regdesc:MCPWMGEN0TSTMPBREG}
\regfield{(reserved)}{16}{16}{0000000000000000}%
\regfield{MCPWM\_GEN0\_B}{16}{0}{{0}}%
\reglabel{Reset}\regnewline%
\begin{regdesc}\begin{reglist}
\label{fielddesc:MCPWMGEN0B}\item [MCPWM\_GEN0\_B] PWM 生成器 0 时间戳寄存器 B 的影子寄存器。 (R/W)
\end{reglist}\end{regdesc}
\end{register}


\begin{register}{H}{MCPWM\_GEN0\_CFG0\_REG}{0x0048}\label{regdesc:MCPWMGEN0CFG0REG}
\regfield{(reserved)}{22}{10}{0000000000000000000000}%
\regfield{MCPWM\_GEN0\_T1\_SEL}{3}{7}{{0}}%
\regfield{MCPWM\_GEN0\_T0\_SEL}{3}{4}{{0}}%
\regfield{MCPWM\_GEN0\_CFG\_UPMETHOD}{4}{0}{{0}}%
\reglabel{Reset}\regnewline%
\begin{regdesc}\begin{reglist}
\label{fielddesc:MCPWMGEN0CFGUPMETHOD}\item [MCPWM\_GEN0\_CFG\_UPMETHOD] PWM 生成器 0 有效配置寄存器的更新方式。\\所有 bit 值为 0：立即更新\\bit0 为 1：发生 TEZ 事件时更新\\bit1 为 1：发生 TEP 事件时更新\\bit2 为 1：发生同 步时间时更新\\bit3 为 1：关闭更新\\ (R/W)
\label{fielddesc:MCPWMGEN0T0SEL}\item [MCPWM\_GEN0\_T0\_SEL] 选择 PWM 生成器 0 event\_t0 的信号源，立即生效。\\0：fault\_event0\\1：fault\_event1\\2：fault\_event2\\3：sync\_taken\\4：无\\(R/W)
\label{fielddesc:MCPWMGEN0T1SEL}\item [MCPWM\_GEN0\_T1\_SEL] 选择 PWM 生成器 0 event\_t1 的信号源，立即生效。\\0：fault\_event0\\1：fault\_event1\\2：fault\_event2\\3：sync\_taken\\4：无\\ (R/W)
\end{reglist}\end{regdesc}
\end{register}


\begin{register}{H}{MCPWM\_GEN0\_FORCE\_REG}{0x004C}\label{regdesc:MCPWMGEN0FORCEREG}
\regfield{(reserved)}{16}{16}{0000000000000000}%
\regfield{MCPWM\_GEN0\_B\_NCIFORCE\_MODE}{2}{14}{{0}}%
\regfield{MCPWM\_GEN0\_B\_NCIFORCE}{1}{13}{{0}}%
\regfield{MCPWM\_GEN0\_A\_NCIFORCE\_MODE}{2}{11}{{0}}%
\regfield{MCPWM\_GEN0\_A\_NCIFORCE}{1}{10}{{0}}%
\regfield{MCPWM\_GEN0\_B\_CNTUFORCE\_MODE}{2}{8}{{0}}%
\regfield{MCPWM\_GEN0\_A\_CNTUFORCE\_MODE}{2}{6}{{0}}%
\regfield{MCPWM\_GEN0\_CNTUFORCE\_UPMETHOD}{6}{0}{{0x20}}%
\reglabel{Reset}\regnewline%
\begin{regdesc}\begin{reglist}
\label{fielddesc:MCPWMGEN0CNTUFORCEUPMETHOD}\item [MCPWM\_GEN0\_CNTUFORCE\_UPMETHOD] 生成器 0 的连续软件强制事件更新方式。\\所有 bit 为 0 时：立即更新\\bit0 为 1：发生 TEZ 事件时更新\\bit1 为 1：发生 TEP 事件时更新\\bit2 为 1：发生 TEA 事件时更新\\bit3 为 1：发生 TEB 事件时更新\\bit4：发生同步事件时更新\\bit5 为 1：关闭更新\\本文档中，TEA/B 指定时器值为寄存器 A/B 的值时生成的事件。\\(R/W)
\label{fielddesc:MCPWMGEN0ACNTUFORCEMODE}\item [MCPWM\_GEN0\_A\_CNTUFORCE\_MODE] 设置 PWM0A 的连续软件强制模式。\\0：关闭\\1：低电平\\2：高电平\\3：关闭\\ (R/W)
\label{fielddesc:MCPWMGEN0BCNTUFORCEMODE}\item [MCPWM\_GEN0\_B\_CNTUFORCE\_MODE] 设置 PWM0B 的连续软件强制模式。\\0：关闭\\1：低电平\\2：高电平\\3：关闭\\ (R/W)
\label{fielddesc:MCPWMGEN0ANCIFORCE}\item [MCPWM\_GEN0\_A\_NCIFORCE] 该位的值取反时将触发 PWM0A 上的非连续即时软件强制事件。(R/W)
\label{fielddesc:MCPWMGEN0ANCIFORCEMODE}\item [MCPWM\_GEN0\_A\_NCIFORCE\_MODE] 设置用于 PWM0A 的非连续即时软件强制模式。\\0：关闭\\1：低电平\\2：高电平\\3：关闭\\(R/W)
\item [见下页...]
\end{reglist}\end{regdesc}
\end{register}
\addtocounter{Regfloat}{-1}
\begin{register}{H}{MCPWM\_GEN0\_FORCE\_REG}{0x0034}
\begin{regdesc}\begin{reglist}
\item [...接上页]
\label{fielddesc:MCPWMGEN0BNCIFORCE}\item [MCPWM\_GEN0\_B\_NCIFORCE] 该位的值取反时将触发PWM0B上的非连续即时软件强制事件。(R/W)
\label{fielddesc:MCPWMGEN0BNCIFORCEMODE}\item [MCPWM\_GEN0\_B\_NCIFORCE\_MODE] 设置用于 PWM0B 的非连续即时软件强制模式。\\0：关闭\\1：低电平\\2：高电平\\3：关闭\\ (R/W)
\end{reglist}\end{regdesc}
\end{register}


\begin{register}{H}{MCPWM\_GEN0\_A\_REG}{0x0050}\label{regdesc:MCPWMGEN0AREG}
\regfield{(reserved)}{8}{24}{00000000}%
\regfield{MCPWM\_GEN0\_A\_DT1}{2}{22}{{0}}%
\regfield{MCPWM\_GEN0\_A\_DT0}{2}{20}{{0}}%
\regfield{MCPWM\_GEN0\_A\_DTEB}{2}{18}{{0}}%
\regfield{MCPWM\_GEN0\_A\_DTEA}{2}{16}{{0}}%
\regfield{MCPWM\_GEN0\_A\_DTEP}{2}{14}{{0}}%
\regfield{MCPWM\_GEN0\_A\_DTEZ}{2}{12}{{0}}%
\regfield{MCPWM\_GEN0\_A\_UT1}{2}{10}{{0}}%
\regfield{MCPWM\_GEN0\_A\_UT0}{2}{8}{{0}}%
\regfield{MCPWM\_GEN0\_A\_UTEB}{2}{6}{{0}}%
\regfield{MCPWM\_GEN0\_A\_UTEA}{2}{4}{{0}}%
\regfield{MCPWM\_GEN0\_A\_UTEP}{2}{2}{{0}}%
\regfield{MCPWM\_GEN0\_A\_UTEZ}{2}{0}{{0}}%
\reglabel{Reset}\regnewline%
\begin{regdesc}\begin{reglist}
\label{fielddesc:MCPWMGEN0AUTEZ}\item [MCPWM\_GEN0\_A\_UTEZ] 定时器递增时，TEZ 事件在 PWM0A 上触发的操作。\\0：波形无改变\\1：拉低\\2：拉高\\3：取反\\ (R/W)
\label{fielddesc:MCPWMGEN0AUTEP}\item [MCPWM\_GEN0\_A\_UTEP] 定时器递增时，TEP 事件在 PWM0A 上触发的操作。具体配置值请参考 \hyperref[fielddesc:MCPWMGEN0AUTEZ]{MCPWM\_GEN0\_A\_UTEZ}。 (R/W)
\label{fielddesc:MCPWMGEN0AUTEA}\item [MCPWM\_GEN0\_A\_UTEA] 定时器递增时，TEA 事件在 PWM0A 上触发的操作。具体配置值请参考 \hyperref[fielddesc:MCPWMGEN0AUTEZ]{MCPWM\_GEN0\_A\_UTEZ}。(R/W)
\label{fielddesc:MCPWMGEN0AUTEB}\item [MCPWM\_GEN0\_A\_UTEB] 定时器递增时，TEB 事件在 PWM0A 上触发的操作。具体配置值请参考 \hyperref[fielddesc:MCPWMGEN0AUTEZ]{MCPWM\_GEN0\_A\_UTEZ}。 (R/W)
\label{fielddesc:MCPWMGEN0AUT0}\item [MCPWM\_GEN0\_A\_UT0] 定时器递增时，event\_t0 在 PWM0A 上触发的操作。具体配置值请参考 \hyperref[fielddesc:MCPWMGEN0AUTEZ]{MCPWM\_GEN0\_A\_UTEZ}。 (R/W)
\label{fielddesc:MCPWMGEN0AUT1}\item [MCPWM\_GEN0\_A\_UT1] 定时器递增时，event\_t1 在 PWM0A 上触发的操作。具体配置值请参考 \hyperref[fielddesc:MCPWMGEN0AUTEZ]{MCPWM\_GEN0\_A\_UTEZ}。 (R/W)
\label{fielddesc:MCPWMGEN0ADTEZ}\item [MCPWM\_GEN0\_A\_DTEZ] 定时器递减时，TEZ 事件在 PWM0A 上触发的操作。具体配置值请参考 \hyperref[fielddesc:MCPWMGEN0AUTEZ]{MCPWM\_GEN0\_A\_UTEZ}。 (R/W)
\label{fielddesc:MCPWMGEN0ADTEP}\item [MCPWM\_GEN0\_A\_DTEP] 定时器递减时，TEP 事件在 PWM0A 上触发的操作。具体配置值请参考 \hyperref[fielddesc:MCPWMGEN0AUTEZ]{MCPWM\_GEN0\_A\_UTEZ}。 (R/W)
\label{fielddesc:MCPWMGEN0ADTEA}\item [MCPWM\_GEN0\_A\_DTEA] 定时器递减时，TEA 事件在 PWM0A 上触发的操作。具体配置值请参考 \hyperref[fielddesc:MCPWMGEN0AUTEZ]{MCPWM\_GEN0\_A\_UTEZ}。 (R/W)
\label{fielddesc:MCPWMGEN0ADTEB}\item [MCPWM\_GEN0\_A\_DTEB] 定时器递减时，TEB 事件在 PWM0A 上触发的操作。具体配置值请参考 \hyperref[fielddesc:MCPWMGEN0AUTEZ]{MCPWM\_GEN0\_A\_UTEZ}。 (R/W)
\label{fielddesc:MCPWMGEN0ADT0}\item [MCPWM\_GEN0\_A\_DT0] 定时器递减时，event\_t0 在 PWM0A 上触发的操作。具体配置值请参考 \hyperref[fielddesc:MCPWMGEN0AUTEZ]{MCPWM\_GEN0\_A\_UTEZ}。 (R/W)
\label{fielddesc:MCPWMGEN0ADT1}\item [MCPWM\_GEN0\_A\_DT1] 定时器递减时，event\_t1 在 PWM0A 上触发的操作。具体配置值请参考 \hyperref[fielddesc:MCPWMGEN0AUTEZ]{MCPWM\_GEN0\_A\_UTEZ}。 (R/W)
\end{reglist}\end{regdesc}
\end{register}


\begin{register}{H}{MCPWM\_GEN0\_B\_REG}{0x0054}\label{regdesc:MCPWMGEN0BREG}
\regfield{(reserved)}{8}{24}{00000000}%
\regfield{MCPWM\_GEN0\_B\_DT1}{2}{22}{{0}}%
\regfield{MCPWM\_GEN0\_B\_DT0}{2}{20}{{0}}%
\regfield{MCPWM\_GEN0\_B\_DTEB}{2}{18}{{0}}%
\regfield{MCPWM\_GEN0\_B\_DTEA}{2}{16}{{0}}%
\regfield{MCPWM\_GEN0\_B\_DTEP}{2}{14}{{0}}%
\regfield{MCPWM\_GEN0\_B\_DTEZ}{2}{12}{{0}}%
\regfield{MCPWM\_GEN0\_B\_UT1}{2}{10}{{0}}%
\regfield{MCPWM\_GEN0\_B\_UT0}{2}{8}{{0}}%
\regfield{MCPWM\_GEN0\_B\_UTEB}{2}{6}{{0}}%
\regfield{MCPWM\_GEN0\_B\_UTEA}{2}{4}{{0}}%
\regfield{MCPWM\_GEN0\_B\_UTEP}{2}{2}{{0}}%
\regfield{MCPWM\_GEN0\_B\_UTEZ}{2}{0}{{0}}%
\reglabel{Reset}\regnewline%
\begin{regdesc}\begin{reglist}
\label{fielddesc:MCPWMGEN0BUTEZ}\item [MCPWM\_GEN0\_B\_UTEZ] 定时器递增时，TEZ 在 PWM0B 上触发的操作。\\0：波形无改变\\1：拉低\\2：拉高\\3：取反\\ (R/W)
\label{fielddesc:MCPWMGEN0BUTEP}\item [MCPWM\_GEN0\_B\_UTEP] 定时器递增时，TEP 在 PWM0B 上触发的操作。具体配置值请参考 \hyperref[fielddesc:MCPWMGEN0BUTEZ]{MCPWM\_GEN0\_B\_UTEZ}。 (R/W)
\label{fielddesc:MCPWMGEN0BUTEA}\item [MCPWM\_GEN0\_B\_UTEA] 定时器递增时，TEA 在 PWM0B 上触发的操作。具体配置值请参考 \hyperref[fielddesc:MCPWMGEN0BUTEZ]{MCPWM\_GEN0\_B\_UTEZ}。 (R/W)
\label{fielddesc:MCPWMGEN0BUTEB}\item [MCPWM\_GEN0\_B\_UTEB] 定时器递增时，TEB 在 PWM0B 上触发的操作。具体配置值请参考 \hyperref[fielddesc:MCPWMGEN0BUTEZ]{MCPWM\_GEN0\_B\_UTEZ}。 (R/W)
\label{fielddesc:MCPWMGEN0BUT0}\item [MCPWM\_GEN0\_B\_UT0] 定时器递增时，event\_t0 在 PWM0B 上触发的操作。具体配置值请参考 \hyperref[fielddesc:MCPWMGEN0BUTEZ]{MCPWM\_GEN0\_B\_UTEZ}。 (R/W)
\label{fielddesc:MCPWMGEN0BUT1}\item [MCPWM\_GEN0\_B\_UT1] 定时器递增时，event\_t1 在 PWM0B 上触发的操作。具体配置值请参考 \hyperref[fielddesc:MCPWMGEN0BUTEZ]{MCPWM\_GEN0\_B\_UTEZ}。 (R/W)
\label{fielddesc:MCPWMGEN0BDTEZ}\item [MCPWM\_GEN0\_B\_DTEZ] 定时器递减时，TEZ 事件在 PWM0B 上触发的操作。具体配置值请参考 \hyperref[fielddesc:MCPWMGEN0BUTEZ]{MCPWM\_GEN0\_B\_UTEZ}。 (R/W)
\label{fielddesc:MCPWMGEN0BDTEP}\item [MCPWM\_GEN0\_B\_DTEP] 定时器递减时，TEP 事件在 PWM0B 上触发的操作。具体配置值请参考 \hyperref[fielddesc:MCPWMGEN0BUTEZ]{MCPWM\_GEN0\_B\_UTEZ}。 (R/W)
\label{fielddesc:MCPWMGEN0BDTEA}\item [MCPWM\_GEN0\_B\_DTEA] 定时器递减时，TEA 事件在 PWM0B 上触发的操作。具体配置值请参考 \hyperref[fielddesc:MCPWMGEN0BUTEZ]{MCPWM\_GEN0\_B\_UTEZ}。 (R/W)
\label{fielddesc:MCPWMGEN0BDTEB}\item [MCPWM\_GEN0\_B\_DTEB] 定时器递减时，TEB 事件在 PWM0B 上触发的操作。具体配置值请参考 \hyperref[fielddesc:MCPWMGEN0BUTEZ]{MCPWM\_GEN0\_B\_UTEZ}。 (R/W)
\label{fielddesc:MCPWMGEN0BDT0}\item [MCPWM\_GEN0\_B\_DT0] 定时器递减时，event\_t0 事件在 PWM0B 上触发的操作。具体配置值请参考 \hyperref[fielddesc:MCPWMGEN0BUTEZ]{MCPWM\_GEN0\_B\_UTEZ}。 (R/W)
\label{fielddesc:MCPWMGEN0BDT1}\item [MCPWM\_GEN0\_B\_DT1] 定时器递减时，event\_t1 在 PWM0B 上触发的操作。具体配置值请参考 \hyperref[fielddesc:MCPWMGEN0BUTEZ]{MCPWM\_GEN0\_B\_UTEZ}。 (R/W)
\end{reglist}\end{regdesc}
\end{register}


\begin{register}{H}{MCPWM\_DT0\_CFG\_REG}{0x0058}\label{regdesc:MCPWMDT0CFGREG}
\regfield{(reserved)}{14}{18}{00000000000000}%
\regfield{MCPWM\_DT0\_CLK\_SEL}{1}{17}{{0}}%
\regfield{MCPWM\_DT0\_B\_OUTBYPASS}{1}{16}{{1}}%
\regfield{MCPWM\_DT0\_A\_OUTBYPASS}{1}{15}{{1}}%
\regfield{MCPWM\_DT0\_FED\_OUTINVERT}{1}{14}{{0}}%
\regfield{MCPWM\_DT0\_RED\_OUTINVERT}{1}{13}{{0}}%
\regfield{MCPWM\_DT0\_FED\_INSEL}{1}{12}{{0}}%
\regfield{MCPWM\_DT0\_RED\_INSEL}{1}{11}{{0}}%
\regfield{MCPWM\_DT0\_B\_OUTSWAP}{1}{10}{{0}}%
\regfield{MCPWM\_DT0\_A\_OUTSWAP}{1}{9}{{0}}%
\regfield{MCPWM\_DT0\_DEB\_MODE}{1}{8}{{0}}%
\regfield{MCPWM\_DT0\_RED\_UPMETHOD}{4}{4}{{0}}%
\regfield{MCPWM\_DT0\_FED\_UPMETHOD}{4}{0}{{0}}%
\reglabel{Reset}\regnewline%
\begin{regdesc}\begin{reglist}
\label{fielddesc:MCPWMDT0FEDUPMETHOD}\item [MCPWM\_DT0\_FED\_UPMETHOD] 下降沿延迟有效寄存器的更新方式。\\
所有bit值为0：立即更新\\
bit0 为 1：发生 TEZ 事件时更新\\
bit1 为 1：发生 TEP 事件时更新\\
bit2 为 1：发生同步时间时更新\\
bit3 为 1：关闭更新\\
(R/W)
\label{fielddesc:MCPWMDT0REDUPMETHOD}\item [MCPWM\_DT0\_RED\_UPMETHOD] 上升沿延迟有效寄存器的更新方式。具体配置值请参考 \hyperref[fielddesc:MCPWMDT0FEDUPMETHOD]{MCPWM\_DT0\_FED\_UPMETHOD}。 (R/W)
\label{fielddesc:MCPWMDT0DEBMODE}\item [MCPWM\_DT0\_DEB\_MODE] 配置表 \ref{tab:mcpwm-s-dt}中的 S8 开关，详细配置信息请参考表 \ref{tab:mcpwm-mod-dt}。
(R/W)
\label{fielddesc:MCPWMDT0AOUTSWAP}\item [MCPWM\_DT0\_A\_OUTSWAP] 配置表 \ref{tab:mcpwm-s-dt} 中的 S6 开关，详细配置信息请参考表 \ref{tab:mcpwm-mod-dt}。 (R/W)
\label{fielddesc:MCPWMDT0BOUTSWAP}\item [MCPWM\_DT0\_B\_OUTSWAP] 配置表 \ref{tab:mcpwm-s-dt} 中的 S7 开关，详细配置信息请参考表 \ref{tab:mcpwm-mod-dt}。 (R/W)
\label{fielddesc:MCPWMDT0REDINSEL}\item [MCPWM\_DT0\_RED\_INSEL] 配置表 \ref{tab:mcpwm-s-dt} 中的 S4 开关，详细配置信息请参考表 \ref{tab:mcpwm-mod-dt}。 (R/W)
\label{fielddesc:MCPWMDT0FEDINSEL}\item [MCPWM\_DT0\_FED\_INSEL] 配置表 \ref{tab:mcpwm-s-dt} 中的 S5 开关，详细配置信息请参考表 \ref{tab:mcpwm-mod-dt}。 (R/W)
\label{fielddesc:MCPWMDT0REDOUTINVERT}\item [MCPWM\_DT0\_RED\_OUTINVERT] 配置表 \ref{tab:mcpwm-s-dt} 中的 S2 开关，详细配置信息请参考表 \ref{tab:mcpwm-mod-dt}。 (R/W)
\label{fielddesc:MCPWMDT0FEDOUTINVERT}\item [MCPWM\_DT0\_FED\_OUTINVERT] 配置表 \ref{tab:mcpwm-s-dt} 中的 S3 开关，详细配置信息请参考表 \ref{tab:mcpwm-mod-dt}。(R/W)
\label{fielddesc:MCPWMDT0AOUTBYPASS}\item [MCPWM\_DT0\_A\_OUTBYPASS] 配置表 \ref{tab:mcpwm-s-dt} 中的 S1 开关，详细配置信息请参考表 \ref{tab:mcpwm-mod-dt}。 (R/W)
\label{fielddesc:MCPWMDT0BOUTBYPASS}\item [MCPWM\_DT0\_B\_OUTBYPASS] 配置表 \ref{tab:mcpwm-s-dt} 中的 S0 开关，详细配置信息请参考表 \ref{tab:mcpwm-mod-dt}。 (R/W)
\label{fielddesc:MCPWMDT0CLKSEL}\item [MCPWM\_DT0\_CLK\_SEL] 选择死区时间生成器 0 的时钟。\\0：PWM\_CLK\\1：PT\_CLK\\ (R/W)
\end{reglist}\end{regdesc}
\end{register}


\begin{register}{H}{MCPWM\_DT0\_FED\_CFG\_REG}{0x005C}\label{regdesc:MCPWMDT0FEDCFGREG}
\regfield{(reserved)}{16}{16}{0000000000000000}%
\regfield{MCPWM\_DT0\_FED}{16}{0}{{0}}%
\reglabel{Reset}\regnewline%
\begin{regdesc}\begin{reglist}
\label{fielddesc:MCPWMDT0FED}\item [MCPWM\_DT0\_FED] FED 的影子寄存器。 (R/W)
\end{reglist}\end{regdesc}
\end{register}


\begin{register}{H}{MCPWM\_DT0\_RED\_CFG\_REG}{0x0060}\label{regdesc:MCPWMDT0REDCFGREG}
\regfield{(reserved)}{16}{16}{0000000000000000}%
\regfield{MCPWM\_DT0\_RED}{16}{0}{{0}}%
\reglabel{Reset}\regnewline%
\begin{regdesc}\begin{reglist}
\label{fielddesc:MCPWMDT0RED}\item [MCPWM\_DT0\_RED] RED 的影子寄存器。(R/W)
\end{reglist}\end{regdesc}
\end{register}


\begin{register}{H}{MCPWM\_CARRIER0\_CFG\_REG}{0x0064}\label{regdesc:MCPWMCARRIER0CFGREG}
\regfield{(reserved)}{18}{14}{000000000000000000}%
\regfield{MCPWM\_CARRIER0\_IN\_INVERT}{1}{13}{{0}}%
\regfield{MCPWM\_CARRIER0\_OUT\_INVERT}{1}{12}{{0}}%
\regfield{MCPWM\_CARRIER0\_OSHTWTH}{4}{8}{{0}}%
\regfield{MCPWM\_CARRIER0\_DUTY}{3}{5}{{0}}%
\regfield{MCPWM\_CARRIER0\_PRESCALE}{4}{1}{{0}}%
\regfield{MCPWM\_CARRIER0\_EN}{1}{0}{{0}}%
\reglabel{Reset}\regnewline%
\begin{regdesc}\begin{reglist}
\label{fielddesc:MCPWMCARRIER0EN}\item [MCPWM\_CARRIER0\_EN] 置 1 时，使能载波 0 的功能。清零时，载波 0 被绕过。(R/W)
\label{fielddesc:MCPWMCARRIER0PRESCALE}\item [MCPWM\_CARRIER0\_PRESCALE] PWM 载波 0 时钟 (PC\_CLK) 的预分频值。PC\_CLK 周期 = PWM\_CLK 周期 * (PWM\_CARRIER0\_PRESCALE + 1)。 (R/W)
\label{fielddesc:MCPWMCARRIER0DUTY}\item [MCPWM\_CARRIER0\_DUTY] 选择载波占空比。占空比 = PWM\_CARRIER0\_DUTY/8。 (R/W)
\label{fielddesc:MCPWMCARRIER0OSHTWTH}\item [MCPWM\_CARRIER0\_OSHTWTH] 载波第一个脉冲的宽度，单位为载波周期。 (R/W)
\label{fielddesc:MCPWMCARRIER0OUTINVERT}\item [MCPWM\_CARRIER0\_OUT\_INVERT] 置 1 时，将此模块的 PWM0A 和 PWM0B 输出反相。(R/W)
\label{fielddesc:MCPWMCARRIER0ININVERT}\item [MCPWM\_CARRIER0\_IN\_INVERT] 置 1 时，将此模块的 PWM0A 和 PWM0B 输入反相。 (R/W)
\end{reglist}\end{regdesc}
\end{register}


\begin{register}{H}{MCPWM\_FH0\_CFG0\_REG}{0x0068}\label{regdesc:MCPWMFH0CFG0REG}
\regfield{(reserved)}{8}{24}{00000000}%
\regfield{MCPWM\_FH0\_B\_OST\_U}{2}{22}{{0}}%
\regfield{MCPWM\_FH0\_B\_OST\_D}{2}{20}{{0}}%
\regfield{MCPWM\_FH0\_B\_CBC\_U}{2}{18}{{0}}%
\regfield{MCPWM\_FH0\_B\_CBC\_D}{2}{16}{{0}}%
\regfield{MCPWM\_FH0\_A\_OST\_U}{2}{14}{{0}}%
\regfield{MCPWM\_FH0\_A\_OST\_D}{2}{12}{{0}}%
\regfield{MCPWM\_FH0\_A\_CBC\_U}{2}{10}{{0}}%
\regfield{MCPWM\_FH0\_A\_CBC\_D}{2}{8}{{0}}%
\regfield{MCPWM\_FH0\_F0\_OST}{1}{7}{{0}}%
\regfield{MCPWM\_FH0\_F1\_OST}{1}{6}{{0}}%
\regfield{MCPWM\_FH0\_F2\_OST}{1}{5}{{0}}%
\regfield{MCPWM\_FH0\_SW\_OST}{1}{4}{{0}}%
\regfield{MCPWM\_FH0\_F0\_CBC}{1}{3}{{0}}%
\regfield{MCPWM\_FH0\_F1\_CBC}{1}{2}{{0}}%
\regfield{MCPWM\_FH0\_F2\_CBC}{1}{1}{{0}}%
\regfield{MCPWM\_FH0\_SW\_CBC}{1}{0}{{0}}%
\reglabel{Reset}\regnewline%
\begin{regdesc}\begin{reglist}
\label{fielddesc:MCPWMFH0SWCBC}\item [MCPWM\_FH0\_SW\_CBC] 使能软件强制逐周期模式操作。\\0：关闭\\1：使能\\ (R/W)
\label{fielddesc:MCPWMFH0F2CBC}\item [MCPWM\_FH0\_F2\_CBC] 设置 fault\_event2 是否触发逐周期模式操作。\\0：关闭\\1：使能\\ (R/W)
\label{fielddesc:MCPWMFH0F1CBC}\item [MCPWM\_FH0\_F1\_CBC] 设置 fault\_event1 是否触发逐周期模式操作。\\0：关闭\\1：使能\\(R/W)
\label{fielddesc:MCPWMFH0F0CBC}\item [MCPWM\_FH0\_F0\_CBC] 设置 fault\_event0 是否触发逐周期模式操作。\\0：关闭\\1：使能\\ (R/W)
\label{fielddesc:MCPWMFH0SWOST}\item [MCPWM\_FH0\_SW\_OST] 软件强制一次性模式操作的使能寄存器。\\0：关闭\\1：使能\\ (R/W)
\label{fielddesc:MCPWMFH0F2OST}\item [MCPWM\_FH0\_F2\_OST] 设置 fault\_event2 是否触发一次性模式操作。\\0：关闭\\1：使能\\ (R/W)
\label{fielddesc:MCPWMFH0F1OST}\item [MCPWM\_FH0\_F1\_OST] 设置 fault\_event1 是否触发一次性模式操作。\\0：关闭\\1：使能\\ (R/W)
\label{fielddesc:MCPWMFH0F0OST}\item [MCPWM\_FH0\_F0\_OST] 设置 fault\_event0 是否触发一次性模式操作。\\0：关闭\\1：使能\\ (R/W)
\item [见下页...]
\end{reglist}\end{regdesc}
\end{register}
\addtocounter{Regfloat}{-1}
\begin{register}{H}{MCPWM\_FH0\_CFG0\_REG}{0x0034}
\begin{regdesc}\begin{reglist}
\item [...接上页]
\label{fielddesc:MCPWMFH0ACBCD}\item [MCPWM\_FH0\_A\_CBC\_D] 定时器递减计数并且发生故障事件时，PWM0A 上的逐周期模式操作。\\0：无\\1：强制拉低\\2：强制拉高\\3：取反\\ (R/W)
\label{fielddesc:MCPWMFH0ACBCU}\item [MCPWM\_FH0\_A\_CBC\_U] 定时器递增计数并且发生故障事件时，PWM0A 上的逐周期模式操作。具体配置值请参考 \hyperref[fielddesc:MCPWMFH0ACBCD]{MCPWM\_FH0\_A\_CBC\_D}。 (R/W)
\label{fielddesc:MCPWMFH0AOSTD}\item [MCPWM\_FH0\_A\_OST\_D] 定时器递减计数并且发生故障事件时，PWM0A 上的一次性模式操作。具体配置值请参考 \hyperref[fielddesc:MCPWMFH0ACBCD]{MCPWM\_FH0\_A\_CBC\_D}。 (R/W)
\label{fielddesc:MCPWMFH0AOSTU}\item [MCPWM\_FH0\_A\_OST\_U] 定时器递增计数并且发生故障事件时，PWM0A 上的一次性模式操作。具体配置值请参考 \hyperref[fielddesc:MCPWMFH0ACBCD]{MCPWM\_FH0\_A\_CBC\_D}。 (R/W)
\label{fielddesc:MCPWMFH0BCBCD}\item [MCPWM\_FH0\_B\_CBC\_D] 定时器递减计数并且发生故障事件时，PWM0B 上的逐周期模式操作。具体配置值请参考 \hyperref[fielddesc:MCPWMFH0ACBCD]{MCPWM\_FH0\_A\_CBC\_D}。 (R/W)
\label{fielddesc:MCPWMFH0BCBCU}\item [MCPWM\_FH0\_B\_CBC\_U] 定时器递增计数并且发生故障事件时，PWM0B 上的逐周期模式操作。具体配置值请参考 \hyperref[fielddesc:MCPWMFH0ACBCD]{MCPWM\_FH0\_A\_CBC\_D}。 (R/W)
\label{fielddesc:MCPWMFH0BOSTD}\item [MCPWM\_FH0\_B\_OST\_D] 定时器递减计数并且发生故障事件时，PWM0B 上的一次性模式操作。具体配置值请参考 \hyperref[fielddesc:MCPWMFH0ACBCD]{MCPWM\_FH0\_A\_CBC\_D}。 (R/W)
\label{fielddesc:MCPWMFH0BOSTU}\item [MCPWM\_FH0\_B\_OST\_U] 定时器递增计数并且发生故障事件时，PWM0B 上的一次性模式操作。具体配置值请参考 \hyperref[fielddesc:MCPWMFH0ACBCD]{MCPWM\_FH0\_A\_CBC\_D}。 (R/W)
\end{reglist}\end{regdesc}
\end{register}


\begin{register}{H}{MCPWM\_FH0\_CFG1\_REG}{0x006C}\label{regdesc:MCPWMFH0CFG1REG}
\regfield{(reserved)}{27}{5}{000000000000000000000000000}%
\regfield{MCPWM\_FH0\_FORCE\_OST}{1}{4}{{0}}%
\regfield{MCPWM\_FH0\_FORCE\_CBC}{1}{3}{{0}}%
\regfield{MCPWM\_FH0\_CBCPULSE}{2}{1}{{0}}%
\regfield{MCPWM\_FH0\_CLR\_OST}{1}{0}{{0}}%
\reglabel{Reset}\regnewline%
\begin{regdesc}\begin{reglist}
\label{fielddesc:MCPWMFH0CLROST}\item [MCPWM\_FH0\_CLR\_OST] 上升沿时将清除持续一次性模式操作。 (R/W)
\label{fielddesc:MCPWMFH0CBCPULSE}\item [MCPWM\_FH0\_CBCPULSE] 设置逐周期模式操作更新的时间点。\\bit0 为 1：发生 TEZ 事件时\\bit1 为 1：发生 TEP 事件时\\ (R/W)
\label{fielddesc:MCPWMFH0FORCECBC}\item [MCPWM\_FH0\_FORCE\_CBC] 通过软件将此位的值取反，可触发逐周期模式的操作。 (R/W)
\label{fielddesc:MCPWMFH0FORCEOST}\item [MCPWM\_FH0\_FORCE\_OST] 通过软件将此位的值取反，可触发一次性模式的操作。 (R/W)
\end{reglist}\end{regdesc}
\end{register}


\begin{register}{H}{MCPWM\_FH0\_STATUS\_REG}{0x0070}\label{regdesc:MCPWMFH0STATUSREG}
\regfield{(reserved)}{30}{2}{000000000000000000000000000000}%
\regfield{MCPWM\_FH0\_OST\_ON}{1}{1}{{0}}%
\regfield{MCPWM\_FH0\_CBC\_ON}{1}{0}{{0}}%
\reglabel{Reset}\regnewline%
\begin{regdesc}\begin{reglist}
\label{fielddesc:MCPWMFH0CBCON}\item [MCPWM\_FH0\_CBC\_ON] 表示是否正在进行逐周期模式的操作。此字段由硬件置 1 和清零。\\
0：不存在正在进行的逐周期模式操作\\
1：逐周期模式的操作正在进行\\
(RO)
\label{fielddesc:MCPWMFH0OSTON}\item [MCPWM\_FH0\_OST\_ON] 表示是否正在进行一次性模式的操作。此字段由硬件置 1 和清零。\\
0：不存在正在进行的一次性模式操作\\
1：一次性模式的操作正在进行\\
(RO)
\end{reglist}\end{regdesc}
\end{register}


\begin{register}{H}{MCPWM\_GEN1\_STMP\_CFG\_REG}{0x0074}\label{regdesc:MCPWMGEN1STMPCFGREG}
\regfield{(reserved)}{22}{10}{0000000000000000000000}%
\regfield{MCPWM\_GEN1\_B\_SHDW\_FULL}{1}{9}{{0}}%
\regfield{MCPWM\_GEN1\_A\_SHDW\_FULL}{1}{8}{{0}}%
\regfield{MCPWM\_GEN1\_B\_UPMETHOD}{4}{4}{{0}}%
\regfield{MCPWM\_GEN1\_A\_UPMETHOD}{4}{0}{{0}}%
\reglabel{Reset}\regnewline%
\begin{regdesc}\begin{reglist}
\label{fielddesc:MCPWMGEN1AUPMETHOD}\item [MCPWM\_GEN1\_A\_UPMETHOD] PWM 生成器 1 时间戳寄存器 A 有效寄存器的更新方式。\\所有 bit 值 为 0：立即更新\\bit0 为 1：发生 TEZ 事件时更新\\bit1 为 1：发生 TEP 事件时更新\\bit2 为 1：发生同步时间时更新\\bit3 为 1：关闭更新\\ (R/W)
\label{fielddesc:MCPWMGEN1BUPMETHOD}\item [MCPWM\_GEN1\_B\_UPMETHOD] PWM 生成器 1 时间戳寄存器 B 有效寄存器的更新方式。具体配置值请参考 \hyperref[fielddesc:MCPWMGEN1AUPMETHOD]{MCPWM\_GEN1\_A\_UPMETHOD}。 (R/W)
\label{fielddesc:MCPWMGEN1ASHDWFULL}\item [MCPWM\_GEN1\_A\_SHDW\_FULL] 配置是否在某个特定时间将影子寄存器中的值写入对应的有效寄存器中。此字段由硬件置 1 或复位。\\
0：将影子寄存器中最新的值写入有效寄存器 A 中\\
1：将值写入 PWM 生成器 1 时间戳寄存器 A 的影子寄存器中，等待传输给有效寄存器 A\\
(R/SC/WTC)
\label{fielddesc:MCPWMGEN1BSHDWFULL}\item [MCPWM\_GEN1\_B\_SHDW\_FULL] 配置是否在某个特定时间将影子寄存器中的值写入对应的有效寄存器中。此字段由硬件置 1 或复位。\\
0：将影子寄存器中最新的值写入有效寄存器 B 中\\
1：将值写入 PWM 生成器 1 时间戳寄存器 B 的影子寄存器中，等待传输给有效寄存器 A\\
(R/SC/WTC)
\end{reglist}\end{regdesc}
\end{register}


\begin{register}{H}{MCPWM\_GEN1\_TSTMP\_A\_REG}{0x0078}\label{regdesc:MCPWMGEN1TSTMPAREG}
\regfield{(reserved)}{16}{16}{0000000000000000}%
\regfield{MCPWM\_GEN1\_A}{16}{0}{{0}}%
\reglabel{Reset}\regnewline%
\begin{regdesc}\begin{reglist}
\label{fielddesc:MCPWMGEN1A}\item [MCPWM\_GEN1\_A] PWM 生成器 1 时间戳寄存器 A 的影子寄存器。(R/W)
\end{reglist}\end{regdesc}
\end{register}


\begin{register}{H}{MCPWM\_GEN1\_TSTMP\_B\_REG}{0x007C}\label{regdesc:MCPWMGEN1TSTMPBREG}
\regfield{(reserved)}{16}{16}{0000000000000000}%
\regfield{MCPWM\_GEN1\_B}{16}{0}{{0}}%
\reglabel{Reset}\regnewline%
\begin{regdesc}\begin{reglist}
\label{fielddesc:MCPWMGEN1B}\item [MCPWM\_GEN1\_B] PWM 生成器 1 时间戳寄存器 B 的影子寄存器。 (R/W)
\end{reglist}\end{regdesc}
\end{register}


\begin{register}{H}{MCPWM\_GEN1\_CFG0\_REG}{0x0080}\label{regdesc:MCPWMGEN1CFG0REG}
\regfield{(reserved)}{22}{10}{0000000000000000000000}%
\regfield{MCPWM\_GEN1\_T1\_SEL}{3}{7}{{0}}%
\regfield{MCPWM\_GEN1\_T0\_SEL}{3}{4}{{0}}%
\regfield{MCPWM\_GEN1\_CFG\_UPMETHOD}{4}{0}{{0}}%
\reglabel{Reset}\regnewline%
\begin{regdesc}\begin{reglist}
\label{fielddesc:MCPWMGEN1CFGUPMETHOD}\item [MCPWM\_GEN1\_CFG\_UPMETHOD] 配置 PWM 生成器 1 有效寄存器的更新方式。\\
所有 bit 值为 0：立即更新\\
bit0 为 1：发生 TEZ 事件时更新\\
bit1 为 1：发生 TEP 事件时更新\\
bit2 为 1：发生同步时间时更新\\
bit3 为 1：关闭更新\\
(R/W)
\label{fielddesc:MCPWMGEN1T0SEL}\item [MCPWM\_GEN1\_T0\_SEL] 配置 PWM 生成器 1 event\_t0 的信号源，立即生效\\
0：fault\_event0\\
1：fault\_event1\\
2：fault\_event2\\
3：sync\_taken\\
4：无\\
(R/W)
\label{fielddesc:MCPWMGEN1T1SEL}\item [MCPWM\_GEN1\_T1\_SEL] 配置 PWM 生成器 1 event\_t1 的信号源，立即生效。\\
0：fault\_event0\\
1：fault\_event1\\
2：fault\_event2\\
3：sync\_taken\\
4：无\\
(R/W)
\end{reglist}\end{regdesc}
\end{register}


\begin{register}{H}{MCPWM\_GEN1\_FORCE\_REG}{0x0084}\label{regdesc:MCPWMGEN1FORCEREG}
\regfield{(reserved)}{16}{16}{0000000000000000}%
\regfield{MCPWM\_GEN1\_B\_NCIFORCE\_MODE}{2}{14}{{0}}%
\regfield{MCPWM\_GEN1\_B\_NCIFORCE}{1}{13}{{0}}%
\regfield{MCPWM\_GEN1\_A\_NCIFORCE\_MODE}{2}{11}{{0}}%
\regfield{MCPWM\_GEN1\_A\_NCIFORCE}{1}{10}{{0}}%
\regfield{MCPWM\_GEN1\_B\_CNTUFORCE\_MODE}{2}{8}{{0}}%
\regfield{MCPWM\_GEN1\_A\_CNTUFORCE\_MODE}{2}{6}{{0}}%
\regfield{MCPWM\_GEN1\_CNTUFORCE\_UPMETHOD}{6}{0}{{0x20}}%
\reglabel{Reset}\regnewline%
\begin{regdesc}\begin{reglist}
\label{fielddesc:MCPWMGEN1CNTUFORCEUPMETHOD}\item [MCPWM\_GEN1\_CNTUFORCE\_UPMETHOD] PWM 生成器 1 持续软件强制的更新方式。\\
0：立即更新\\
bit0 为 1：发生 TEZ 事件时更新\\
bit1 为 1：发生 TEP 事件时更新\\
bit2 为 1：发生 TEA 事件时更新\\bit3 为 1：发生 TEB 事件时更新\\
bit4 为 1：发生同步时间时更新\\
bit5 为 1：关闭更新\\
本文档中的 TEA/B 表示定时器值等于寄存器 A/B 生成的事件。\\
(R/W)
\label{fielddesc:MCPWMGEN1ACNTUFORCEMODE}\item [MCPWM\_GEN1\_A\_CNTUFORCE\_MODE] 用于 PWM1A 的连续软件强制事件。\\
0：关闭\\
1：拉低\\
2：拉高\\
3：关闭\\
(R/W)
\label{fielddesc:MCPWMGEN1BCNTUFORCEMODE}\item [MCPWM\_GEN1\_B\_CNTUFORCE\_MODE] 用于 PWM1B 的连续软件强制事件。\\
0：关闭\\
1：拉低\\
2：拉高\\
3：关闭\\
(R/W)
\label{fielddesc:MCPWMGEN1ANCIFORCE}\item [MCPWM\_GEN1\_A\_NCIFORCE] 该位的值取反时将触发 PWM1A 上的非连续即时软件强制事件。 (R/W)
\label{fielddesc:MCPWMGEN1ANCIFORCEMODE}\item [MCPWM\_GEN1\_A\_NCIFORCE\_MODE] 用于 PWM1A 的非持续性即时软件强制事件。\\
0：关闭\\
1：拉低\\
2：拉高\\
3：关闭\\
(R/W)
\item [见下页...]
\end{reglist}\end{regdesc}
\end{register}
\addtocounter{Regfloat}{-1}
\begin{register}{H}{MCPWM\_GEN1\_FORCE\_REG}{0x0034}
\begin{regdesc}\begin{reglist}
\item [...接上页]
\label{fielddesc:MCPWMGEN1BNCIFORCE}\item [MCPWM\_GEN1\_B\_NCIFORCE] 该位的值取反时将触发PWM1B上的非连续即时软件强制事件。 (R/W)
\label{fielddesc:MCPWMGEN1BNCIFORCEMODE}\item [MCPWM\_GEN1\_B\_NCIFORCE\_MODE] 用于 PWM1B 的非持续性即时软件强制事件。\\
0：关闭\\
1：拉低\\
2：拉高\\
3：关闭\\
(R/W)
\end{reglist}\end{regdesc}
\end{register}


\begin{register}{H}{MCPWM\_GEN1\_A\_REG}{0x0088}\label{regdesc:MCPWMGEN1AREG}
\regfield{(reserved)}{8}{24}{00000000}%
\regfield{MCPWM\_GEN1\_A\_DT1}{2}{22}{{0}}%
\regfield{MCPWM\_GEN1\_A\_DT0}{2}{20}{{0}}%
\regfield{MCPWM\_GEN1\_A\_DTEB}{2}{18}{{0}}%
\regfield{MCPWM\_GEN1\_A\_DTEA}{2}{16}{{0}}%
\regfield{MCPWM\_GEN1\_A\_DTEP}{2}{14}{{0}}%
\regfield{MCPWM\_GEN1\_A\_DTEZ}{2}{12}{{0}}%
\regfield{MCPWM\_GEN1\_A\_UT1}{2}{10}{{0}}%
\regfield{MCPWM\_GEN1\_A\_UT0}{2}{8}{{0}}%
\regfield{MCPWM\_GEN1\_A\_UTEB}{2}{6}{{0}}%
\regfield{MCPWM\_GEN1\_A\_UTEA}{2}{4}{{0}}%
\regfield{MCPWM\_GEN1\_A\_UTEP}{2}{2}{{0}}%
\regfield{MCPWM\_GEN1\_A\_UTEZ}{2}{0}{{0}}%
\reglabel{Reset}\regnewline%
\begin{regdesc}\begin{reglist}
\label{fielddesc:MCPWMGEN1AUTEZ}\item [MCPWM\_GEN1\_A\_UTEZ] 定时器递增计数时，TEZ 在 PWM1A 上触发的操作。\\
0：无\\
1：拉低\\
2：拉高\\
3：取反\\
(R/W)
\label{fielddesc:MCPWMGEN1AUTEP}\item [MCPWM\_GEN1\_A\_UTEP] 定时器递增计数时，TEP 在 PWM1A 上触发的操作。\\
0：无\\
1：拉低\\
2：拉高\\
3：取反\\
(R/W)
\label{fielddesc:MCPWMGEN1AUTEA}\item [MCPWM\_GEN1\_A\_UTEA] 定时器递增计数时，TEA 在 PWM1A 上触发的操作。\\
0：无\\
1：拉低\\
2：拉高\\
3：取反\\
(R/W)
\label{fielddesc:MCPWMGEN1AUTEB}\item [MCPWM\_GEN1\_A\_UTEB] 定时器递增计数时，TEB 在 PWM1A 上触发的操作。\\
0：无\\
1：拉低\\
2：拉高\\
3：取反\\
(R/W)
\label{fielddesc:MCPWMGEN1AUT0}\item [MCPWM\_GEN1\_A\_UT0] 定时器递增计数时，event\_t0 在 PWM1A 上触发的操作。\\
0：无\\
1：拉低\\
2：拉高\\
3：取反\\
(R/W)
\item [见下页...]
\end{reglist}\end{regdesc}
\end{register}
\addtocounter{Regfloat}{-1}
\begin{register}{H}{MCPWM\_GEN1\_A\_REG}{0x0034}
\begin{regdesc}\begin{reglist}
\item [...接上页]
\label{fielddesc:MCPWMGEN1AUT1}\item [MCPWM\_GEN1\_A\_UT1] 定时器递增计数时，event\_t1 在 PWM1A 上触发的操作。\\
0：无\\
1：拉低\\
2：拉高\\
3：取反\\
(R/W)
\label{fielddesc:MCPWMGEN1ADTEZ}\item [MCPWM\_GEN1\_A\_DTEZ] 定时器递减计数时，TEZ 在 PWM1A 上触发的操作。\\
0：无\\
1：拉低\\
2：拉高\\
3：取反\\
(R/W)
\label{fielddesc:MCPWMGEN1ADTEP}\item [MCPWM\_GEN1\_A\_DTEP] 定时器递减计数时，TEP 在 PWM1A 上触发的操作。\\
0：无\\
1：拉低\\
2：拉高\\
3：取反\\
(R/W)
\label{fielddesc:MCPWMGEN1ADTEA}\item [MCPWM\_GEN1\_A\_DTEA] 定时器递减计数时，TEA 在 PWM1A 上触发的操作。\\
0：无\\
1：拉低\\
2：拉高\\
3：取反\\
(R/W)
\label{fielddesc:MCPWMGEN1ADTEB}\item [MCPWM\_GEN1\_A\_DTEB] 定时器递减计数时，TEB 在 PWM1A 上触发的操作。\\
0：无\\
1：拉低\\
2：拉高\\
3：取反\\
(R/W)
\label{fielddesc:MCPWMGEN1ADT0}\item [MCPWM\_GEN1\_A\_DT0] 定时器递减计数时，event\_t0 在 PWM1A 上触发的操作。\\
0：无\\
1：拉低\\
2：拉高\\
3：取反\\
(R/W)
\label{fielddesc:MCPWMGEN1ADT1}\item [MCPWM\_GEN1\_A\_DT1] 定时器递减计数时，event\_t1 在 PWM1A 上触发的操作。\\
0：无\\
1：拉低\\
2：拉高\\
3：取反\\
(R/W)
\end{reglist}\end{regdesc}
\end{register}


\begin{register}{H}{MCPWM\_GEN1\_B\_REG}{0x008C}\label{regdesc:MCPWMGEN1BREG}
\regfield{(reserved)}{8}{24}{00000000}%
\regfield{MCPWM\_GEN1\_B\_DT1}{2}{22}{{0}}%
\regfield{MCPWM\_GEN1\_B\_DT0}{2}{20}{{0}}%
\regfield{MCPWM\_GEN1\_B\_DTEB}{2}{18}{{0}}%
\regfield{MCPWM\_GEN1\_B\_DTEA}{2}{16}{{0}}%
\regfield{MCPWM\_GEN1\_B\_DTEP}{2}{14}{{0}}%
\regfield{MCPWM\_GEN1\_B\_DTEZ}{2}{12}{{0}}%
\regfield{MCPWM\_GEN1\_B\_UT1}{2}{10}{{0}}%
\regfield{MCPWM\_GEN1\_B\_UT0}{2}{8}{{0}}%
\regfield{MCPWM\_GEN1\_B\_UTEB}{2}{6}{{0}}%
\regfield{MCPWM\_GEN1\_B\_UTEA}{2}{4}{{0}}%
\regfield{MCPWM\_GEN1\_B\_UTEP}{2}{2}{{0}}%
\regfield{MCPWM\_GEN1\_B\_UTEZ}{2}{0}{{0}}%
\reglabel{Reset}\regnewline%
\begin{regdesc}\begin{reglist}
\label{fielddesc:MCPWMGEN1BUTEZ}\item [MCPWM\_GEN1\_B\_UTEZ] 定时器递增计数时，TEZ 在 PWM1B 上触发的操作。\\
0：无\\
1：拉低\\
2：拉高\\
3：取反\\
(R/W)
\label{fielddesc:MCPWMGEN1BUTEP}\item [MCPWM\_GEN1\_B\_UTEP] 定时器递增计数时，TEP 在 PWM1B 上触发的操作。\\
0：无\\
1：拉低\\
2：拉高\\
3：取反\\
(R/W)
\label{fielddesc:MCPWMGEN1BUTEA}\item [MCPWM\_GEN1\_B\_UTEA] 定时器递增计数时，TEA 在 PWM1B 上触发的操作。\\
0：无\\
1：拉低\\
2：拉高\\
3：取反\\
(R/W)
\label{fielddesc:MCPWMGEN1BUTEB}\item [MCPWM\_GEN1\_B\_UTEB] 定时器递增计数时，TEB 在 PWM1B 上触发的操作。\\
0：无\\
1：拉低\\
2：拉高\\
3：取反\\
(R/W)
\label{fielddesc:MCPWMGEN1BUT0}\item [MCPWM\_GEN1\_B\_UT0] 定时器递增计数时，event\_t0 在 PWM1B 上触发的操作。\\
0：无\\
1：拉低\\
2：拉高\\
3：取反\\
(R/W)
\item [见下页...]
\end{reglist}\end{regdesc}
\end{register}
\addtocounter{Regfloat}{-1}
\begin{register}{H}{MCPWM\_GEN1\_B\_REG}{0x0034}
\begin{regdesc}\begin{reglist}
\item [...接上页]
\label{fielddesc:MCPWMGEN1BUT1}\item [MCPWM\_GEN1\_B\_UT1] 定时器递增计数时，event\_t1 在 PWM1B 上触发的操作。\\
0：无\\
1：拉低\\
2：拉高\\
3：取反\\
(R/W)
\label{fielddesc:MCPWMGEN1BDTEZ}\item [MCPWM\_GEN1\_B\_DTEZ] 定时器递减计数时，TEZ 在 PWM1B 上触发的操作。\\
0：无\\
1：拉低\\
2：拉高\\
3：取反\\
(R/W)
\label{fielddesc:MCPWMGEN1BDTEP}\item [MCPWM\_GEN1\_B\_DTEP] 定时器递减计数时，TEP 在 PWM1B 上触发的操作。\\
0：无\\
1：拉低\\
2：拉高\\
3：取反\\
(R/W)
\label{fielddesc:MCPWMGEN1BDTEA}\item [MCPWM\_GEN1\_B\_DTEA] 定时器递减计数时，TEA 在 PWM1B 上触发的操作。\\
0：无\\
1：拉低\\
2：拉高\\
3：取反\\
(R/W)
\label{fielddesc:MCPWMGEN1BDTEB}\item [MCPWM\_GEN1\_B\_DTEB] 定时器递减计数时，TEB 在 PWM1B 上触发的操作。\\
0：无\\
1：拉低\\
2：拉高\\
3：取反\\
(R/W)
\label{fielddesc:MCPWMGEN1BDT0}\item [MCPWM\_GEN1\_B\_DT0] 定时器递减计数时，event\_t0 在 PWM1B 上触发的操作。\\
0：无\\
1：拉低\\
2：拉高\\
3：取反\\
(R/W)
\label{fielddesc:MCPWMGEN1BDT1}\item [MCPWM\_GEN1\_B\_DT1] 定时器递减计数时，event\_t1 在 PWM1B 上触发的操作。\\
0：无\\
1：拉低\\
2：拉高\\
3：取反\\
(R/W)
\end{reglist}\end{regdesc}
\end{register}


\begin{register}{H}{MCPWM\_DT1\_CFG\_REG}{0x0090}\label{regdesc:MCPWMDT1CFGREG}
\regfield{(reserved)}{14}{18}{00000000000000}%
\regfield{MCPWM\_DT1\_CLK\_SEL}{1}{17}{{0}}%
\regfield{MCPWM\_DT1\_B\_OUTBYPASS}{1}{16}{{1}}%
\regfield{MCPWM\_DT1\_A\_OUTBYPASS}{1}{15}{{1}}%
\regfield{MCPWM\_DT1\_FED\_OUTINVERT}{1}{14}{{0}}%
\regfield{MCPWM\_DT1\_RED\_OUTINVERT}{1}{13}{{0}}%
\regfield{MCPWM\_DT1\_FED\_INSEL}{1}{12}{{0}}%
\regfield{MCPWM\_DT1\_RED\_INSEL}{1}{11}{{0}}%
\regfield{MCPWM\_DT1\_B\_OUTSWAP}{1}{10}{{0}}%
\regfield{MCPWM\_DT1\_A\_OUTSWAP}{1}{9}{{0}}%
\regfield{MCPWM\_DT1\_DEB\_MODE}{1}{8}{{0}}%
\regfield{MCPWM\_DT1\_RED\_UPMETHOD}{4}{4}{{0}}%
\regfield{MCPWM\_DT1\_FED\_UPMETHOD}{4}{0}{{0}}%
\reglabel{Reset}\regnewline%
\begin{regdesc}\begin{reglist}
\label{fielddesc:MCPWMDT1FEDUPMETHOD}\item [MCPWM\_DT1\_FED\_UPMETHOD] FED(下降沿延迟)有效寄存器的更新方式。\\
0：立即更新\\
bit0 为 1：发生 TEZ 事件时更新\\
bit1 为 1：发生 TEP 事件时更新\\
bit2 为 1：发生同步事件时更新\\
bit3 为 1：关闭更新\\
(R/W)
\label{fielddesc:MCPWMDT1REDUPMETHOD}\item [MCPWM\_DT1\_RED\_UPMETHOD] RED(上升沿延迟)有效寄存器的更新方式。\\
0：立即更新\\
bit0 为 1：发生 TEZ 事件时更新\\
bit1 为 1：发生 TEP 事件时更新\\
bit2 为 1：发生同步事件时更新\\
bit3 为 1：关闭更新\\
(R/W)
\label{fielddesc:MCPWMDT1DEBMODE}\item [MCPWM\_DT1\_DEB\_MODE] 配置表 \ref{tab:mcpwm-s-dt}中的 S8 开关，详细配置信息请参考表 \ref{tab:mcpwm-mod-dt}。
(R/W)
\label{fielddesc:MCPWMDT1AOUTSWAP}\item [MCPWM\_DT1\_A\_OUTSWAP] 配置表 \ref{tab:mcpwm-s-dt} 中的 S6 开关，详细配置信息请参考表 \ref{tab:mcpwm-mod-dt}。 (R/W)
\label{fielddesc:MCPWMDT1BOUTSWAP}\item [MCPWM\_DT1\_B\_OUTSWAP] 配置表 \ref{tab:mcpwm-s-dt} 中的 S7 开关，详细配置信息请参考表 \ref{tab:mcpwm-mod-dt}。 (R/W)
\label{fielddesc:MCPWMDT1REDINSEL}\item [MCPWM\_DT1\_RED\_INSEL] 配置表 \ref{tab:mcpwm-s-dt} 中的 S4 开关，详细配置信息请参考表 \ref{tab:mcpwm-mod-dt}。 (R/W)
\label{fielddesc:MCPWMDT1FEDINSEL}\item [MCPWM\_DT1\_FED\_INSEL] 配置表 \ref{tab:mcpwm-s-dt} 中的 S5 开关，详细配置信息请参考表 \ref{tab:mcpwm-mod-dt}。 (R/W)
\label{fielddesc:MCPWMDT1REDOUTINVERT}\item [MCPWM\_DT1\_RED\_OUTINVERT] 配置表 \ref{tab:mcpwm-s-dt} 中的 S2 开关，详细配置信息请参考表 \ref{tab:mcpwm-mod-dt}。 (R/W)
\label{fielddesc:MCPWMDT1FEDOUTINVERT}\item [MCPWM\_DT1\_FED\_OUTINVERT] 配置表 \ref{tab:mcpwm-s-dt} 中的 S3 开关，详细配置信息请参考表 \ref{tab:mcpwm-mod-dt}。 (R/W)
\label{fielddesc:MCPWMDT1AOUTBYPASS}\item [MCPWM\_DT1\_A\_OUTBYPASS] 配置表 \ref{tab:mcpwm-s-dt} 中的 S1 开关，详细配置信息请参考表 \ref{tab:mcpwm-mod-dt}。 (R/W)
\label{fielddesc:MCPWMDT1BOUTBYPASS}\item [MCPWM\_DT1\_B\_OUTBYPASS] 配置表 \ref{tab:mcpwm-s-dt} 中的 S0 开关，详细配置信息请参考表 \ref{tab:mcpwm-mod-dt}。 (R/W)
\label{fielddesc:MCPWMDT1CLKSEL}\item [MCPWM\_DT1\_CLK\_SEL] 设置死区时间生成器时钟。\\
0：PWM\_CLK\\
1：PT\_CLK\\
(R/W)
\end{reglist}\end{regdesc}
\end{register}


\begin{register}{H}{MCPWM\_DT1\_FED\_CFG\_REG}{0x0094}\label{regdesc:MCPWMDT1FEDCFGREG}
\regfield{(reserved)}{16}{16}{0000000000000000}%
\regfield{MCPWM\_DT1\_FED}{16}{0}{{0}}%
\reglabel{Reset}\regnewline%
\begin{regdesc}\begin{reglist}
\label{fielddesc:MCPWMDT1FED}\item [MCPWM\_DT1\_FED] FED 影子寄存器。 (R/W)
\end{reglist}\end{regdesc}
\end{register}


\begin{register}{H}{MCPWM\_DT1\_RED\_CFG\_REG}{0x0098}\label{regdesc:MCPWMDT1REDCFGREG}
\regfield{(reserved)}{16}{16}{0000000000000000}%
\regfield{MCPWM\_DT1\_RED}{16}{0}{{0}}%
\reglabel{Reset}\regnewline%
\begin{regdesc}\begin{reglist}
\label{fielddesc:MCPWMDT1RED}\item [MCPWM\_DT1\_RED] RED 影子寄存器。 (R/W)
\end{reglist}\end{regdesc}
\end{register}


\begin{register}{H}{MCPWM\_CARRIER1\_CFG\_REG}{0x009C}\label{regdesc:MCPWMCARRIER1CFGREG}
\regfield{(reserved)}{18}{14}{000000000000000000}%
\regfield{MCPWM\_CARRIER1\_IN\_INVERT}{1}{13}{{0}}%
\regfield{MCPWM\_CARRIER1\_OUT\_INVERT}{1}{12}{{0}}%
\regfield{MCPWM\_CARRIER1\_OSHTWTH}{4}{8}{{0}}%
\regfield{MCPWM\_CARRIER1\_DUTY}{3}{5}{{0}}%
\regfield{MCPWM\_CARRIER1\_PRESCALE}{4}{1}{{0}}%
\regfield{MCPWM\_CARRIER1\_EN}{1}{0}{{0}}%
\reglabel{Reset}\regnewline%
\begin{regdesc}\begin{reglist}
\label{fielddesc:MCPWMCARRIER1EN}\item [MCPWM\_CARRIER1\_EN] 置 1 时，使能载波 1 功能。此位清零时，绕过载波 1。 (R/W)
\label{fielddesc:MCPWMCARRIER1PRESCALE}\item [MCPWM\_CARRIER1\_PRESCALE] 配置 PWM 载波 1 时钟 (PC\_CLK) 预分频值。PC\_CLK 周期 = PWM\_CLK 周期 * (PWM\_CARRIER0\_PRESCALE + 1)。 (R/W)
\label{fielddesc:MCPWMCARRIER1DUTY}\item [MCPWM\_CARRIER1\_DUTY] 设置载波占空比。占空比 = PWM\_CARRIER0\_DUTY / 8。 (R/W)
\label{fielddesc:MCPWMCARRIER1OSHTWTH}\item [MCPWM\_CARRIER1\_OSHTWTH] 配置载波第一个脉冲的宽度，单位为载波周期。(R/W)
\label{fielddesc:MCPWMCARRIER1OUTINVERT}\item [MCPWM\_CARRIER1\_OUT\_INVERT] 置 1 时，将此模块的 PWM1A 和 PWM1B 输出反相。 (R/W)
\label{fielddesc:MCPWMCARRIER1ININVERT}\item [MCPWM\_CARRIER1\_IN\_INVERT] 置 1 时，将此模块的 PWM1A 和 PWM1B 输入反相。 (R/W)
\end{reglist}\end{regdesc}
\end{register}


\begin{register}{H}{MCPWM\_FH1\_CFG0\_REG}{0x00A0}\label{regdesc:MCPWMFH1CFG0REG}
\regfield{(reserved)}{8}{24}{00000000}%
\regfield{MCPWM\_FH1\_B\_OST\_U}{2}{22}{{0}}%
\regfield{MCPWM\_FH1\_B\_OST\_D}{2}{20}{{0}}%
\regfield{MCPWM\_FH1\_B\_CBC\_U}{2}{18}{{0}}%
\regfield{MCPWM\_FH1\_B\_CBC\_D}{2}{16}{{0}}%
\regfield{MCPWM\_FH1\_A\_OST\_U}{2}{14}{{0}}%
\regfield{MCPWM\_FH1\_A\_OST\_D}{2}{12}{{0}}%
\regfield{MCPWM\_FH1\_A\_CBC\_U}{2}{10}{{0}}%
\regfield{MCPWM\_FH1\_A\_CBC\_D}{2}{8}{{0}}%
\regfield{MCPWM\_FH1\_F0\_OST}{1}{7}{{0}}%
\regfield{MCPWM\_FH1\_F1\_OST}{1}{6}{{0}}%
\regfield{MCPWM\_FH1\_F2\_OST}{1}{5}{{0}}%
\regfield{MCPWM\_FH1\_SW\_OST}{1}{4}{{0}}%
\regfield{MCPWM\_FH1\_F0\_CBC}{1}{3}{{0}}%
\regfield{MCPWM\_FH1\_F1\_CBC}{1}{2}{{0}}%
\regfield{MCPWM\_FH1\_F2\_CBC}{1}{1}{{0}}%
\regfield{MCPWM\_FH1\_SW\_CBC}{1}{0}{{0}}%
\reglabel{Reset}\regnewline%
\begin{regdesc}\begin{reglist}
\label{fielddesc:MCPWMFH1SWCBC}\item [MCPWM\_FH1\_SW\_CBC] 软件强制逐周期模式操作的使能寄存器。\\
0：关闭\\
1：使能\\
(R/W)
\label{fielddesc:MCPWMFH1F2CBC}\item [MCPWM\_FH1\_F2\_CBC] 设置 fault\_event2 触发逐周期模式操作。 \\
0：关闭\\
1：使能\\
(R/W)
\label{fielddesc:MCPWMFH1F1CBC}\item [MCPWM\_FH1\_F1\_CBC] 设置 fault\_event1 触发逐周期模式操作。 \\
0：关闭\\
1：使能\\
(R/W)
\label{fielddesc:MCPWMFH1F0CBC}\item [MCPWM\_FH1\_F0\_CBC] 设置 fault\_event0 触发逐周期模式操作。\\
0：关闭\\
1：使能\\
(R/W)
\label{fielddesc:MCPWMFH1SWOST}\item [MCPWM\_FH1\_SW\_OST] 软件强制一次性模式操作的使能寄存器。\\
0：关闭\\
1：使能\\
(R/W)
\label{fielddesc:MCPWMFH1F2OST}\item [MCPWM\_FH1\_F2\_OST] 设置 fault\_event2 触发一次性模式操作。\\
0：关闭\\
1：使能\\
(R/W)
\label{fielddesc:MCPWMFH1F1OST}\item [MCPWM\_FH1\_F1\_OST] 设置 fault\_event1 触发一次性模式操作。\\
0：关闭\\
1：使能\\
(R/W)
\label{fielddesc:MCPWMFH1F0OST}\item [MCPWM\_FH1\_F0\_OST] 设置 fault\_event0 触发一次性模式操作。\\
0：关闭\\
1：使能\\
(R/W)
\item [见下页...]
\end{reglist}\end{regdesc}
\end{register}
\addtocounter{Regfloat}{-1}
\begin{register}{H}{MCPWM\_FH1\_CFG0\_REG}{0x0034}
\begin{regdesc}\begin{reglist}
\item [...接上页]
\label{fielddesc:MCPWMFH1ACBCD}\item [MCPWM\_FH1\_A\_CBC\_D] 定时器递减计数并且发生故障事件时，PWM1A 上的逐周期模式操作。\\
0：无\\
1：强制拉低\\
2：强制拉高\\
3：取反\\
(R/W)
\label{fielddesc:MCPWMFH1ACBCU}\item [MCPWM\_FH1\_A\_CBC\_U] 定时器递增计数并且发生故障事件时，PWM1A 上的逐周期模式操作。\\
0：无\\
1：强制拉低\\
2：强制拉高\\
3：取反\\
(R/W)
\label{fielddesc:MCPWMFH1AOSTD}\item [MCPWM\_FH1\_A\_OST\_D] 定时器递减计数并且发生故障事件时，PWM1A 上的一次性模式操作。\\
0：无\\
1：强制拉低\\
2：强制拉高\\
3：取反\\
(R/W)
\label{fielddesc:MCPWMFH1AOSTU}\item [MCPWM\_FH1\_A\_OST\_U] 定时器递增计数并且发生故障事件时，PWM1A 上的一次性模式操作。\\
0：无\\
1：强制拉低\\
2：强制拉高\\
3：取反\\
(R/W)
\label{fielddesc:MCPWMFH1BCBCD}\item [MCPWM\_FH1\_B\_CBC\_D] 定时器递减计数并且发生故障事件时，PWM1B 上的逐周期模式操作。\\
0：无\\
1：强制拉低\\
2：强制拉高\\
3：取反\\
(R/W)
\label{fielddesc:MCPWMFH1BCBCU}\item [MCPWM\_FH1\_B\_CBC\_U] 定时器递增计数并且发生故障事件时，PWM1B 上的逐周期模式操作。\\
0：无\\
1：强制拉低\\
2：强制拉高\\
3：取反\\
(R/W)
\item [见下页...]
\end{reglist}\end{regdesc}
\end{register}
\addtocounter{Regfloat}{-1}
\begin{register}{H}{MCPWM\_FH1\_CFG0\_REG}{0x0034}
\begin{regdesc}\begin{reglist}
\item [...接上页]
\label{fielddesc:MCPWMFH1BOSTD}\item [MCPWM\_FH1\_B\_OST\_D] 定时器递减计数并且发生故障事件时，PWM1B 上的一次性模式操作。\\
0：无\\
1：强制拉低\\
2：强制拉高\\
3：取反\\
(R/W)
\label{fielddesc:MCPWMFH1BOSTU}\item [MCPWM\_FH1\_B\_OST\_U] 定时器递增计数并且发生故障事件时，PWM1B 上的一次性模式操作。\\
0：无\\
1：强制拉低\\
2：强制拉高\\
3：取反\\
(R/W)
\end{reglist}\end{regdesc}
\end{register}


\begin{register}{H}{MCPWM\_FH1\_CFG1\_REG}{0x00A4}\label{regdesc:MCPWMFH1CFG1REG}
\regfield{(reserved)}{27}{5}{000000000000000000000000000}%
\regfield{MCPWM\_FH1\_FORCE\_OST}{1}{4}{{0}}%
\regfield{MCPWM\_FH1\_FORCE\_CBC}{1}{3}{{0}}%
\regfield{MCPWM\_FH1\_CBCPULSE}{2}{1}{{0}}%
\regfield{MCPWM\_FH1\_CLR\_OST}{1}{0}{{0}}%
\reglabel{Reset}\regnewline%
\begin{regdesc}\begin{reglist}
\label{fielddesc:MCPWMFH1CLROST}\item [MCPWM\_FH1\_CLR\_OST] 置 1 清除正在进行的一次性模式操作。 (R/W)
\label{fielddesc:MCPWMFH1CBCPULSE}\item [MCPWM\_FH1\_CBCPULSE] 设置逐周期模式操作的更新方式。\\
bit0 为 1：发生 TEZ 事件时更新\\
bit1 为 1：发生 TEP 事件时更新\\
(R/W)
\label{fielddesc:MCPWMFH1FORCECBC}\item [MCPWM\_FH1\_FORCE\_CBC] 通过软件将该位取反时，触发逐周期模式操作。 (R/W)
\label{fielddesc:MCPWMFH1FORCEOST}\item [MCPWM\_FH1\_FORCE\_OST] 通过软件将该位取反时，触发一次性模式操作。 (R/W)
\end{reglist}\end{regdesc}
\end{register}


\begin{register}{H}{MCPWM\_FH1\_STATUS\_REG}{0x00A8}\label{regdesc:MCPWMFH1STATUSREG}
\regfield{(reserved)}{30}{2}{000000000000000000000000000000}%
\regfield{MCPWM\_FH1\_OST\_ON}{1}{1}{{0}}%
\regfield{MCPWM\_FH1\_CBC\_ON}{1}{0}{{0}}%
\reglabel{Reset}\regnewline%
\begin{regdesc}\begin{reglist}
\label{fielddesc:MCPWMFH1CBCON}\item [MCPWM\_FH1\_CBC\_ON] 表示是否正在进行逐周期模式的操作。此字段由硬件置 1 和清零。\\
0：不存在正在进行的逐周期模式操作\\
1：逐周期模式的操作正在进行\\
(RO)
\label{fielddesc:MCPWMFH1OSTON}\item [MCPWM\_FH1\_OST\_ON] 表示是否正在进行一次性模式的操作。此字段由硬件置 1 和清零。\\
0：不存在正在进行的一次性模式操作\\
1：一次性模式的操作正在进行\\
(RO)
\end{reglist}\end{regdesc}
\end{register}


\begin{register}{H}{MCPWM\_FH2\_STATUS\_REG}{0x00E0}\label{regdesc:MCPWMFH2STATUSREG}
\regfield{(reserved)}{30}{2}{000000000000000000000000000000}%
\regfield{MCPWM\_FH2\_OST\_ON}{1}{1}{{0}}%
\regfield{MCPWM\_FH2\_CBC\_ON}{1}{0}{{0}}%
\reglabel{Reset}\regnewline%
\begin{regdesc}\begin{reglist}
\label{fielddesc:MCPWMFH2CBCON}\item [MCPWM\_FH2\_CBC\_ON] 表示是否正在进行逐周期模式的操作。此字段由硬件置 1 和清零。\\
0：不存在正在进行的逐周期模式操作\\
1：逐周期模式的操作正在进行\\
(RO)
\label{fielddesc:MCPWMFH2OSTON}\item [MCPWM\_FH2\_OST\_ON] 表示是否正在进行一次性模式的操作。此字段由硬件置 1 和清零。\\
0：不存在正在进行的一次性模式操作\\
1：一次性模式的操作正在进行\\
(RO)
\end{reglist}\end{regdesc}
\end{register}


\begin{register}{H}{MCPWM\_GEN2\_STMP\_CFG\_REG}{0x00AC}\label{regdesc:MCPWMGEN2STMPCFGREG}
\regfield{(reserved)}{22}{10}{0000000000000000000000}%
\regfield{MCPWM\_GEN2\_B\_SHDW\_FULL}{1}{9}{{0}}%
\regfield{MCPWM\_GEN2\_A\_SHDW\_FULL}{1}{8}{{0}}%
\regfield{MCPWM\_GEN2\_B\_UPMETHOD}{4}{4}{{0}}%
\regfield{MCPWM\_GEN2\_A\_UPMETHOD}{4}{0}{{0}}%
\reglabel{Reset}\regnewline%
\begin{regdesc}\begin{reglist}
\label{fielddesc:MCPWMGEN2AUPMETHOD}\item [MCPWM\_GEN2\_A\_UPMETHOD] 生成器 2 时间戳寄存器 A 有效寄存器的更新方式。\\
所有 bit 值 为 0：立即更新\\
bit0 为 1：发生 TEZ 事件时更新\\
bit1 为 1：发生 TEP 事件时更新\\
bit2 为 1：发生同步时间时更新\\
bit3 为 1：关闭更新\\
(R/W)
\label{fielddesc:MCPWMGEN2BUPMETHOD}\item [MCPWM\_GEN2\_B\_UPMETHOD] 生成器 2 时间戳寄存器 B 有效寄存器的更新方式。\\
所有 bit 值 为 0：立即更新\\
bit0 为 1：发生 TEZ 事件时更新\\
bit1 为 1：发生 TEP 事件时更新\\
bit2 为 1：发生同步时间时更新\\
bit3 为 1：关闭更新\\
(R/W)
\label{fielddesc:MCPWMGEN2ASHDWFULL}\item [MCPWM\_GEN2\_A\_SHDW\_FULL] 配置是否在某个特定时间将影子寄存器中的值写入对应的有效寄存器中。此字段由硬件置 1 或复位。\\
0：将影子寄存器中最新的值写入有效寄存器 A 中\\
1：将值写入 PWM 生成器 2 时间戳寄存器 A 的影子寄存器中，等待传输给有效寄存器 A。\\
(R/SC/WTC)
\label{fielddesc:MCPWMGEN2BSHDWFULL}\item [MCPWM\_GEN2\_B\_SHDW\_FULL] 配置是否在某个特定时间将影子寄存器中的值写入对应的有效寄存器中。此字段由硬件置 1 或复位。\\
0：将影子寄存器中最新的值写入有效寄存器 B 中\\
1：将值写入 PWM 生成器 2 时间戳寄存器 B 的影子寄存器中，等待传输给有效寄存器 B。\\
(R/SC/WTC)
\end{reglist}\end{regdesc}
\end{register}


\begin{register}{H}{MCPWM\_GEN2\_STMP\_A\_REG}{0x00B0}\label{regdesc:MCPWMGEN2STMPAREG}
\regfield{(reserved)}{16}{16}{0000000000000000}%
\regfield{MCPWM\_GEN2\_A}{16}{0}{{0}}%
\reglabel{Reset}\regnewline%
\begin{regdesc}\begin{reglist}
\label{fielddesc:MCPWMGEN2A}\item [MCPWM\_GEN2\_A] PWM 生成器 2 时间戳 A 的影子寄存器。 (R/W)
\end{reglist}\end{regdesc}
\end{register}


\begin{register}{H}{MCPWM\_GEN2\_STMP\_B\_REG}{0x00B4}\label{regdesc:MCPWMGEN2STMPBREG}
\regfield{(reserved)}{16}{16}{0000000000000000}%
\regfield{MCPWM\_GEN2\_B}{16}{0}{{0}}%
\reglabel{Reset}\regnewline%
\begin{regdesc}\begin{reglist}
\label{fielddesc:MCPWMGEN2B}\item [MCPWM\_GEN2\_B] PWM 生成器 2 时间戳 B 的影子寄存器。 (R/W)
\end{reglist}\end{regdesc}
\end{register}


\begin{register}{H}{MCPWM\_GEN2\_CFG0\_REG}{0x00B8}\label{regdesc:MCPWMGEN2CFG0REG}
\regfield{(reserved)}{22}{10}{0000000000000000000000}%
\regfield{MCPWM\_GEN2\_T1\_SEL}{3}{7}{{0}}%
\regfield{MCPWM\_GEN2\_T0\_SEL}{3}{4}{{0}}%
\regfield{MCPWM\_GEN2\_CFG\_UPMETHOD}{4}{0}{{0}}%
\reglabel{Reset}\regnewline%
\begin{regdesc}\begin{reglist}
\label{fielddesc:MCPWMGEN2CFGUPMETHOD}\item [MCPWM\_GEN2\_CFG\_UPMETHOD] PWM 生成器 2 有效配置寄存器的更新方式。\\
所有 bit 值为 0：立即更新\\
bit0 为 1：发生 TEZ 事件时更新\\
bit1 为 1：发生 TEP 事件时更新\\
bit2 为 1：发生同 步时间时更新\\
bit3 为 1：关闭更新\\
(R/W)
\label{fielddesc:MCPWMGEN2T0SEL}\item [MCPWM\_GEN2\_T0\_SEL] 设置 PWM 操作器 2 event\_t0 信号源，立即生效。\\
0：fault\_event0\\
1：fault\_event1\\
2：fault\_event2\\
3：sync\_taken\\
4：无\\
(R/W)
\label{fielddesc:MCPWMGEN2T1SEL}\item [MCPWM\_GEN2\_T1\_SEL] 设置 PWM 操作器 2 event\_t1 信号源， 立即生效。\\
0：fault\_event0\\
1：fault\_event1\\
2：fault\_event2\\
3：sync\_taken\\
4：无\\
(R/W)
\end{reglist}\end{regdesc}
\end{register}


\begin{register}{H}{MCPWM\_GEN2\_FORCE\_REG}{0x00BC}\label{regdesc:MCPWMGEN2FORCEREG}
\regfield{(reserved)}{16}{16}{0000000000000000}%
\regfield{MCPWM\_GEN2\_B\_NCIFORCE\_MODE}{2}{14}{{0}}%
\regfield{MCPWM\_GEN2\_B\_NCIFORCE}{1}{13}{{0}}%
\regfield{MCPWM\_GEN2\_A\_NCIFORCE\_MODE}{2}{11}{{0}}%
\regfield{MCPWM\_GEN2\_A\_NCIFORCE}{1}{10}{{0}}%
\regfield{MCPWM\_GEN2\_B\_CNTUFORCE\_MODE}{2}{8}{{0}}%
\regfield{MCPWM\_GEN2\_A\_CNTUFORCE\_MODE}{2}{6}{{0}}%
\regfield{MCPWM\_GEN2\_CNTUFORCE\_UPMETHOD}{6}{0}{{0x20}}%
\reglabel{Reset}\regnewline%
\begin{regdesc}\begin{reglist}
\label{fielddesc:MCPWMGEN2CNTUFORCEUPMETHOD}\item [MCPWM\_GEN2\_CNTUFORCE\_UPMETHOD] 配置 PWM 生成器 2 的持续性软件强制事件的更新方式。\\
0：立即更新\\
bit0 为 1：发生 TEZ 事件时更新\\
bit1 为 1：发生 TEP 事件时更新\\
bit2 为 1：发生 TEA 事件时更新\\
bit3 为 1：发生 TEB 事件时更新\\
bit4 为 1：发生同步时间时更新\\
bit5 为 1：关闭更新\\
本文档中的 TEA/B 表示定时器值等于寄存器 A/B 生成的事件\\
(R/W)
\label{fielddesc:MCPWMGEN2ACNTUFORCEMODE}\item [MCPWM\_GEN2\_A\_CNTUFORCE\_MODE] 配置 PWM2A 的持续性即时软件强制事件。\\
0：关闭\\
1：拉低\\
2：拉高\\
3：关闭\\
 (R/W)
\label{fielddesc:MCPWMGEN2BCNTUFORCEMODE}\item [MCPWM\_GEN2\_B\_CNTUFORCE\_MODE] 配置 PWM2B 的持续性即时软件强制事件。\\
0：关闭\\
1：拉低\\
2：拉高\\
3：关闭\\
 (R/W)
\label{fielddesc:MCPWMGEN2ANCIFORCE}\item [MCPWM\_GEN2\_A\_NCIFORCE] 配置是否触发 PWM2A 上的非连续即时软件强制事件。\\
0：无效\\
1：触发 PWM2A 上的非连续即时软件强制事件\\
(R/W)
\item [见下页...]
\end{reglist}\end{regdesc}
\end{register}
\addtocounter{Regfloat}{-1}
\begin{register}{H}{MCPWM\_GEN2\_FORCE\_REG}{0x0034}
\begin{regdesc}\begin{reglist}
\item [...接上页]
\label{fielddesc:MCPWMGEN2ANCIFORCEMODE}\item [MCPWM\_GEN2\_A\_NCIFORCE\_MODE] 配置 PWM2A 的非持续性即时软件强制事件。\\
0：关闭\\
1：拉低\\
2：拉高\\
3：关闭\\
(R/W)
\label{fielddesc:MCPWMGEN2BNCIFORCE}\item [MCPWM\_GEN2\_B\_NCIFORCE] 配置是否触发 PWM2B 上的非连续即时软件强制事件。\\
0：无效\\
1：触发 PWM2B 上的非连续即时软件强制事件\\
(R/W)
\label{fielddesc:MCPWMGEN2BNCIFORCEMODE}\item [MCPWM\_GEN2\_B\_NCIFORCE\_MODE] 配置 PWM2B 的非持续性即时软件强制模式。\\
0：关闭\\
1：拉低\\
2：拉高\\
3：关闭\\
(R/W)
\end{reglist}\end{regdesc}
\end{register}


\begin{register}{H}{MCPWM\_GEN2\_A\_REG}{0x00C0}\label{regdesc:MCPWMGEN2AREG}
\regfield{(reserved)}{8}{24}{00000000}%
\regfield{MCPWM\_GEN2\_A\_DT1}{2}{22}{{0}}%
\regfield{MCPWM\_GEN2\_A\_DT0}{2}{20}{{0}}%
\regfield{MCPWM\_GEN2\_A\_DTEB}{2}{18}{{0}}%
\regfield{MCPWM\_GEN2\_A\_DTEA}{2}{16}{{0}}%
\regfield{MCPWM\_GEN2\_A\_DTEP}{2}{14}{{0}}%
\regfield{MCPWM\_GEN2\_A\_DTEZ}{2}{12}{{0}}%
\regfield{MCPWM\_GEN2\_A\_UT1}{2}{10}{{0}}%
\regfield{MCPWM\_GEN2\_A\_UT0}{2}{8}{{0}}%
\regfield{MCPWM\_GEN2\_A\_UTEB}{2}{6}{{0}}%
\regfield{MCPWM\_GEN2\_A\_UTEA}{2}{4}{{0}}%
\regfield{MCPWM\_GEN2\_A\_UTEP}{2}{2}{{0}}%
\regfield{MCPWM\_GEN2\_A\_UTEZ}{2}{0}{{0}}%
\reglabel{Reset}\regnewline%
\begin{regdesc}\begin{reglist}
\label{fielddesc:MCPWMGEN2AUTEZ}\item [MCPWM\_GEN2\_A\_UTEZ] 配置定时器递增时，TEZ 在 PWM2A 上触发的操作。\\
0：无\\
1：拉低\\
2：拉高\\
3：取反\\
(R/W)
\label{fielddesc:MCPWMGEN2AUTEP}\item [MCPWM\_GEN2\_A\_UTEP] 配置定时器递增时，TEP 在 PWM2A 上触发的操作。\\
0：无\\
1：拉低\\
2：拉高\\
3：取反\\
(R/W)
\label{fielddesc:MCPWMGEN2AUTEA}\item [MCPWM\_GEN2\_A\_UTEA] 配置定时器递增时，TEA 在 PWM2A 上触发的操作。\\
0：无\\
1：拉低\\
2：拉高\\
3：取反\\
(R/W)
\label{fielddesc:MCPWMGEN2AUTEB}\item [MCPWM\_GEN2\_A\_UTEB] 配置定时器递增时，TEB 在 PWM2A 上触发的操作。\\
0：无\\
1：拉低\\
2：拉高\\
3：取反\\
(R/W)
\label{fielddesc:MCPWMGEN2AUT0}\item [MCPWM\_GEN2\_A\_UT0] 配置定时器递增时，event\_t0 在 PWM2A 上触发的操作。\\
0：无\\
1：拉低\\
2：拉高\\
3：取反\\
(R/W)
\item [见下页...]
\end{reglist}\end{regdesc}
\end{register}
\addtocounter{Regfloat}{-1}
\begin{register}{H}{MCPWM\_GEN2\_A\_REG}{0x0034}
\begin{regdesc}\begin{reglist}
\item [...接上页]
\label{fielddesc:MCPWMGEN2AUT1}\item [MCPWM\_GEN2\_A\_UT1] 配置定时器递增时，event\_t1 在 PWM2A 上触发的操作。\\
0：无\\
1：拉低\\
2：拉高\\
3：取反\\
(R/W)
\label{fielddesc:MCPWMGEN2ADTEZ}\item [MCPWM\_GEN2\_A\_DTEZ] 配置定时器递减时，TEZ 在 PWM2A 上触发的操作。\\
0：无\\
1：拉低\\
2：拉高\\
3：取反\\
(R/W)
\label{fielddesc:MCPWMGEN2ADTEP}\item [MCPWM\_GEN2\_A\_DTEP] 配置定时器递减时，TEP 在 PWM2A 上触发的操作。\\
0：无\\
1：拉低\\
2：拉高\\
3：取反\\
(R/W)
\label{fielddesc:MCPWMGEN2ADTEA}\item [MCPWM\_GEN2\_A\_DTEA] 配置定时器递减时，TEA 在 PWM2A 上触发的操作。\\
0：无\\
1：拉低\\
2：拉高\\
3：取反\\
(R/W)
\label{fielddesc:MCPWMGEN2ADTEB}\item [MCPWM\_GEN2\_A\_DTEB] 配置定时器递减时，TEB 在 PWM2A 上触发的操作。\\
0：无\\
1：拉低\\
2：拉高\\
3：取反\\
(R/W)
\label{fielddesc:MCPWMGEN2ADT0}\item [MCPWM\_GEN2\_A\_DT0] 配置定时器递减时，event\_t0 在 PWM2A 上触发的操作。\\
0：无\\
1：拉低\\
2：拉高\\
3：取反\\
(R/W)
\label{fielddesc:MCPWMGEN2ADT1}\item [MCPWM\_GEN2\_A\_DT1] 配置定时器递减时，event\_t1 在 PWM2A 上触发的操作。\\
0：无\\
1：拉低\\
2：拉高\\
3：取反\\
(R/W)
\end{reglist}\end{regdesc}
\end{register}


\begin{register}{H}{MCPWM\_GEN2\_B\_REG}{0x00C4}\label{regdesc:MCPWMGEN2BREG}
\regfield{(reserved)}{8}{24}{00000000}%
\regfield{MCPWM\_GEN2\_B\_DT1}{2}{22}{{0}}%
\regfield{MCPWM\_GEN2\_B\_DT0}{2}{20}{{0}}%
\regfield{MCPWM\_GEN2\_B\_DTEB}{2}{18}{{0}}%
\regfield{MCPWM\_GEN2\_B\_DTEA}{2}{16}{{0}}%
\regfield{MCPWM\_GEN2\_B\_DTEP}{2}{14}{{0}}%
\regfield{MCPWM\_GEN2\_B\_DTEZ}{2}{12}{{0}}%
\regfield{MCPWM\_GEN2\_B\_UT1}{2}{10}{{0}}%
\regfield{MCPWM\_GEN2\_B\_UT0}{2}{8}{{0}}%
\regfield{MCPWM\_GEN2\_B\_UTEB}{2}{6}{{0}}%
\regfield{MCPWM\_GEN2\_B\_UTEA}{2}{4}{{0}}%
\regfield{MCPWM\_GEN2\_B\_UTEP}{2}{2}{{0}}%
\regfield{MCPWM\_GEN2\_B\_UTEZ}{2}{0}{{0}}%
\reglabel{Reset}\regnewline%
\begin{regdesc}\begin{reglist}
\label{fielddesc:MCPWMGEN2BUTEZ}\item [MCPWM\_GEN2\_B\_UTEZ] 配置定时器递增时，TEZ 在 PWM2B 上触发的操作。\\
0：无\\
1：拉低\\
2：拉高\\
3：取反 \\
(R/W)
\label{fielddesc:MCPWMGEN2BUTEP}\item [MCPWM\_GEN2\_B\_UTEP] 配置定时器递增时，TEP 在 PWM2B 上触发的操作。\\
0：无\\
1：拉低\\
2：拉高\\
3：取反 \\
(R/W)
\label{fielddesc:MCPWMGEN2BUTEA}\item [MCPWM\_GEN2\_B\_UTEA] 配置定时器递增时，TEA 在 PWM2B 上触发的操作。\\
0：无\\
1：拉低\\
2：拉高\\
3：取反 \\
(R/W)
\label{fielddesc:MCPWMGEN2BUTEB}\item [MCPWM\_GEN2\_B\_UTEB] 配置定时器递增时，TEB 在 PWM2B 上触发的操作。\\
0：无\\
1：拉低\\
2：拉高\\
3：取反 \\
(R/W)
\label{fielddesc:MCPWMGEN2BUT0}\item [MCPWM\_GEN2\_B\_UT0] 配置定时器递增时，event\_t0 在 PWM2B 上触发的操作。\\
0：无\\
1：拉低\\
2：拉高\\
3：取反 \\
(R/W)
\label{fielddesc:MCPWMGEN2BUT1}\item [MCPWM\_GEN2\_B\_UT1] 配置定时器递增时，event\_t1 在 PWM2B 上触发的操作。\\
0：无\\
1：拉低\\
2：拉高\\
3：取反 \\
(R/W)
\item [见下页...]
\end{reglist}\end{regdesc}
\end{register}
\addtocounter{Regfloat}{-1}
\begin{register}{H}{MCPWM\_GEN2\_B\_REG}{0x0034}
\begin{regdesc}\begin{reglist}
\item [...接上页]
\label{fielddesc:MCPWMGEN2BDTEZ}\item [MCPWM\_GEN2\_B\_DTEZ] 配置定时器递减时，TEZ 在 PWM2B 上触发的操作。\\
0：无\\
1：拉低\\
2：拉高\\
3：取反 \\
(R/W)
\label{fielddesc:MCPWMGEN2BDTEP}\item [MCPWM\_GEN2\_B\_DTEP] 配置定时器递减时，TEP 在 PWM2B 上触发的操作。\\
0：无\\
1：拉低\\
2：拉高\\
3：取反 \\
(R/W)
\label{fielddesc:MCPWMGEN2BDTEA}\item [MCPWM\_GEN2\_B\_DTEA] 配置定时器递减时，TEA 在 PWM2B 上触发的操作。\\
0：无\\
1：拉低\\
2：拉高\\
3：取反 \\
(R/W)
\label{fielddesc:MCPWMGEN2BDTEB}\item [MCPWM\_GEN2\_B\_DTEB] 配置定时器递减时，TEB 在 PWM2B 上触发的操作。\\
0：无\\
1：拉低\\
2：拉高\\
3：取反 \\
(R/W)
\label{fielddesc:MCPWMGEN2BDT0}\item [MCPWM\_GEN2\_B\_DT0] 配置定时器递减时，event\_t0 在 PWM2B 上触发的操作。\\
0：无\\
1：拉低\\
2：拉高\\
3：取反 \\
(R/W)
\label{fielddesc:MCPWMGEN2BDT1}\item [MCPWM\_GEN2\_B\_DT1] 配置定时器递减时，event\_t1 在 PWM2B 上触发的操作。\\
0：无\\
1：拉低\\
2：拉高\\
3：取反 \\
(R/W)
\end{reglist}\end{regdesc}
\end{register}


\begin{register}{H}{MCPWM\_DT2\_CFG\_REG}{0x00C8}\label{regdesc:MCPWMDT2CFGREG}
\regfield{(reserved)}{14}{18}{00000000000000}%
\regfield{MCPWM\_DT2\_CLK\_SEL}{1}{17}{{0}}%
\regfield{MCPWM\_DT2\_B\_OUTBYPASS}{1}{16}{{1}}%
\regfield{MCPWM\_DT2\_A\_OUTBYPASS}{1}{15}{{1}}%
\regfield{MCPWM\_DT2\_FED\_OUTINVERT}{1}{14}{{0}}%
\regfield{MCPWM\_DT2\_RED\_OUTINVERT}{1}{13}{{0}}%
\regfield{MCPWM\_DT2\_FED\_INSEL}{1}{12}{{0}}%
\regfield{MCPWM\_DT2\_RED\_INSEL}{1}{11}{{0}}%
\regfield{MCPWM\_DT2\_B\_OUTSWAP}{1}{10}{{0}}%
\regfield{MCPWM\_DT2\_A\_OUTSWAP}{1}{9}{{0}}%
\regfield{MCPWM\_DT2\_DEB\_MODE}{1}{8}{{0}}%
\regfield{MCPWM\_DT2\_RED\_UPMETHOD}{4}{4}{{0}}%
\regfield{MCPWM\_DT2\_FED\_UPMETHOD}{4}{0}{{0}}%
\reglabel{Reset}\regnewline%
\begin{regdesc}\begin{reglist}
\label{fielddesc:MCPWMDT2FEDUPMETHOD}\item [MCPWM\_DT2\_FED\_UPMETHOD] FED(下降沿延迟)有效寄存器的更新方式。\\
0：立即更新\\
bit0 为 1：发生 TEZ 事件时更新\\
bit1 为 1：发生 TEP 事件时更新\\
bit2 为 1：发生同步事件时更新\\
bit3 为 1：关闭更新\\
(R/W)
\label{fielddesc:MCPWMDT2REDUPMETHOD}\item [MCPWM\_DT2\_RED\_UPMETHOD] RED(上升沿延迟)有效寄存器的更新方式。\\
0：立即更新\\
bit0 为 1：发生 TEZ 事件时更新\\
bit1 为 1：发生 TEP 事件时更新\\
bit2 为 1：发生同步事件时更新\\
bit3 为 1：关闭更新\\
(R/W)
\label{fielddesc:MCPWMDT2DEBMODE}\item [MCPWM\_DT2\_DEB\_MODE] 配置表 \ref{tab:mcpwm-s-dt}中的 S8 开关，详细配置信息请参考表 \ref{tab:mcpwm-mod-dt}。
(R/W)
\label{fielddesc:MCPWMDT2AOUTSWAP}\item [MCPWM\_DT2\_A\_OUTSWAP] 配置表 \ref{tab:mcpwm-s-dt} 中的 S6 开关，详细配置信息请参考表 \ref{tab:mcpwm-mod-dt}。 (R/W)
\label{fielddesc:MCPWMDT2BOUTSWAP}\item [MCPWM\_DT2\_B\_OUTSWAP] 配置表 \ref{tab:mcpwm-s-dt} 中的 S7 开关，详细配置信息请参考表 \ref{tab:mcpwm-mod-dt}。 (R/W)
\label{fielddesc:MCPWMDT2REDINSEL}\item [MCPWM\_DT2\_RED\_INSEL] 配置表 \ref{tab:mcpwm-s-dt} 中的 S4 开关，详细配置信息请参考表 \ref{tab:mcpwm-mod-dt}。 (R/W)
\label{fielddesc:MCPWMDT2FEDINSEL}\item [MCPWM\_DT2\_FED\_INSEL] 配置表 \ref{tab:mcpwm-s-dt} 中的 S5 开关，详细配置信息请参考表 \ref{tab:mcpwm-mod-dt}。 (R/W)
\label{fielddesc:MCPWMDT2REDOUTINVERT}\item [MCPWM\_DT2\_RED\_OUTINVERT] 配置表 \ref{tab:mcpwm-s-dt} 中的 S2 开关，详细配置信息请参考表 \ref{tab:mcpwm-mod-dt}。 (R/W)
\label{fielddesc:MCPWMDT2FEDOUTINVERT}\item [MCPWM\_DT2\_FED\_OUTINVERT] 配置表 \ref{tab:mcpwm-s-dt} 中的 S3 开关，详细配置信息请参考表 \ref{tab:mcpwm-mod-dt}。 (R/W)
\label{fielddesc:MCPWMDT2AOUTBYPASS}\item [MCPWM\_DT2\_A\_OUTBYPASS] 配置表 \ref{tab:mcpwm-s-dt} 中的 S1 开关，详细配置信息请参考表 \ref{tab:mcpwm-mod-dt}。 (R/W)
\label{fielddesc:MCPWMDT2BOUTBYPASS}\item [MCPWM\_DT2\_B\_OUTBYPASS] 配置表 \ref{tab:mcpwm-s-dt} 中的 S0 开关，详细配置信息请参考表 \ref{tab:mcpwm-mod-dt}。 (R/W)
\label{fielddesc:MCPWMDT2CLKSEL}\item [MCPWM\_DT2\_CLK\_SEL] 配置死区时间生成器时钟。\\
0：PWM\_CLK\\
1：PT\_CLK\\
(R/W)
\end{reglist}\end{regdesc}
\end{register}


\begin{register}{H}{MCPWM\_DT2\_FED\_CFG\_REG}{0x00CC}\label{regdesc:MCPWMDT2FEDCFGREG}
\regfield{(reserved)}{16}{16}{0000000000000000}%
\regfield{MCPWM\_DT2\_FED}{16}{0}{{0}}%
\reglabel{Reset}\regnewline%
\begin{regdesc}\begin{reglist}
\label{fielddesc:MCPWMDT2FED}\item [MCPWM\_DT2\_FED] FED 影子寄存器。 (R/W)
\end{reglist}\end{regdesc}
\end{register}


\begin{register}{H}{MCPWM\_DT2\_RED\_CFG\_REG}{0x00D0}\label{regdesc:MCPWMDT2REDCFGREG}
\regfield{(reserved)}{16}{16}{0000000000000000}%
\regfield{MCPWM\_DT2\_RED}{16}{0}{{0}}%
\reglabel{Reset}\regnewline%
\begin{regdesc}\begin{reglist}
\label{fielddesc:MCPWMDT2RED}\item [MCPWM\_DT2\_RED] RED 影子寄存器。 (R/W)
\end{reglist}\end{regdesc}
\end{register}


\begin{register}{H}{MCPWM\_CARRIER2\_CFG\_REG}{0x00D4}\label{regdesc:MCPWMCARRIER2CFGREG}
\regfield{(reserved)}{18}{14}{000000000000000000}%
\regfield{MCPWM\_CARRIER2\_IN\_INVERT}{1}{13}{{0}}%
\regfield{MCPWM\_CARRIER2\_OUT\_INVERT}{1}{12}{{0}}%
\regfield{MCPWM\_CARRIER2\_OSHTWTH}{4}{8}{{0}}%
\regfield{MCPWM\_CARRIER2\_DUTY}{3}{5}{{0}}%
\regfield{MCPWM\_CARRIER2\_PRESCALE}{4}{1}{{0}}%
\regfield{MCPWM\_CARRIER2\_EN}{1}{0}{{0}}%
\reglabel{Reset}\regnewline%
\begin{regdesc}\begin{reglist}
\label{fielddesc:MCPWMCARRIER2EN}\item [MCPWM\_CARRIER2\_EN] 配置是否使能载波 2 功能。\\
0：绕过载波 2 \\
1：使能载波 2 功能\\
(R/W)
\label{fielddesc:MCPWMCARRIER2PRESCALE}\item [MCPWM\_CARRIER2\_PRESCALE] 配置 PWM 载 波2 时钟 (PC\_CLK) 的预分频值。PC\_CLK 周期 = PWM\_CLK 周期 * (PWM\_CARRIER0\_PRESCALE + 1)。 (R/W)
\label{fielddesc:MCPWMCARRIER2DUTY}\item [MCPWM\_CARRIER2\_DUTY] 配置载波占空比。占空比 = PWM\_CARRIER0\_DUTY / 8。 (R/W)
\label{fielddesc:MCPWMCARRIER2OSHTWTH}\item [MCPWM\_CARRIER2\_OSHTWTH] 配置载波第一个脉冲的宽度，单位为载波周期。 (R/W)
\label{fielddesc:MCPWMCARRIER2OUTINVERT}\item [MCPWM\_CARRIER2\_OUT\_INVERT] 配置是否将此模块的 PWM2A 和 PWM2B 输出反相。\\
0：无效\\
1：反相\\
(R/W)
\label{fielddesc:MCPWMCARRIER2ININVERT}\item [MCPWM\_CARRIER2\_IN\_INVERT] 配置是否将此模块的 PWM2A 和 PWM2B 输入反相。\\
0：无效\\
1：反相\\
(R/W)
\end{reglist}\end{regdesc}
\end{register}


\begin{register}{H}{MCPWM\_FH2\_CFG0\_REG}{0x00D8}\label{regdesc:MCPWMFH2CFG0REG}
\regfield{(reserved)}{8}{24}{00000000}%
\regfield{MCPWM\_FH2\_B\_OST\_U}{2}{22}{{0}}%
\regfield{MCPWM\_FH2\_B\_OST\_D}{2}{20}{{0}}%
\regfield{MCPWM\_FH2\_B\_CBC\_U}{2}{18}{{0}}%
\regfield{MCPWM\_FH2\_B\_CBC\_D}{2}{16}{{0}}%
\regfield{MCPWM\_FH2\_A\_OST\_U}{2}{14}{{0}}%
\regfield{MCPWM\_FH2\_A\_OST\_D}{2}{12}{{0}}%
\regfield{MCPWM\_FH2\_A\_CBC\_U}{2}{10}{{0}}%
\regfield{MCPWM\_FH2\_A\_CBC\_D}{2}{8}{{0}}%
\regfield{MCPWM\_FH2\_F0\_OST}{1}{7}{{0}}%
\regfield{MCPWM\_FH2\_F1\_OST}{1}{6}{{0}}%
\regfield{MCPWM\_FH2\_F2\_OST}{1}{5}{{0}}%
\regfield{MCPWM\_FH2\_SW\_OST}{1}{4}{{0}}%
\regfield{MCPWM\_FH2\_F0\_CBC}{1}{3}{{0}}%
\regfield{MCPWM\_FH2\_F1\_CBC}{1}{2}{{0}}%
\regfield{MCPWM\_FH2\_F2\_CBC}{1}{1}{{0}}%
\regfield{MCPWM\_FH2\_SW\_CBC}{1}{0}{{0}}%
\reglabel{Reset}\regnewline%
\begin{regdesc}\begin{reglist}
\label{fielddesc:MCPWMFH2SWCBC}\item [MCPWM\_FH2\_SW\_CBC] 软件强制逐周期模式操作的使能寄存器。0：关闭。1：使能。 (R/W)
\label{fielddesc:MCPWMFH2F2CBC}\item [MCPWM\_FH2\_F2\_CBC] 设置 fault\_event2 触发逐周期模式操作。0：关闭。1：使能。 (R/W)
\label{fielddesc:MCPWMFH2F1CBC}\item [MCPWM\_FH2\_F1\_CBC] 设置 fault\_event1 触发逐周期模式操作。0：关闭。1：使能。 (R/W)
\label{fielddesc:MCPWMFH2F0CBC}\item [MCPWM\_FH2\_F0\_CBC] 设置 fault\_event0 触发逐周期模式操作。0：关闭。1：使能。 (R/W)
\label{fielddesc:MCPWMFH2SWOST}\item [MCPWM\_FH2\_SW\_OST] 软件强制一次性模式操作的使能寄存器。0：关闭。1：使能。 (R/W)
\label{fielddesc:MCPWMFH2F2OST}\item [MCPWM\_FH2\_F2\_OST] 设置 fault\_event2 触发一次性模式操作。0：关闭。1：使能。 (R/W)
\label{fielddesc:MCPWMFH2F1OST}\item [MCPWM\_FH2\_F1\_OST] 设置 fault\_event1 触发一次性模式操作。0：关闭。1：使能。 (R/W)
\label{fielddesc:MCPWMFH2F0OST}\item [MCPWM\_FH2\_F0\_OST] 设置 fault\_event0 触发一次性模式操作。0：关闭。1：使能。 (R/W)
\label{fielddesc:MCPWMFH2ACBCD}\item [MCPWM\_FH2\_A\_CBC\_D] 发生故障事件并且定时器递减时，PWM2A 上的逐周期模式操作。0：无。1：强制拉低。2：强制拉高。3：取反。 (R/W)
\label{fielddesc:MCPWMFH2ACBCU}\item [MCPWM\_FH2\_A\_CBC\_U] 发生故障事件并且定时器递增时，PWM2A 上的逐周期模式操作。0：无。1：强制拉低。2：强制拉高。3：取反。 (R/W)
\label{fielddesc:MCPWMFH2AOSTD}\item [MCPWM\_FH2\_A\_OST\_D] 发生故障事件并且定时器递减时，PWM2A 上的一次性模式操作。0：无。1：强制拉低。2：强制拉高。3：取反。 (R/W)
\label{fielddesc:MCPWMFH2AOSTU}\item [MCPWM\_FH2\_A\_OST\_U] 发生故障事件并且定时器递增时，PWM2A 上的一次性模式操作。0：无。1：强制拉低。2：强制拉高。3：取反。 (R/W)
\label{fielddesc:MCPWMFH2BCBCD}\item [MCPWM\_FH2\_B\_CBC\_D] 发生故障事件并且定时器递减时，PWM2B 上的逐周期模式操作。0：无。1：强制拉低。2：强制拉高。3：取反。 (R/W)
\label{fielddesc:MCPWMFH2BCBCU}\item [MCPWM\_FH2\_B\_CBC\_U] 发生故障事件并且定时器递增时，PWM2B 上的逐周期模式操作。0：无。1：强制拉低。2：强制拉高。3：取反。 (R/W)
\label{fielddesc:MCPWMFH2BOSTD}\item [MCPWM\_FH2\_B\_OST\_D] 发生故障事件并且定时器递减时，PWM2B 上的一次性模式操作。0：无。1：强制拉低。2：强制拉高。3：取反。 (R/W)
\label{fielddesc:MCPWMFH2BOSTU}\item [MCPWM\_FH2\_B\_OST\_U] 发生故障事件并且定时器递增时，PWM2B 上的一次性模式操作。0：无。1：强制拉低。2：强制拉高。3：取反。 (R/W)
\end{reglist}\end{regdesc}
\end{register}


\begin{register}{H}{MCPWM\_FH2\_CFG1\_REG}{0x00DC}\label{regdesc:MCPWMFH2CFG1REG}
\regfield{(reserved)}{27}{5}{000000000000000000000000000}%
\regfield{MCPWM\_TZ2\_FORCE\_OST}{1}{4}{{0}}%
\regfield{MCPWM\_TZ2\_FORCE\_CBC}{1}{3}{{0}}%
\regfield{MCPWM\_TZ2\_CBCPULSE}{2}{1}{{0}}%
\regfield{MCPWM\_FH2\_CLR\_OST}{1}{0}{{0}}%
\reglabel{Reset}\regnewline%
\begin{regdesc}\begin{reglist}
\label{fielddesc:MCPWMFH2CLROST}\item [MCPWM\_FH2\_CLR\_OST] 配置是否使上升沿清除正在进行的一次性模式的操作。\\
0: 不清除\\
1: 清除\\
(R/W)
\label{fielddesc:MCPWMTZ2CBCPULSE}\item [MCPWM\_TZ2\_CBCPULSE] 配置逐周期模式的更新方式。\\
bit0 为 1：发生 TEZ 事件时\\
bit1 为 1：发生 TEP 事件时\\
(R/W)
\label{fielddesc:MCPWMTZ2FORCECBC}\item [MCPWM\_TZ2\_FORCE\_CBC] 配置是否触发逐周期模式操作。\\
0：无效\\
1：触发逐周期模式操作\\
(R/W)
\label{fielddesc:MCPWMTZ2FORCEOST}\item [MCPWM\_TZ2\_FORCE\_OST] 配置是否触发一次性模式操作。\\
0：无效\\
1：触发一次性模式操作\\
(R/W)
\end{reglist}\end{regdesc}
\end{register}


\begin{register}{H}{MCPWM\_FAULT\_DETECT\_REG}{0x00E4}\label{regdesc:MCPWMFAULTDETECTREG}
\regfield{(reserved)}{23}{9}{00000000000000000000000}%
\regfield{MCPWM\_EVENT\_F2}{1}{8}{{0}}%
\regfield{MCPWM\_EVENT\_F1}{1}{7}{{0}}%
\regfield{MCPWM\_EVENT\_F0}{1}{6}{{0}}%
\regfield{MCPWM\_F2\_POLE}{1}{5}{{0}}%
\regfield{MCPWM\_F1\_POLE}{1}{4}{{0}}%
\regfield{MCPWM\_F0\_POLE}{1}{3}{{0}}%
\regfield{MCPWM\_F2\_EN}{1}{2}{{0}}%
\regfield{MCPWM\_F1\_EN}{1}{1}{{0}}%
\regfield{MCPWM\_F0\_EN}{1}{0}{{0}}%
\reglabel{Reset}\regnewline%
\begin{regdesc}\begin{reglist}
\label{fielddesc:MCPWMF0EN}\item [MCPWM\_F0\_EN] 配置是否使能生成 fault\_event0。\\
0：无效\\
1：使能\\
(R/W)
\label{fielddesc:MCPWMF1EN}\item [MCPWM\_F1\_EN] 配置是否使能生成 fault\_event1。\\
0：无效\\
1：使能\\
(R/W)
\label{fielddesc:MCPWMF2EN}\item [MCPWM\_F2\_EN] 配置是否使能生成 fault\_event2。\\
0：无效\\
1：使能\\
(R/W)
\label{fielddesc:MCPWMF0POLE}\item [MCPWM\_F0\_POLE] 配置来自 GPIO 矩阵的 FAULT0 信号源触发 fault\_event0 时极性。\\
0：低电平触发\\
1：高电平触发\\
(R/W)
\label{fielddesc:MCPWMF1POLE}\item [MCPWM\_F1\_POLE] 配置来自 GPIO 矩阵的 FAULT1 信号源触发 fault\_event1 时极性。\\
0：低电平触发\\
1：高电平触发\\
(R/W)
\label{fielddesc:MCPWMF2POLE}\item [MCPWM\_F2\_POLE] 配置来自 GPIO 矩阵的 FAULT2 信号源触发 fault\_event2 时极性。\\
0：低电平触发\\
1：高电平触发\\
(R/W)
\label{fielddesc:MCPWMEVENTF0}\item [MCPWM\_EVENT\_F0] 表示 fault\_event0 事件的状态。此字段由硬件置 1 和清零。\\
0：fault\_event0 事件停止\\
1：fault\_event0 事件持续\\
(RO)
\label{fielddesc:MCPWMEVENTF1}\item [MCPWM\_EVENT\_F1] 表示 fault\_event1 事件的状态。此字段由硬件置 1 和清零。\\
0：fault\_event1 事件停止\\
1：fault\_event1 事件持续\\
(RO)
\label{fielddesc:MCPWMEVENTF2}\item [MCPWM\_EVENT\_F2] 表示 fault\_event2 事件的状态。此字段由硬件置 1 和清零。\\
0：fault\_event2 事件停止\\
1：fault\_event2 事件持续\\
(RO)
\end{reglist}\end{regdesc}
\end{register}


\begin{register}{H}{MCPWM\_CAP\_TIMER\_CFG\_REG}{0x00E8}\label{regdesc:MCPWMCAPTIMERCFGREG}
\regfield{(reserved)}{26}{6}{00000000000000000000000000}%
\regfield{MCPWM\_CAP\_SYNC\_SW}{1}{5}{{0}}%
\regfield{MCPWM\_CAP\_SYNCI\_SEL}{3}{2}{{0}}%
\regfield{MCPWM\_CAP\_SYNCI\_EN}{1}{1}{{0}}%
\regfield{MCPWM\_CAP\_TIMER\_EN}{1}{0}{{0}}%
\reglabel{Reset}\regnewline%
\begin{regdesc}\begin{reglist}
\label{fielddesc:MCPWMCAPTIMEREN}\item [MCPWM\_CAP\_TIMER\_EN] 配置是否使能捕获定时器在 APB\_CLK 下的递增。\\
0：无效\\
1：使能捕获定时器在 APB\_CLK 下的递增\\
(R/W)
\label{fielddesc:MCPWMCAPSYNCIEN}\item [MCPWM\_CAP\_SYNCI\_EN] 配置是否使能捕获定时器同步。\\
0：无效\\
1：使能捕获定时器同步\\
(R/W)
\label{fielddesc:MCPWMCAPSYNCISEL}\item [MCPWM\_CAP\_SYNCI\_SEL] 配置捕获模块的同步输入。\\
0：无\\
1：定时器 0 的同步输出\\
2：定时器 1 的同步输出\\
3：定时器 2 的同步输出\\
4：来自 GPIO 矩阵的 SYNC0\\
5：来自 GPIO 矩阵的 SYNC1\\
6：来自 GPIO 矩阵的 SYNC2\\
(R/W)
\label{fielddesc:MCPWMCAPSYNCSW}\item [MCPWM\_CAP\_SYNC\_SW] 当 \hyperref[fielddesc:MCPWMCAPSYNCIEN]{MCPWM\_CAP\_SYNCI\_EN} 置位时，配置是否同步捕获定时器，在捕获定时器中写入相位寄存器的值。\\
0：无效\\
1：同步捕获定时器\\
(WT)
\end{reglist}\end{regdesc}
\end{register}


\begin{register}{H}{MCPWM\_CAP\_TIMER\_PHASE\_REG}{0x00EC}\label{regdesc:MCPWMCAPTIMERPHASEREG}
\regfield{MCPWM\_CAP\_TIMER\_PHASE}{32}{0}{{0}}%
\reglabel{Reset}\regnewline%
\begin{regdesc}\begin{reglist}
\label{fielddesc:MCPWMCAPTIMERPHASE}\item [MCPWM\_CAP\_TIMER\_PHASE] 捕获定时器同步操作的相位值。(R/W)
\end{reglist}\end{regdesc}
\end{register}


\begin{register}{H}{MCPWM\_CAP\_CH0\_CFG\_REG}{0x00F0}\label{regdesc:MCPWMCAPCH0CFGREG}
\regfield{(reserved)}{19}{13}{0000000000000000000}%
\regfield{MCPWM\_CAP0\_SW}{1}{12}{{0}}%
\regfield{MCPWM\_CAP0\_IN\_INVERT}{1}{11}{{0}}%
\regfield{MCPWM\_CAP0\_PRESCALE}{8}{3}{{0}}%
\regfield{MCPWM\_CAP0\_MODE}{2}{1}{{0}}%
\regfield{MCPWM\_CAP0\_EN}{1}{0}{{0}}%
\reglabel{Reset}\regnewline%
\begin{regdesc}\begin{reglist}
\label{fielddesc:MCPWMCAP0EN}\item [MCPWM\_CAP0\_EN] 配置是否使能信道 0 上的捕获事件。\\
0：不使能\\
1：使能\\
(R/W)
\label{fielddesc:MCPWMCAP0MODE}\item [MCPWM\_CAP0\_MODE] 配置预分频后信道 0 上的捕获沿。\\
bit0 为 1：使能下降沿捕获。\\
bit1 为 1：使能上升沿捕获。 \\
(R/W)
\label{fielddesc:MCPWMCAP0PRESCALE}\item [MCPWM\_CAP0\_PRESCALE] 配置 CAP0上升沿的预分频值。预分频值 = PWM\_CAP0\_PRESCALE + 1。 (R/W)
\label{fielddesc:MCPWMCAP0ININVERT}\item [MCPWM\_CAP0\_IN\_INVERT] 配置是否在预分频前反相来自 GPIO 矩阵的 CAP0。\\
0：无效\\
1：在预分频前反相来自 GPIO 矩阵的 CAP0\\
(R/W)
\label{fielddesc:MCPWMCAP0SW}\item [MCPWM\_CAP0\_SW] 配置是否触发信道 0 上的软件强制捕获事件。\\
0：不触发\\
1：触发\\
(WT)
\end{reglist}\end{regdesc}
\end{register}


\begin{register}{H}{MCPWM\_CAP\_CH1\_CFG\_REG}{0x00F4}\label{regdesc:MCPWMCAPCH1CFGREG}
\regfield{(reserved)}{19}{13}{0000000000000000000}%
\regfield{MCPWM\_CAP1\_SW}{1}{12}{{0}}%
\regfield{MCPWM\_CAP1\_IN\_INVERT}{1}{11}{{0}}%
\regfield{MCPWM\_CAP1\_PRESCALE}{8}{3}{{0}}%
\regfield{MCPWM\_CAP1\_MODE}{2}{1}{{0}}%
\regfield{MCPWM\_CAP1\_EN}{1}{0}{{0}}%
\reglabel{Reset}\regnewline%
\begin{regdesc}\begin{reglist}
\label{fielddesc:MCPWMCAP1EN}\item [MCPWM\_CAP1\_EN] 配置是否使能信道 1 上的捕获事件。\\
0：不使能\\
1：使能\\
(R/W)
\label{fielddesc:MCPWMCAP1MODE}\item [MCPWM\_CAP1\_MODE] 配置预分频后信道 1 上的捕获沿。\\
bit0 为 1：使能下降沿捕获。\\
bit1 为 1：使能上升沿捕获。 \\
(R/W)
\label{fielddesc:MCPWMCAP1PRESCALE}\item [MCPWM\_CAP1\_PRESCALE] 配置 CAP1 上升沿的预分频值。预分频值 = PWM\_CAP1\_PRESCALE + 1。 (R/W)
\label{fielddesc:MCPWMCAP1ININVERT}\item [MCPWM\_CAP1\_IN\_INVERT] 配置是否在预分频前反相来自 GPIO 矩阵的 CAP1。\\
0：无效\\
1：在预分频前反相来自 GPIO 矩阵的 CAP1\\
(R/W)
\label{fielddesc:MCPWMCAP1SW}\item [MCPWM\_CAP1\_SW] 配置是否触发信道 1 上的软件强制捕获事件。\\
0：不触发\\
1：触发\\
(WT)
\end{reglist}\end{regdesc}
\end{register}


\begin{register}{H}{MCPWM\_CAP\_CH2\_CFG\_REG}{0x00F8}\label{regdesc:MCPWMCAPCH2CFGREG}
\regfield{(reserved)}{19}{13}{0000000000000000000}%
\regfield{MCPWM\_CAP2\_SW}{1}{12}{{0}}%
\regfield{MCPWM\_CAP2\_IN\_INVERT}{1}{11}{{0}}%
\regfield{MCPWM\_CAP2\_PRESCALE}{8}{3}{{0}}%
\regfield{MCPWM\_CAP2\_MODE}{2}{1}{{0}}%
\regfield{MCPWM\_CAP2\_EN}{1}{0}{{0}}%
\reglabel{Reset}\regnewline%
\begin{regdesc}\begin{reglist}
\label{fielddesc:MCPWMCAP2EN}\item [MCPWM\_CAP2\_EN] 配置是否使能信道 2 上的捕获事件。\\
0：不使能\\
1：使能\\
(R/W)
\label{fielddesc:MCPWMCAP2MODE}\item [MCPWM\_CAP2\_MODE] 配置预分频后信道 2 上的捕获沿。\\
bit0 为 1：使能下降沿捕获。\\
bit1 为 1：使能上升沿捕获。\\
(R/W)
\label{fielddesc:MCPWMCAP2PRESCALE}\item [MCPWM\_CAP2\_PRESCALE] 配置 CAP2 上 升 沿 的 预 分 频 值。 该 预 分 频 值 = PWM\_CAP2\_PRESCALE + 1. (R/W)
\label{fielddesc:MCPWMCAP2ININVERT}\item [MCPWM\_CAP2\_IN\_INVERT] 配置是否在预分频前反相来自 GPIO 矩阵的 CAP2。\\
0：无效\\
1：在预分频前反相来自 GPIO 矩阵的 CAP2\\
(R/W)
\label{fielddesc:MCPWMCAP2SW}\item [MCPWM\_CAP2\_SW] 配置是否触发信道 2 上的软件强制捕获事件。\\
0：不触发\\
1：触发\\
(WT)
\end{reglist}\end{regdesc}
\end{register}


\begin{register}{H}{MCPWM\_CAP\_CH0\_REG}{0x00FC}\label{regdesc:MCPWMCAPCH0REG}
\regfield{MCPWM\_CAP0\_VALUE}{32}{0}{{0}}%
\reglabel{Reset}\regnewline%
\begin{regdesc}\begin{reglist}
\label{fielddesc:MCPWMCAP0VALUE}\item [MCPWM\_CAP0\_VALUE] 表示信道 0 上一次捕获的值。 (RO)
\end{reglist}\end{regdesc}
\end{register}


\begin{register}{H}{MCPWM\_CAP\_CH1\_REG}{0x0100}\label{regdesc:MCPWMCAPCH1REG}
\regfield{MCPWM\_CAP1\_VALUE}{32}{0}{{0}}%
\reglabel{Reset}\regnewline%
\begin{regdesc}\begin{reglist}
\label{fielddesc:MCPWMCAP1VALUE}\item [MCPWM\_CAP1\_VALUE] 表示信道 1 上一次捕获的值。 (RO)
\end{reglist}\end{regdesc}
\end{register}


\begin{register}{H}{MCPWM\_CAP\_CH2\_REG}{0x0104}\label{regdesc:MCPWMCAPCH2REG}
\regfield{MCPWM\_CAP2\_VALUE}{32}{0}{{0}}%
\reglabel{Reset}\regnewline%
\begin{regdesc}\begin{reglist}
\label{fielddesc:MCPWMCAP2VALUE}\item [MCPWM\_CAP2\_VALUE] 表示信道 2 上一次捕获的值。(RO)
\end{reglist}\end{regdesc}
\end{register}


\begin{register}{H}{MCPWM\_CAP\_STATUS\_REG}{0x0108}\label{regdesc:MCPWMCAPSTATUSREG}
\regfield{(reserved)}{29}{3}{00000000000000000000000000000}%
\regfield{MCPWM\_CAP2\_EDGE}{1}{2}{{0}}%
\regfield{MCPWM\_CAP1\_EDGE}{1}{1}{{0}}%
\regfield{MCPWM\_CAP0\_EDGE}{1}{0}{{0}}%
\reglabel{Reset}\regnewline%
\begin{regdesc}\begin{reglist}
\label{fielddesc:MCPWMCAP0EDGE}\item [MCPWM\_CAP0\_EDGE] 信道 0 上一次捕获触发事件的边沿。\\0：上升沿。\\1：下降沿。 \\(RO)
\label{fielddesc:MCPWMCAP1EDGE}\item [MCPWM\_CAP1\_EDGE] 信道 1 上一次捕获触发事件的边沿。\\0：上升沿。\\1：下降沿。 \\(RO)
\label{fielddesc:MCPWMCAP2EDGE}\item [MCPWM\_CAP2\_EDGE] 信道 2 上一次捕获触发事件的边沿。\\0：上升沿。\\1：下降沿。 \\(RO)
\end{reglist}\end{regdesc}
\end{register}


\begin{register}{H}{MCPWM\_UPDATE\_CFG\_REG}{0x010C}\label{regdesc:MCPWMUPDATECFGREG}
\regfield{(reserved)}{24}{8}{000000000000000000000000}%
\regfield{MCPWM\_OP2\_FORCE\_UP}{1}{7}{{0}}%
\regfield{MCPWM\_OP2\_UP\_EN}{1}{6}{{1}}%
\regfield{MCPWM\_OP1\_FORCE\_UP}{1}{5}{{0}}%
\regfield{MCPWM\_OP1\_UP\_EN}{1}{4}{{1}}%
\regfield{MCPWM\_OP0\_FORCE\_UP}{1}{3}{{0}}%
\regfield{MCPWM\_OP0\_UP\_EN}{1}{2}{{1}}%
\regfield{MCPWM\_GLOBAL\_FORCE\_UP}{1}{1}{{0}}%
\regfield{MCPWM\_GLOBAL\_UP\_EN}{1}{0}{{1}}%
\reglabel{Reset}\regnewline%
\begin{regdesc}\begin{reglist}
\label{fielddesc:MCPWMGLOBALUPEN}\item [MCPWM\_GLOBAL\_UP\_EN] 配置是否更新 MCPWM 模块的所有有效寄存器。 \\
0：无效
1：更新 MCPWM 模块的所有有效寄存器 \\
(R/W)
\label{fielddesc:MCPWMGLOBALFORCEUP}\item [MCPWM\_GLOBAL\_FORCE\_UP] 配置是否强制更新 MCPWM 模块的所有有效寄存器。\\
0：无效
1：强制更新 MCPWM 模块的所有有效寄存器 \\
(R/W)
\label{fielddesc:MCPWMOP0UPEN}\item [MCPWM\_OP0\_UP\_EN] 当 \hyperref[fielddesc:MCPWMGLOBALUPEN]{MCPWM\_GLOBAL\_UP\_EN} 置 1 时，配置是否更新 PWM 操作器 0 中的有效寄存器。 \\
0：无效
1：更新 PWM 操作器 0 中的有效寄存器 \\
(R/W)
\label{fielddesc:MCPWMOP0FORCEUP}\item [MCPWM\_OP0\_FORCE\_UP] 配置是否强制更新 PWM 操作器 0 中的有效寄存器。\\
0：无效
1：强制更新 \\
(R/W)
\label{fielddesc:MCPWMOP1UPEN}\item [MCPWM\_OP1\_UP\_EN] 当 \hyperref[fielddesc:MCPWMGLOBALUPEN]{MCPWM\_GLOBAL\_UP\_EN} 置 1 时，配置是否更新 PWM 操作器 1 中的有效寄存器。 \\
0：无效
1：更新 PWM 操作器 1 中的有效寄存器 \\
(R/W)
\label{fielddesc:MCPWMOP1FORCEUP}\item [MCPWM\_OP1\_FORCE\_UP] 配置是否强制更新 PWM 操作器 1 中的有效寄存器。\\
0：无效
1：强制更新 \\
(R/W)
\label{fielddesc:MCPWMOP2UPEN}\item [MCPWM\_OP2\_UP\_EN] 当 \hyperref[fielddesc:MCPWMGLOBALUPEN]{MCPWM\_GLOBAL\_UP\_EN} 置 1 时，配置是否更新 PWM 操作器 2 中的有效寄存器。 \\
0：无效
1：更新 PWM 操作器 2 中的有效寄存器 \\
(R/W)
\label{fielddesc:MCPWMOP2FORCEUP}\item [MCPWM\_OP2\_FORCE\_UP] 配置是否强制更新 PWM 操作器 2 中的有效寄存器。\\
0：无效
1：强制更新 \\
(R/W)
\end{reglist}\end{regdesc}
\end{register}


\begin{register}{H}{MCPWM\_INT\_ENA\_REG}{0x0110}\label{regdesc:MCPWMINTENAREG}
\regfield{(reserved)}{2}{30}{00}%
\regfield{MCPWM\_CAP2\_INT\_ENA}{1}{29}{{0}}%
\regfield{MCPWM\_CAP1\_INT\_ENA}{1}{28}{{0}}%
\regfield{MCPWM\_CAP0\_INT\_ENA}{1}{27}{{0}}%
\regfield{MCPWM\_TZ2\_OST\_INT\_ENA}{1}{26}{{0}}%
\regfield{MCPWM\_TZ1\_OST\_INT\_ENA}{1}{25}{{0}}%
\regfield{MCPWM\_TZ0\_OST\_INT\_ENA}{1}{24}{{0}}%
\regfield{MCPWM\_TZ2\_CBC\_INT\_ENA}{1}{23}{{0}}%
\regfield{MCPWM\_TZ1\_CBC\_INT\_ENA}{1}{22}{{0}}%
\regfield{MCPWM\_TZ0\_CBC\_INT\_ENA}{1}{21}{{0}}%
\regfield{MCPWM\_CMPR2\_TEB\_INT\_ENA}{1}{20}{{0}}%
\regfield{MCPWM\_CMPR1\_TEB\_INT\_ENA}{1}{19}{{0}}%
\regfield{MCPWM\_CMPR0\_TEB\_INT\_ENA}{1}{18}{{0}}%
\regfield{MCPWM\_CMPR2\_TEA\_INT\_ENA}{1}{17}{{0}}%
\regfield{MCPWM\_CMPR1\_TEA\_INT\_ENA}{1}{16}{{0}}%
\regfield{MCPWM\_CMPR0\_TEA\_INT\_ENA}{1}{15}{{0}}%
\regfield{MCPWM\_FAULT2\_CLR\_INT\_ENA}{1}{14}{{0}}%
\regfield{MCPWM\_FAULT1\_CLR\_INT\_ENA}{1}{13}{{0}}%
\regfield{MCPWM\_FAULT0\_CLR\_INT\_ENA}{1}{12}{{0}}%
\regfield{MCPWM\_FAULT2\_INT\_ENA}{1}{11}{{0}}%
\regfield{MCPWM\_FAULT1\_INT\_ENA}{1}{10}{{0}}%
\regfield{MCPWM\_FAULT0\_INT\_ENA}{1}{9}{{0}}%
\regfield{MCPWM\_TIMER2\_TEP\_INT\_ENA}{1}{8}{{0}}%
\regfield{MCPWM\_TIMER1\_TEP\_INT\_ENA}{1}{7}{{0}}%
\regfield{MCPWM\_TIMER0\_TEP\_INT\_ENA}{1}{6}{{0}}%
\regfield{MCPWM\_TIMER2\_TEZ\_INT\_ENA}{1}{5}{{0}}%
\regfield{MCPWM\_TIMER1\_TEZ\_INT\_ENA}{1}{4}{{0}}%
\regfield{MCPWM\_TIMER0\_TEZ\_INT\_ENA}{1}{3}{{0}}%
\regfield{MCPWM\_TIMER2\_STOP\_INT\_ENA}{1}{2}{{0}}%
\regfield{MCPWM\_TIMER1\_STOP\_INT\_ENA}{1}{1}{{0}}%
\regfield{MCPWM\_TIMER0\_STOP\_INT\_ENA}{1}{0}{{0}}%
\reglabel{Reset}\regnewline%
\begin{regdesc}\begin{reglist}
\label{fielddesc:MCPWMTIMER0STOPINTENA}\item [MCPWM\_TIMER0\_STOP\_INT\_ENA] 写 1 使能 \hyperref[sec:mcpwm-interrupts]{MCPWM\_TIMER0\_STOP\_INT}。 (R/W)
\label{fielddesc:MCPWMTIMER1STOPINTENA}\item [MCPWM\_TIMER1\_STOP\_INT\_ENA] 写 1 使能 \hyperref[sec:mcpwm-interrupts]{MCPWM\_TIMER1\_STOP\_INT}。 (R/W)
\label{fielddesc:MCPWMTIMER2STOPINTENA}\item [MCPWM\_TIMER2\_STOP\_INT\_ENA] 写 1 使能 \hyperref[sec:mcpwm-interrupts]{MCPWM\_TIMER2\_STOP\_INT}。 (R/W)
\label{fielddesc:MCPWMTIMER0TEZINTENA}\item [MCPWM\_TIMER0\_TEZ\_INT\_ENA] 写 1 使能 \hyperref[sec:mcpwm-interrupts]{MCPWM\_TIMER0\_TEZ\_INT}。 (R/W)
\label{fielddesc:MCPWMTIMER1TEZINTENA}\item [MCPWM\_TIMER1\_TEZ\_INT\_ENA] 写 1 使能 \hyperref[sec:mcpwm-interrupts]{MCPWM\_TIMER1\_TEZ\_INT}。 (R/W)
\label{fielddesc:MCPWMTIMER2TEZINTENA}\item [MCPWM\_TIMER2\_TEZ\_INT\_ENA] 写 1 使能 \hyperref[sec:mcpwm-interrupts]{MCPWM\_TIMER2\_TEZ\_INT}。 (R/W)
\label{fielddesc:MCPWMTIMER0TEPINTENA}\item [MCPWM\_TIMER0\_TEP\_INT\_ENA] 写 1 使能 \hyperref[sec:mcpwm-interrupts]{MCPWM\_TIMER0\_TEP\_INT}。 (R/W)
\label{fielddesc:MCPWMTIMER1TEPINTENA}\item [MCPWM\_TIMER1\_TEP\_INT\_ENA] 写 1 使能 \hyperref[sec:mcpwm-interrupts]{MCPWM\_TIMER1\_TEP\_INT}。 (R/W)
\label{fielddesc:MCPWMTIMER2TEPINTENA}\item [MCPWM\_TIMER2\_TEP\_INT\_ENA] 写 1 使能 \hyperref[sec:mcpwm-interrupts]{MCPWM\_TIMER2\_TEP\_INT}。 (R/W)
\label{fielddesc:MCPWMFAULT0INTENA}\item [MCPWM\_FAULT0\_INT\_ENA] 写 1 使能 \hyperref[sec:mcpwm-interrupts]{MCPWM\_FAULT0\_INT}。 (R/W)
\label{fielddesc:MCPWMFAULT1INTENA}\item [MCPWM\_FAULT1\_INT\_ENA] 写 1 使能 \hyperref[sec:mcpwm-interrupts]{MCPWM\_FAULT1\_INT}。 (R/W)
\label{fielddesc:MCPWMFAULT2INTENA}\item [MCPWM\_FAULT2\_INT\_ENA] 写 1 使能 \hyperref[sec:mcpwm-interrupts]{MCPWM\_FAULT2\_INT}。 (R/W)
\label{fielddesc:MCPWMFAULT0CLRINTENA}\item [MCPWM\_FAULT0\_CLR\_INT\_ENA] 写 1 使能 \hyperref[sec:mcpwm-interrupts]{MCPWM\_FAULT0\_CLR\_INT}。 (R/W)
\label{fielddesc:MCPWMFAULT1CLRINTENA}\item [MCPWM\_FAULT1\_CLR\_INT\_ENA] 写 1 使能 \hyperref[sec:mcpwm-interrupts]{MCPWM\_FAULT1\_CLR\_INT}。 (R/W)
\label{fielddesc:MCPWMFAULT2CLRINTENA}\item [MCPWM\_FAULT2\_CLR\_INT\_ENA] 写 1 使能 \hyperref[sec:mcpwm-interrupts]{MCPWM\_FAULT2\_CLR\_INT}。(R/W)
\label{fielddesc:MCPWMCMPR0TEAINTENA}\item [MCPWM\_CMPR0\_TEA\_INT\_ENA] 写 1 使能 \hyperref[sec:mcpwm-interrupts]{MCPWM\_CMPR0\_TEA\_INT}。 (R/W)
\label{fielddesc:MCPWMCMPR1TEAINTENA}\item [MCPWM\_CMPR1\_TEA\_INT\_ENA] 写 1 使能 \hyperref[sec:mcpwm-interrupts]{MCPWM\_CMPR1\_TEA\_INT}。 (R/W)
\label{fielddesc:MCPWMCMPR2TEAINTENA}\item [MCPWM\_CMPR2\_TEA\_INT\_ENA] 写 1 使能 \hyperref[sec:mcpwm-interrupts]{MCPWM\_CMPR2\_TEA\_INT}。 (R/W)
\item [见下页...]
\end{reglist}\end{regdesc}
\end{register}
\addtocounter{Regfloat}{-1}
\begin{register}{H}{MCPWM\_INT\_ENA\_REG}{0x0110}
\begin{regdesc}\begin{reglist}
\item [...接上页]
\label{fielddesc:MCPWMCMPR0TEBINTENA}\item [MCPWM\_CMPR0\_TEB\_INT\_ENA] 写 1 使能 \hyperref[sec:mcpwm-interrupts]{MCPWM\_CMPR0\_TEB\_INT}。 (R/W)
\label{fielddesc:MCPWMCMPR1TEBINTENA}\item [MCPWM\_CMPR1\_TEB\_INT\_ENA] 写 1 使能 \hyperref[sec:mcpwm-interrupts]{MCPWM\_CMPR1\_TEB\_INT}。 (R/W)
\label{fielddesc:MCPWMCMPR2TEBINTENA}\item [MCPWM\_CMPR2\_TEB\_INT\_ENA] 写 1 使能 \hyperref[sec:mcpwm-interrupts]{MCPWM\_CMPR2\_TEB\_INT}。 (R/W)
\label{fielddesc:MCPWMTZ0CBCINTENA}\item [MCPWM\_TZ0\_CBC\_INT\_ENA] 写 1 使能 \hyperref[sec:mcpwm-interrupts]{MCPWM\_TZ0\_CBC\_INT}。 (R/W)
\label{fielddesc:MCPWMTZ1CBCINTENA}\item [MCPWM\_TZ1\_CBC\_INT\_ENA] 写 1 使能 \hyperref[sec:mcpwm-interrupts]{MCPWM\_TZ1\_CBC\_INT}。 (R/W)
\label{fielddesc:MCPWMTZ2CBCINTENA}\item [MCPWM\_TZ2\_CBC\_INT\_ENA] 写 1 使能 \hyperref[sec:mcpwm-interrupts]{MCPWM\_TZ2\_CBC\_INT}。 (R/W)
\label{fielddesc:MCPWMTZ0OSTINTENA}\item [MCPWM\_TZ0\_OST\_INT\_ENA] 写 1 使能 \hyperref[sec:mcpwm-interrupts]{MCPWM\_TZ0\_OST\_INT}。 (R/W)
\label{fielddesc:MCPWMTZ1OSTINTENA}\item [MCPWM\_TZ1\_OST\_INT\_ENA] 写 1 使能 \hyperref[sec:mcpwm-interrupts]{MCPWM\_TZ1\_OST\_INT}。 (R/W)
\label{fielddesc:MCPWMTZ2OSTINTENA}\item [MCPWM\_TZ2\_OST\_INT\_ENA] 写 1 使能 \hyperref[sec:mcpwm-interrupts]{MCPWM\_TZ2\_OST\_INT}。 (R/W)
\label{fielddesc:MCPWMCAP0INTENA}\item [MCPWM\_CAP0\_INT\_ENA] 写 1 使能 \hyperref[sec:mcpwm-interrupts]{MCPWM\_CAP0\_INT}。 (R/W)
\label{fielddesc:MCPWMCAP1INTENA}\item [MCPWM\_CAP1\_INT\_ENA] 写 1 使能 \hyperref[sec:mcpwm-interrupts]{MCPWM\_CAP1\_INT}。 (R/W)
\label{fielddesc:MCPWMCAP2INTENA}\item [MCPWM\_CAP2\_INT\_ENA] 写 1 使能 \hyperref[sec:mcpwm-interrupts]{MCPWM\_CAP2\_INT}。 (R/W)
\end{reglist}\end{regdesc}
\end{register}


\begin{register}{H}{MCPWM\_INT\_RAW\_REG}{0x0114}\label{regdesc:MCPWMINTRAWREG}
\regfield{(reserved)}{2}{30}{00}%
\regfield{MCPWM\_CAP2\_INT\_RAW}{1}{29}{{0}}%
\regfield{MCPWM\_CAP1\_INT\_RAW}{1}{28}{{0}}%
\regfield{MCPWM\_CAP0\_INT\_RAW}{1}{27}{{0}}%
\regfield{MCPWM\_TZ2\_OST\_INT\_RAW}{1}{26}{{0}}%
\regfield{MCPWM\_TZ1\_OST\_INT\_RAW}{1}{25}{{0}}%
\regfield{MCPWM\_TZ0\_OST\_INT\_RAW}{1}{24}{{0}}%
\regfield{MCPWM\_TZ2\_CBC\_INT\_RAW}{1}{23}{{0}}%
\regfield{MCPWM\_TZ1\_CBC\_INT\_RAW}{1}{22}{{0}}%
\regfield{MCPWM\_TZ0\_CBC\_INT\_RAW}{1}{21}{{0}}%
\regfield{MCPWM\_CMPR2\_TEB\_INT\_RAW}{1}{20}{{0}}%
\regfield{MCPWM\_CMPR1\_TEB\_INT\_RAW}{1}{19}{{0}}%
\regfield{MCPWM\_CMPR0\_TEB\_INT\_RAW}{1}{18}{{0}}%
\regfield{MCPWM\_CMPR2\_TEA\_INT\_RAW}{1}{17}{{0}}%
\regfield{MCPWM\_CMPR1\_TEA\_INT\_RAW}{1}{16}{{0}}%
\regfield{MCPWM\_CMPR0\_TEA\_INT\_RAW}{1}{15}{{0}}%
\regfield{MCPWM\_FAULT2\_CLR\_INT\_RAW}{1}{14}{{0}}%
\regfield{MCPWM\_FAULT1\_CLR\_INT\_RAW}{1}{13}{{0}}%
\regfield{MCPWM\_FAULT0\_CLR\_INT\_RAW}{1}{12}{{0}}%
\regfield{MCPWM\_FAULT2\_INT\_RAW}{1}{11}{{0}}%
\regfield{MCPWM\_FAULT1\_INT\_RAW}{1}{10}{{0}}%
\regfield{MCPWM\_FAULT0\_INT\_RAW}{1}{9}{{0}}%
\regfield{MCPWM\_TIMER2\_TEP\_INT\_RAW}{1}{8}{{0}}%
\regfield{MCPWM\_TIMER1\_TEP\_INT\_RAW}{1}{7}{{0}}%
\regfield{MCPWM\_TIMER0\_TEP\_INT\_RAW}{1}{6}{{0}}%
\regfield{MCPWM\_TIMER2\_TEZ\_INT\_RAW}{1}{5}{{0}}%
\regfield{MCPWM\_TIMER1\_TEZ\_INT\_RAW}{1}{4}{{0}}%
\regfield{MCPWM\_TIMER0\_TEZ\_INT\_RAW}{1}{3}{{0}}%
\regfield{MCPWM\_TIMER2\_STOP\_INT\_RAW}{1}{2}{{0}}%
\regfield{MCPWM\_TIMER1\_STOP\_INT\_RAW}{1}{1}{{0}}%
\regfield{MCPWM\_TIMER0\_STOP\_INT\_RAW}{1}{0}{{0}}%
\reglabel{Reset}\regnewline%
\begin{regdesc}\begin{reglist}
\label{fielddesc:MCPWMTIMER0STOPINTRAW}\item [MCPWM\_TIMER0\_STOP\_INT\_RAW] \hyperref[sec:mcpwm-interrupts]{MCPWM\_TIMER0\_STOP\_INT} 的原始状态位。 (R/WTC/SS)
\label{fielddesc:MCPWMTIMER1STOPINTRAW}\item [MCPWM\_TIMER1\_STOP\_INT\_RAW] \hyperref[sec:mcpwm-interrupts]{MCPWM\_TIMER1\_STOP\_INT} 的原始状态位。 (R/WTC/SS)
\label{fielddesc:MCPWMTIMER2STOPINTRAW}\item [MCPWM\_TIMER2\_STOP\_INT\_RAW] \hyperref[sec:mcpwm-interrupts]{MCPWM\_TIMER2\_STOP\_INT} 的原始状态位。 (R/WTC/SS)
\label{fielddesc:MCPWMTIMER0TEZINTRAW}\item [MCPWM\_TIMER0\_TEZ\_INT\_RAW] \hyperref[sec:mcpwm-interrupts]{MCPWM\_TIMER0\_TEZ\_INT} 的原始状态位。 (R/WTC/SS)
\label{fielddesc:MCPWMTIMER1TEZINTRAW}\item [MCPWM\_TIMER1\_TEZ\_INT\_RAW] \hyperref[sec:mcpwm-interrupts]{MCPWM\_TIMER1\_TEZ\_INT} 的原始状态位。 (R/WTC/SS)
\label{fielddesc:MCPWMTIMER2TEZINTRAW}\item [MCPWM\_TIMER2\_TEZ\_INT\_RAW] \hyperref[sec:mcpwm-interrupts]{MCPWM\_TIMER2\_TEZ\_INT} 的原始状态位。 (R/WTC/SS)
\label{fielddesc:MCPWMTIMER0TEPINTRAW}\item [MCPWM\_TIMER0\_TEP\_INT\_RAW] \hyperref[sec:mcpwm-interrupts]{MCPWM\_TIMER0\_TEP\_INT} 的原始状态位。 (R/WTC/SS)
\label{fielddesc:MCPWMTIMER1TEPINTRAW}\item [MCPWM\_TIMER1\_TEP\_INT\_RAW] \hyperref[sec:mcpwm-interrupts]{MCPWM\_TIMER1\_TEP\_INT} 的原始状态位。 (R/WTC/SS)
\label{fielddesc:MCPWMTIMER2TEPINTRAW}\item [MCPWM\_TIMER2\_TEP\_INT\_RAW] \hyperref[sec:mcpwm-interrupts]{MCPWM\_TIMER2\_TEP\_INT} 的原始状态位。 (R/WTC/SS)
\label{fielddesc:MCPWMFAULT0INTRAW}\item [MCPWM\_FAULT0\_INT\_RAW] \hyperref[sec:mcpwm-interrupts]{MCPWM\_FAULT0\_INT} 的原始状态位。 (R/WTC/SS)
\label{fielddesc:MCPWMFAULT1INTRAW}\item [MCPWM\_FAULT1\_INT\_RAW] \hyperref[sec:mcpwm-interrupts]{MCPWM\_FAULT1\_INT} 的原始状态位。 (R/WTC/SS)
\label{fielddesc:MCPWMFAULT2INTRAW}\item [MCPWM\_FAULT2\_INT\_RAW] \hyperref[sec:mcpwm-interrupts]{MCPWM\_FAULT2\_INT} 的原始状态位。 (R/WTC/SS)
\label{fielddesc:MCPWMFAULT0CLRINTRAW}\item [MCPWM\_FAULT0\_CLR\_INT\_RAW] \hyperref[sec:mcpwm-interrupts]{MCPWM\_FAULT0\_CLR\_INT} 的原始状态位。 (R/WTC/SS)
\label{fielddesc:MCPWMFAULT1CLRINTRAW}\item [MCPWM\_FAULT1\_CLR\_INT\_RAW] \hyperref[sec:mcpwm-interrupts]{MCPWM\_FAULT1\_CLR\_INT} 的原始状态位。 (R/WTC/SS)
\label{fielddesc:MCPWMFAULT2CLRINTRAW}\item [MCPWM\_FAULT2\_CLR\_INT\_RAW] \hyperref[sec:mcpwm-interrupts]{MCPWM\_FAULT2\_CLR\_INT} 的原始状态位。 (R/WTC/SS)
\item [见下页...]
\end{reglist}\end{regdesc}
\end{register}
\addtocounter{Regfloat}{-1}
\begin{register}{H}{MCPWM\_INT\_RAW\_REG}{0x0114}
\begin{regdesc}\begin{reglist}
\item [...接上页]
\label{fielddesc:MCPWMCMPR0TEAINTRAW}\item [MCPWM\_CMPR0\_TEA\_INT\_RAW] \hyperref[sec:mcpwm-interrupts]{MCPWM\_CMPR0\_TEA\_INT} 的原始状态位。 (R/WTC/SS)
\label{fielddesc:MCPWMCMPR1TEAINTRAW}\item [MCPWM\_CMPR1\_TEA\_INT\_RAW] \hyperref[sec:mcpwm-interrupts]{MCPWM\_CMPR1\_TEA\_INT} 的原始状态位。 (R/WTC/SS)
\label{fielddesc:MCPWMCMPR2TEAINTRAW}\item [MCPWM\_CMPR2\_TEA\_INT\_RAW] \hyperref[sec:mcpwm-interrupts]{MCPWM\_CMPR2\_TEA\_INT} 的原始状态位。 (R/WTC/SS)
\label{fielddesc:MCPWMCMPR0TEBINTRAW}\item [MCPWM\_CMPR0\_TEB\_INT\_RAW] \hyperref[sec:mcpwm-interrupts]{MCPWM\_CMPR0\_TEB\_INT} 的原始状态位。 (R/WTC/SS)
\label{fielddesc:MCPWMCMPR1TEBINTRAW}\item [MCPWM\_CMPR1\_TEB\_INT\_RAW] \hyperref[sec:mcpwm-interrupts]{MCPWM\_CMPR1\_TEB\_INT} 的原始状态位。 (R/WTC/SS)
\label{fielddesc:MCPWMCMPR2TEBINTRAW}\item [MCPWM\_CMPR2\_TEB\_INT\_RAW] \hyperref[sec:mcpwm-interrupts]{MCPWM\_CMPR2\_TEB\_INT} 的原始状态位。 (R/WTC/SS)
\label{fielddesc:MCPWMTZ0CBCINTRAW}\item [MCPWM\_TZ0\_CBC\_INT\_RAW] \hyperref[sec:mcpwm-interrupts]{MCPWM\_TZ0\_CBC\_INT} 的原始状态位。 (R/WTC/SS)
\label{fielddesc:MCPWMTZ1CBCINTRAW}\item [MCPWM\_TZ1\_CBC\_INT\_RAW] \hyperref[sec:mcpwm-interrupts]{MCPWM\_TZ1\_CBC\_INT} 的原始状态位。 (R/WTC/SS)
\label{fielddesc:MCPWMTZ2CBCINTRAW}\item [MCPWM\_TZ2\_CBC\_INT\_RAW] \hyperref[sec:mcpwm-interrupts]{MCPWM\_TZ2\_CBC\_INT} 的原始状态位。 (R/WTC/SS)
\label{fielddesc:MCPWMTZ0OSTINTRAW}\item [MCPWM\_TZ0\_OST\_INT\_RAW] \hyperref[sec:mcpwm-interrupts]{MCPWM\_TZ0\_OST\_INT} 的原始状态位。 (R/WTC/SS)
\label{fielddesc:MCPWMTZ1OSTINTRAW}\item [MCPWM\_TZ1\_OST\_INT\_RAW] \hyperref[sec:mcpwm-interrupts]{MCPWM\_TZ1\_OST\_INT} 的原始状态位。 (R/WTC/SS)
\label{fielddesc:MCPWMTZ2OSTINTRAW}\item [MCPWM\_TZ2\_OST\_INT\_RAW] \hyperref[sec:mcpwm-interrupts]{MCPWM\_TZ2\_OST\_INT} 的原始状态位。 (R/WTC/SS)
\label{fielddesc:MCPWMCAP0INTRAW}\item [MCPWM\_CAP0\_INT\_RAW] \hyperref[sec:mcpwm-interrupts]{MCPWM\_CAP0\_INT} 的原始状态位。 (R/WTC/SS)
\label{fielddesc:MCPWMCAP1INTRAW}\item [MCPWM\_CAP1\_INT\_RAW] \hyperref[sec:mcpwm-interrupts]{MCPWM\_CAP1\_INT} 的原始状态位。 (R/WTC/SS)
\label{fielddesc:MCPWMCAP2INTRAW}\item [MCPWM\_CAP2\_INT\_RAW] \hyperref[sec:mcpwm-interrupts]{MCPWM\_CAP2\_INT} 的原始状态位。 (R/WTC/SS)
\end{reglist}\end{regdesc}
\end{register}


\begin{register}{H}{MCPWM\_INT\_ST\_REG}{0x0118}\label{regdesc:MCPWMINTSTREG}
\regfield{(reserved)}{2}{30}{00}%
\regfield{MCPWM\_CAP2\_INT\_ST}{1}{29}{{0}}%
\regfield{MCPWM\_CAP1\_INT\_ST}{1}{28}{{0}}%
\regfield{MCPWM\_CAP0\_INT\_ST}{1}{27}{{0}}%
\regfield{MCPWM\_TZ2\_OST\_INT\_ST}{1}{26}{{0}}%
\regfield{MCPWM\_TZ1\_OST\_INT\_ST}{1}{25}{{0}}%
\regfield{MCPWM\_TZ0\_OST\_INT\_ST}{1}{24}{{0}}%
\regfield{MCPWM\_TZ2\_CBC\_INT\_ST}{1}{23}{{0}}%
\regfield{MCPWM\_TZ1\_CBC\_INT\_ST}{1}{22}{{0}}%
\regfield{MCPWM\_TZ0\_CBC\_INT\_ST}{1}{21}{{0}}%
\regfield{MCPWM\_CMPR2\_TEB\_INT\_ST}{1}{20}{{0}}%
\regfield{MCPWM\_CMPR1\_TEB\_INT\_ST}{1}{19}{{0}}%
\regfield{MCPWM\_CMPR0\_TEB\_INT\_ST}{1}{18}{{0}}%
\regfield{MCPWM\_CMPR2\_TEA\_INT\_ST}{1}{17}{{0}}%
\regfield{MCPWM\_CMPR1\_TEA\_INT\_ST}{1}{16}{{0}}%
\regfield{MCPWM\_CMPR0\_TEA\_INT\_ST}{1}{15}{{0}}%
\regfield{MCPWM\_FAULT2\_CLR\_INT\_ST}{1}{14}{{0}}%
\regfield{MCPWM\_FAULT1\_CLR\_INT\_ST}{1}{13}{{0}}%
\regfield{MCPWM\_FAULT0\_CLR\_INT\_ST}{1}{12}{{0}}%
\regfield{MCPWM\_FAULT2\_INT\_ST}{1}{11}{{0}}%
\regfield{MCPWM\_FAULT1\_INT\_ST}{1}{10}{{0}}%
\regfield{MCPWM\_FAULT0\_INT\_ST}{1}{9}{{0}}%
\regfield{MCPWM\_TIMER2\_TEP\_INT\_ST}{1}{8}{{0}}%
\regfield{MCPWM\_TIMER1\_TEP\_INT\_ST}{1}{7}{{0}}%
\regfield{MCPWM\_TIMER0\_TEP\_INT\_ST}{1}{6}{{0}}%
\regfield{MCPWM\_TIMER2\_TEZ\_INT\_ST}{1}{5}{{0}}%
\regfield{MCPWM\_TIMER1\_TEZ\_INT\_ST}{1}{4}{{0}}%
\regfield{MCPWM\_TIMER0\_TEZ\_INT\_ST}{1}{3}{{0}}%
\regfield{MCPWM\_TIMER2\_STOP\_INT\_ST}{1}{2}{{0}}%
\regfield{MCPWM\_TIMER1\_STOP\_INT\_ST}{1}{1}{{0}}%
\regfield{MCPWM\_TIMER0\_STOP\_INT\_ST}{1}{0}{{0}}%
\reglabel{Reset}\regnewline%
\begin{regdesc}\begin{reglist}
\label{fielddesc:MCPWMTIMER0STOPINTST}\item [MCPWM\_TIMER0\_STOP\_INT\_ST] \hyperref[sec:mcpwm-interrupts]{MCPWM\_TIMER0\_STOP\_INT} 的屏蔽状态位。 (RO)
\label{fielddesc:MCPWMTIMER1STOPINTST}\item [MCPWM\_TIMER1\_STOP\_INT\_ST] \hyperref[sec:mcpwm-interrupts]{MCPWM\_TIMER1\_STOP\_INT} 的屏蔽状态位。 (RO)
\label{fielddesc:MCPWMTIMER2STOPINTST}\item [MCPWM\_TIMER2\_STOP\_INT\_ST] \hyperref[sec:mcpwm-interrupts]{MCPWM\_TIMER2\_STOP\_INT} 的屏蔽状态位。 (RO)
\label{fielddesc:MCPWMTIMER0TEZINTST}\item [MCPWM\_TIMER0\_TEZ\_INT\_ST] \hyperref[sec:mcpwm-interrupts]{MCPWM\_TIMER0\_TEZ\_INT} 的屏蔽状态位。 (RO)
\label{fielddesc:MCPWMTIMER1TEZINTST}\item [MCPWM\_TIMER1\_TEZ\_INT\_ST] \hyperref[sec:mcpwm-interrupts]{MCPWM\_TIMER1\_TEZ\_INT} 的屏蔽状态位。 (RO)
\label{fielddesc:MCPWMTIMER2TEZINTST}\item [MCPWM\_TIMER2\_TEZ\_INT\_ST] \hyperref[sec:mcpwm-interrupts]{MCPWM\_TIMER2\_TEZ\_INT} 的屏蔽状态位。 (RO)
\label{fielddesc:MCPWMTIMER0TEPINTST}\item [MCPWM\_TIMER0\_TEP\_INT\_ST] \hyperref[sec:mcpwm-interrupts]{MCPWM\_TIMER0\_TEP\_INT} 的屏蔽状态位。 (RO)
\label{fielddesc:MCPWMTIMER1TEPINTST}\item [MCPWM\_TIMER1\_TEP\_INT\_ST] \hyperref[sec:mcpwm-interrupts]{MCPWM\_TIMER1\_TEP\_INT} 的屏蔽状态位。 (RO)
\label{fielddesc:MCPWMTIMER2TEPINTST}\item [MCPWM\_TIMER2\_TEP\_INT\_ST] \hyperref[sec:mcpwm-interrupts]{MCPWM\_TIMER2\_TEP\_INT} 的屏蔽状态位。 (RO)
\label{fielddesc:MCPWMFAULT0INTST}\item [MCPWM\_FAULT0\_INT\_ST] \hyperref[sec:mcpwm-interrupts]{MCPWM\_FAULT0\_INT\_INT} 的屏蔽状态位。 (RO)
\label{fielddesc:MCPWMFAULT1INTST}\item [MCPWM\_FAULT1\_INT\_ST] \hyperref[sec:mcpwm-interrupts]{MCPWM\_FAULT1\_INT\_INT} 的屏蔽状态位。 (RO)
\label{fielddesc:MCPWMFAULT2INTST}\item [MCPWM\_FAULT2\_INT\_ST] \hyperref[sec:mcpwm-interrupts]{MCPWM\_FAULT2\_INT\_INT} 的屏蔽状态位。 (RO)
\label{fielddesc:MCPWMFAULT0CLRINTST}\item [MCPWM\_FAULT0\_CLR\_INT\_ST] \hyperref[sec:mcpwm-interrupts]{MCPWM\_FAULT0\_CLR\_INT} 的屏蔽状态位。 (RO)
\label{fielddesc:MCPWMFAULT1CLRINTST}\item [MCPWM\_FAULT1\_CLR\_INT\_ST] \hyperref[sec:mcpwm-interrupts]{MCPWM\_FAULT1\_CLR\_INT} 的屏蔽状态位。 (RO)
\label{fielddesc:MCPWMFAULT2CLRINTST}\item [MCPWM\_FAULT2\_CLR\_INT\_ST] \hyperref[sec:mcpwm-interrupts]{MCPWM\_FAULT2\_CLR\_INT} 的屏蔽状态位。 (RO)
\label{fielddesc:MCPWMCMPR0TEAINTST}\item [MCPWM\_CMPR0\_TEA\_INT\_ST] \hyperref[sec:mcpwm-interrupts]{MCPWM\_CMPR0\_TEA\_INT} 的屏蔽状态位。 (RO)
\label{fielddesc:MCPWMCMPR1TEAINTST}\item [MCPWM\_CMPR1\_TEA\_INT\_ST] \hyperref[sec:mcpwm-interrupts]{MCPWM\_CMPR1\_TEA\_INT} 的屏蔽状态位。 (RO)
\label{fielddesc:MCPWMCMPR2TEAINTST}\item [MCPWM\_CMPR2\_TEA\_INT\_ST] \hyperref[sec:mcpwm-interrupts]{MCPWM\_CMPR2\_TEA\_INT} 的屏蔽状态位。 (RO)
\item [见下页...]
\end{reglist}\end{regdesc}
\end{register}
\addtocounter{Regfloat}{-1}
\begin{register}{H}{MCPWM\_INT\_ST\_REG}{0x0118}
\begin{regdesc}\begin{reglist}
\item [...接上页]
\label{fielddesc:MCPWMCMPR0TEBINTST}\item [MCPWM\_CMPR0\_TEB\_INT\_ST] \hyperref[sec:mcpwm-interrupts]{MCPWM\_CMPR0\_TEB\_INT} 的屏蔽状态位。 (RO)
\label{fielddesc:MCPWMCMPR1TEBINTST}\item [MCPWM\_CMPR1\_TEB\_INT\_ST] \hyperref[sec:mcpwm-interrupts]{MCPWM\_CMPR1\_TEB\_INT} 的屏蔽状态位。 (RO)
\label{fielddesc:MCPWMCMPR2TEBINTST}\item [MCPWM\_CMPR2\_TEB\_INT\_ST] \hyperref[sec:mcpwm-interrupts]{MCPWM\_CMPR2\_TEB\_INT} 的屏蔽状态位。 (RO)
\label{fielddesc:MCPWMTZ0CBCINTST}\item [MCPWM\_TZ0\_CBC\_INT\_ST] \hyperref[sec:mcpwm-interrupts]{MCPWM\_TZ0\_CBC\_INT} 的屏蔽状态位。 (RO)
\label{fielddesc:MCPWMTZ1CBCINTST}\item [MCPWM\_TZ1\_CBC\_INT\_ST] \hyperref[sec:mcpwm-interrupts]{MCPWM\_TZ1\_CBC\_INT} 的屏蔽状态位。 (RO)
\label{fielddesc:MCPWMTZ2CBCINTST}\item [MCPWM\_TZ2\_CBC\_INT\_ST] \hyperref[sec:mcpwm-interrupts]{MCPWM\_TZ2\_CBC\_INT} 的屏蔽状态位。 (RO)
\label{fielddesc:MCPWMTZ0OSTINTST}\item [MCPWM\_TZ0\_OST\_INT\_ST] \hyperref[sec:mcpwm-interrupts]{MCPWM\_TZ0\_OST\_INT} 的屏蔽状态位。 (RO)
\label{fielddesc:MCPWMTZ1OSTINTST}\item [MCPWM\_TZ1\_OST\_INT\_ST] \hyperref[sec:mcpwm-interrupts]{MCPWM\_TZ1\_OST\_INT} 的屏蔽状态位。 (RO)
\label{fielddesc:MCPWMTZ2OSTINTST}\item [MCPWM\_TZ2\_OST\_INT\_ST] \hyperref[sec:mcpwm-interrupts]{MCPWM\_TZ2\_OST\_INT} 的屏蔽状态位。 (RO)
\label{fielddesc:MCPWMCAP0INTST}\item [MCPWM\_CAP0\_INT\_ST] \hyperref[sec:mcpwm-interrupts]{MCPWM\_CAP0\_INT} 的屏蔽状态位。 (RO)
\label{fielddesc:MCPWMCAP1INTST}\item [MCPWM\_CAP1\_INT\_ST] \hyperref[sec:mcpwm-interrupts]{MCPWM\_CAP1\_INT} 的屏蔽状态位。 (RO)
\label{fielddesc:MCPWMCAP2INTST}\item [MCPWM\_CAP2\_INT\_ST] \hyperref[sec:mcpwm-interrupts]{MCPWM\_CAP2\_INT} 的屏蔽状态位。 (RO)
\end{reglist}\end{regdesc}
\end{register}


\begin{register}{H}{MCPWM\_INT\_CLR\_REG}{0x011C}\label{regdesc:MCPWMINTCLRREG}
\regfield{(reserved)}{2}{30}{00}%
\regfield{MCPWM\_CAP2\_INT\_CLR}{1}{29}{{0}}%
\regfield{MCPWM\_CAP1\_INT\_CLR}{1}{28}{{0}}%
\regfield{MCPWM\_CAP0\_INT\_CLR}{1}{27}{{0}}%
\regfield{MCPWM\_TZ2\_OST\_INT\_CLR}{1}{26}{{0}}%
\regfield{MCPWM\_TZ1\_OST\_INT\_CLR}{1}{25}{{0}}%
\regfield{MCPWM\_TZ0\_OST\_INT\_CLR}{1}{24}{{0}}%
\regfield{MCPWM\_TZ2\_CBC\_INT\_CLR}{1}{23}{{0}}%
\regfield{MCPWM\_TZ1\_CBC\_INT\_CLR}{1}{22}{{0}}%
\regfield{MCPWM\_TZ0\_CBC\_INT\_CLR}{1}{21}{{0}}%
\regfield{MCPWM\_CMPR2\_TEB\_INT\_CLR}{1}{20}{{0}}%
\regfield{MCPWM\_CMPR1\_TEB\_INT\_CLR}{1}{19}{{0}}%
\regfield{MCPWM\_CMPR0\_TEB\_INT\_CLR}{1}{18}{{0}}%
\regfield{MCPWM\_CMPR2\_TEA\_INT\_CLR}{1}{17}{{0}}%
\regfield{MCPWM\_CMPR1\_TEA\_INT\_CLR}{1}{16}{{0}}%
\regfield{MCPWM\_CMPR0\_TEA\_INT\_CLR}{1}{15}{{0}}%
\regfield{MCPWM\_FAULT2\_CLR\_INT\_CLR}{1}{14}{{0}}%
\regfield{MCPWM\_FAULT1\_CLR\_INT\_CLR}{1}{13}{{0}}%
\regfield{MCPWM\_FAULT0\_CLR\_INT\_CLR}{1}{12}{{0}}%
\regfield{MCPWM\_FAULT2\_INT\_CLR}{1}{11}{{0}}%
\regfield{MCPWM\_FAULT1\_INT\_CLR}{1}{10}{{0}}%
\regfield{MCPWM\_FAULT0\_INT\_CLR}{1}{9}{{0}}%
\regfield{MCPWM\_TIMER2\_TEP\_INT\_CLR}{1}{8}{{0}}%
\regfield{MCPWM\_TIMER1\_TEP\_INT\_CLR}{1}{7}{{0}}%
\regfield{MCPWM\_TIMER0\_TEP\_INT\_CLR}{1}{6}{{0}}%
\regfield{MCPWM\_TIMER2\_TEZ\_INT\_CLR}{1}{5}{{0}}%
\regfield{MCPWM\_TIMER1\_TEZ\_INT\_CLR}{1}{4}{{0}}%
\regfield{MCPWM\_TIMER0\_TEZ\_INT\_CLR}{1}{3}{{0}}%
\regfield{MCPWM\_TIMER2\_STOP\_INT\_CLR}{1}{2}{{0}}%
\regfield{MCPWM\_TIMER1\_STOP\_INT\_CLR}{1}{1}{{0}}%
\regfield{MCPWM\_TIMER0\_STOP\_INT\_CLR}{1}{0}{{0}}%
\reglabel{Reset}\regnewline%
\begin{regdesc}\begin{reglist}
\label{fielddesc:MCPWMTIMER0STOPINTCLR}\item [MCPWM\_TIMER0\_STOP\_INT\_CLR] 写 1 清除 \hyperref[sec:mcpwm-interrupts]{MCPWM\_TIMER0\_STOP\_INT}。 (WT)
\label{fielddesc:MCPWMTIMER1STOPINTCLR}\item [MCPWM\_TIMER1\_STOP\_INT\_CLR] 写 1 清除 \hyperref[sec:mcpwm-interrupts]{MCPWM\_TIMER1\_STOP\_INT}。 (WT)
\label{fielddesc:MCPWMTIMER2STOPINTCLR}\item [MCPWM\_TIMER2\_STOP\_INT\_CLR] 写 1 清除 \hyperref[sec:mcpwm-interrupts]{MCPWM\_TIMER2\_STOP\_INT}。 (WT)
\label{fielddesc:MCPWMTIMER0TEZINTCLR}\item [MCPWM\_TIMER0\_TEZ\_INT\_CLR] 写 1 清除 \hyperref[sec:mcpwm-interrupts]{MCPWM\_TIMER0\_TEZ\_INT}。 (WT)
\label{fielddesc:MCPWMTIMER1TEZINTCLR}\item [MCPWM\_TIMER1\_TEZ\_INT\_CLR] 写 1 清除 \hyperref[sec:mcpwm-interrupts]{MCPWM\_TIMER1\_TEZ\_INT}。 (WT)
\label{fielddesc:MCPWMTIMER2TEZINTCLR}\item [MCPWM\_TIMER2\_TEZ\_INT\_CLR] 写 1 清除 \hyperref[sec:mcpwm-interrupts]{MCPWM\_TIMER2\_TEZ\_INT}。 (WT)
\label{fielddesc:MCPWMTIMER0TEPINTCLR}\item [MCPWM\_TIMER0\_TEP\_INT\_CLR] 写 1 清除 \hyperref[sec:mcpwm-interrupts]{MCPWM\_TIMER0\_TEP\_INT}。 (WT)
\label{fielddesc:MCPWMTIMER1TEPINTCLR}\item [MCPWM\_TIMER1\_TEP\_INT\_CLR] 写 1 清除 \hyperref[sec:mcpwm-interrupts]{MCPWM\_TIMER1\_TEP\_INT}。 (WT)
\label{fielddesc:MCPWMTIMER2TEPINTCLR}\item [MCPWM\_TIMER2\_TEP\_INT\_CLR] 写 1 清除 \hyperref[sec:mcpwm-interrupts]{MCPWM\_TIMER2\_TEP\_INT}。 (WT)
\label{fielddesc:MCPWMFAULT0INTCLR}\item [MCPWM\_FAULT0\_INT\_CLR] 写 1 清除 \hyperref[sec:mcpwm-interrupts]{MCPWM\_FAULT0\_INT}。 (WT)
\label{fielddesc:MCPWMFAULT1INTCLR}\item [MCPWM\_FAULT1\_INT\_CLR] 写 1 清除 \hyperref[sec:mcpwm-interrupts]{MCPWM\_FAULT1\_INT}。 (WT)
\label{fielddesc:MCPWMFAULT2INTCLR}\item [MCPWM\_FAULT2\_INT\_CLR] 写 1 清除 \hyperref[sec:mcpwm-interrupts]{MCPWM\_FAULT2\_INT}。 (WT)
\label{fielddesc:MCPWMFAULT0CLRINTCLR}\item [MCPWM\_FAULT0\_CLR\_INT\_CLR] 写 1 清除 \hyperref[sec:mcpwm-interrupts]{MCPWM\_FAULT0\_CLR\_INT}。 (WT)
\label{fielddesc:MCPWMFAULT1CLRINTCLR}\item [MCPWM\_FAULT1\_CLR\_INT\_CLR] 写 1 清除 \hyperref[sec:mcpwm-interrupts]{MCPWM\_FAULT1\_CLR\_INT}。 (WT)
\label{fielddesc:MCPWMFAULT2CLRINTCLR}\item [MCPWM\_FAULT2\_CLR\_INT\_CLR] 写 1 清除 \hyperref[sec:mcpwm-interrupts]{MCPWM\_FAULT2\_CLR\_INT}。 (WT)
\item [见下页...]
\end{reglist}\end{regdesc}
\end{register}
\addtocounter{Regfloat}{-1}
\begin{register}{H}{MCPWM\_INT\_CLR\_REG}{0x011C}
\begin{regdesc}\begin{reglist}
\item [...接上页]
\label{fielddesc:MCPWMCMPR0TEAINTCLR}\item [MCPWM\_CMPR0\_TEA\_INT\_CLR] 写 1 清除 \hyperref[sec:mcpwm-interrupts]{MCPWM\_CMPR0\_TEA\_INT}。 (WT)
\label{fielddesc:MCPWMCMPR1TEAINTCLR}\item [MCPWM\_CMPR1\_TEA\_INT\_CLR] 写 1 清除 \hyperref[sec:mcpwm-interrupts]{MCPWM\_CMPR1\_TEA\_INT}。 (WT)
\label{fielddesc:MCPWMCMPR2TEAINTCLR}\item [MCPWM\_CMPR2\_TEA\_INT\_CLR] 写 1 清除 \hyperref[sec:mcpwm-interrupts]{MCPWM\_CMPR2\_TEA\_INT}。 (WT)
\label{fielddesc:MCPWMCMPR0TEBINTCLR}\item [MCPWM\_CMPR0\_TEB\_INT\_CLR] 写 1 清除 \hyperref[sec:mcpwm-interrupts]{MCPWM\_CMPR0\_TEB\_INT}。 (WT)
\label{fielddesc:MCPWMCMPR1TEBINTCLR}\item [MCPWM\_CMPR1\_TEB\_INT\_CLR] 写 1 清除 \hyperref[sec:mcpwm-interrupts]{MCPWM\_CMPR1\_TEB\_INT}。 (WT)
\label{fielddesc:MCPWMCMPR2TEBINTCLR}\item [MCPWM\_CMPR2\_TEB\_INT\_CLR] 写 1 清除 \hyperref[sec:mcpwm-interrupts]{MCPWM\_CMPR2\_TEB\_INT}。 (WT)
\label{fielddesc:MCPWMTZ0CBCINTCLR}\item [MCPWM\_TZ0\_CBC\_INT\_CLR] 写 1 清除 \hyperref[sec:mcpwm-interrupts]{MCPWM\_TZ0\_CBC\_INT}。 (WT)
\label{fielddesc:MCPWMTZ1CBCINTCLR}\item [MCPWM\_TZ1\_CBC\_INT\_CLR] 写 1 清除 \hyperref[sec:mcpwm-interrupts]{MCPWM\_TZ1\_CBC\_INT}。 (WT)
\label{fielddesc:MCPWMTZ2CBCINTCLR}\item [MCPWM\_TZ2\_CBC\_INT\_CLR] 写 1 清除 \hyperref[sec:mcpwm-interrupts]{MCPWM\_TZ2\_CBC\_INT}。 (WT)
\label{fielddesc:MCPWMTZ0OSTINTCLR}\item [MCPWM\_TZ0\_OST\_INT\_CLR] 写 1 清除 \hyperref[sec:mcpwm-interrupts]{MCPWM\_TZ0\_OST\_INT}。 (WT)
\label{fielddesc:MCPWMTZ1OSTINTCLR}\item [MCPWM\_TZ1\_OST\_INT\_CLR] 写 1 清除 \hyperref[sec:mcpwm-interrupts]{MCPWM\_TZ1\_OST\_INT}。 (WT)
\label{fielddesc:MCPWMTZ2OSTINTCLR}\item [MCPWM\_TZ2\_OST\_INT\_CLR] 写 1 清除 \hyperref[sec:mcpwm-interrupts]{MCPWM\_TZ2\_OST\_INT}。 (WT)
\label{fielddesc:MCPWMCAP0INTCLR}\item [MCPWM\_CAP0\_INT\_CLR] 写 1 清除 \hyperref[sec:mcpwm-interrupts]{MCPWM\_CAP0\_INT}。 (WT)
\label{fielddesc:MCPWMCAP1INTCLR}\item [MCPWM\_CAP1\_INT\_CLR] 写 1 清除 \hyperref[sec:mcpwm-interrupts]{MCPWM\_CAP1\_INT}。 (WT)
\label{fielddesc:MCPWMCAP2INTCLR}\item [MCPWM\_CAP2\_INT\_CLR] 写 1 清除 \hyperref[sec:mcpwm-interrupts]{MCPWM\_CAP2\_INT}。 (WT)
\end{reglist}\end{regdesc}
\end{register}


\begin{register}{H}{MCPWM\_EVT\_EN\_REG}{0x0120}\label{regdesc:MCPWMEVTENREG}
\regfield{(reserved)}{2}{30}{00}%
\regfield{MCPWM\_EVT\_CAP2\_EN}{1}{29}{{0}}%
\regfield{MCPWM\_EVT\_CAP1\_EN}{1}{28}{{0}}%
\regfield{MCPWM\_EVT\_CAP0\_EN}{1}{27}{{0}}%
\regfield{MCPWM\_EVT\_TZ2\_OST\_EN}{1}{26}{{0}}%
\regfield{MCPWM\_EVT\_TZ1\_OST\_EN}{1}{25}{{0}}%
\regfield{MCPWM\_EVT\_TZ0\_OST\_EN}{1}{24}{{0}}%
\regfield{MCPWM\_EVT\_TZ2\_CBC\_EN}{1}{23}{{0}}%
\regfield{MCPWM\_EVT\_TZ1\_CBC\_EN}{1}{22}{{0}}%
\regfield{MCPWM\_EVT\_TZ0\_CBC\_EN}{1}{21}{{0}}%
\regfield{MCPWM\_EVT\_F2\_CLR\_EN}{1}{20}{{0}}%
\regfield{MCPWM\_EVT\_F1\_CLR\_EN}{1}{19}{{0}}%
\regfield{MCPWM\_EVT\_F0\_CLR\_EN}{1}{18}{{0}}%
\regfield{MCPWM\_EVT\_F2\_EN}{1}{17}{{0}}%
\regfield{MCPWM\_EVT\_F1\_EN}{1}{16}{{0}}%
\regfield{MCPWM\_EVT\_F0\_EN}{1}{15}{{0}}%
\regfield{MCPWM\_EVT\_OP2\_TEB\_EN}{1}{14}{{0}}%
\regfield{MCPWM\_EVT\_OP1\_TEB\_EN}{1}{13}{{0}}%
\regfield{MCPWM\_EVT\_OP0\_TEB\_EN}{1}{12}{{0}}%
\regfield{MCPWM\_EVT\_OP2\_TEA\_EN}{1}{11}{{0}}%
\regfield{MCPWM\_EVT\_OP1\_TEA\_EN}{1}{10}{{0}}%
\regfield{MCPWM\_EVT\_OP0\_TEA\_EN}{1}{9}{{0}}%
\regfield{MCPWM\_EVT\_TIMER2\_TEP\_EN}{1}{8}{{0}}%
\regfield{MCPWM\_EVT\_TIMER1\_TEP\_EN}{1}{7}{{0}}%
\regfield{MCPWM\_EVT\_TIMER0\_TEP\_EN}{1}{6}{{0}}%
\regfield{MCPWM\_EVT\_TIMER2\_TEZ\_EN}{1}{5}{{0}}%
\regfield{MCPWM\_EVT\_TIMER1\_TEZ\_EN}{1}{4}{{0}}%
\regfield{MCPWM\_EVT\_TIMER0\_TEZ\_EN}{1}{3}{{0}}%
\regfield{MCPWM\_EVT\_TIMER2\_STOP\_EN}{1}{2}{{0}}%
\regfield{MCPWM\_EVT\_TIMER1\_STOP\_EN}{1}{1}{{0}}%
\regfield{MCPWM\_EVT\_TIMER0\_STOP\_EN}{1}{0}{{0}}%
\reglabel{Reset}\regnewline%
\begin{regdesc}\begin{reglist}
\label{fielddesc:MCPWMEVTTIMER0STOPEN}\item [MCPWM\_EVT\_TIMER0\_STOP\_EN] 配置是否使能当定时器 0 停止时的事件生成。\\0：关闭 \\1：使能 \\(R/W)
\label{fielddesc:MCPWMEVTTIMER1STOPEN}\item [MCPWM\_EVT\_TIMER1\_STOP\_EN] 配置是否使能当定时器 1 停止时的事件生成。\\0：关闭 \\1：使能 \\(R/W)
\label{fielddesc:MCPWMEVTTIMER2STOPEN}\item [MCPWM\_EVT\_TIMER2\_STOP\_EN] 配置是否使能当定时器 2 停止时的事件生成。\\0：关闭 \\1：使能 \\(R/W)
\label{fielddesc:MCPWMEVTTIMER0TEZEN}\item [MCPWM\_EVT\_TIMER0\_TEZ\_EN] 配置是否使能当定时器 0 等于零时的事件生成。\\0：关闭 \\1：使能 \\(R/W)
\label{fielddesc:MCPWMEVTTIMER1TEZEN}\item [MCPWM\_EVT\_TIMER1\_TEZ\_EN] 配置是否使能当定时器 1 等于零时的事件生成。\\0：关闭 \\1：使能 \\(R/W)
\label{fielddesc:MCPWMEVTTIMER2TEZEN}\item [MCPWM\_EVT\_TIMER2\_TEZ\_EN] 配置是否使能当定时器 2 等于零时的事件生成。\\0：关闭 \\1：使能 \\(R/W)
\label{fielddesc:MCPWMEVTTIMER0TEPEN}\item [MCPWM\_EVT\_TIMER0\_TEP\_EN] 配置是否使能当定时器 0 等于周期时的事件生成。\\0：关闭 \\1：使能 \\(R/W)
\label{fielddesc:MCPWMEVTTIMER1TEPEN}\item [MCPWM\_EVT\_TIMER1\_TEP\_EN] 配置是否使能当定时器 1 等于周期时的事件生成。\\0：关闭 \\1：使能 \\(R/W)
\item [见下页...]
\end{reglist}\end{regdesc}
\end{register}
\addtocounter{Regfloat}{-1}
\begin{register}{H}{MCPWM\_EVT\_EN\_REG}{0x0120}
\begin{regdesc}\begin{reglist}
\item [...接上页]
\label{fielddesc:MCPWMEVTTIMER2TEPEN}\item [MCPWM\_EVT\_TIMER2\_TEP\_EN] 配置是否使能当定时器 2 等于周期时的事件生成。\\0：关闭 \\1：使能 \\(R/W)
\label{fielddesc:MCPWMEVTOP0TEAEN}\item [MCPWM\_EVT\_OP0\_TEA\_EN] 配置是否使能当生成器 0 等于时间戳 A 时的事件生成。\\0：关闭 \\1：使能 \\(R/W)
\label{fielddesc:MCPWMEVTOP1TEAEN}\item [MCPWM\_EVT\_OP1\_TEA\_EN] 配置是否使能当生成器 1 等于时间戳 A 时的事件生成。\\0：关闭 \\1：使能 \\(R/W)
\label{fielddesc:MCPWMEVTOP2TEAEN}\item [MCPWM\_EVT\_OP2\_TEA\_EN] 配置是否使能当生成器 2 等于时间戳 A 时的事件生成。\\0：关闭 \\1：使能 \\(R/W)
\label{fielddesc:MCPWMEVTOP0TEBEN}\item [MCPWM\_EVT\_OP0\_TEB\_EN] 配置是否使能当生成器 0 等于时间戳 B 时的事件生成。\\0：关闭 \\1：使能 \\(R/W)
\label{fielddesc:MCPWMEVTOP1TEBEN}\item [MCPWM\_EVT\_OP1\_TEB\_EN] 配置是否使能当生成器 1 等于时间戳 B 时的事件生成。\\0：关闭 \\1：使能 \\(R/W)
\label{fielddesc:MCPWMEVTOP2TEBEN}\item [MCPWM\_EVT\_OP2\_TEB\_EN] 配置是否使能当生成器 2 等于时间戳 B 时的事件生成。\\0：关闭 \\1：使能 \\(R/W)
\label{fielddesc:MCPWMEVTF0EN}\item [MCPWM\_EVT\_F0\_EN] 配置是否使能 FAULT0 事件生成。\\0：关闭 \\1：使能 \\(R/W)
\label{fielddesc:MCPWMEVTF1EN}\item [MCPWM\_EVT\_F1\_EN] 配置是否使能 FAULT1 事件生成。\\0：关闭 \\1：使能 \\(R/W)
\item [见下页...]
\end{reglist}\end{regdesc}
\end{register}
\addtocounter{Regfloat}{-1}
\begin{register}{H}{MCPWM\_EVT\_EN\_REG}{0x0120}
\begin{regdesc}\begin{reglist}
\item [...接上页]
\label{fielddesc:MCPWMEVTF2EN}\item [MCPWM\_EVT\_F2\_EN] 配置是否使能 FAULT2 事件生成。\\0：关闭 \\1：使能 \\(R/W)
\label{fielddesc:MCPWMEVTF0CLREN}\item [MCPWM\_EVT\_F0\_CLR\_EN] 配置是否使能 FAULT0 清除事件生成。\\0：关闭 \\1：使能 \\(R/W)
\label{fielddesc:MCPWMEVTF1CLREN}\item [MCPWM\_EVT\_F1\_CLR\_EN] 配置是否使能 FAULT0 清除事件生成。\\0：关闭 \\1：使能 \\(R/W)
\label{fielddesc:MCPWMEVTF2CLREN}\item [MCPWM\_EVT\_F2\_CLR\_EN] 配置是否使能 FAULT0 清除事件生成。\\0：关闭 \\1：使能 \\(R/W)
\label{fielddesc:MCPWMEVTTZ0CBCEN}\item [MCPWM\_EVT\_TZ0\_CBC\_EN] 配置是否使能 CBC0 事件生成。\\0：关闭 \\1：使能 \\(R/W)
\label{fielddesc:MCPWMEVTTZ1CBCEN}\item [MCPWM\_EVT\_TZ1\_CBC\_EN] 配置是否使能 CBC1 事件生成。\\0：关闭 \\1：使能 \\(R/W)
\label{fielddesc:MCPWMEVTTZ2CBCEN}\item [MCPWM\_EVT\_TZ2\_CBC\_EN] 配置是否使能 CBC2 事件生成。\\0：关闭 \\1：使能 \\(R/W)
\label{fielddesc:MCPWMEVTTZ0OSTEN}\item [MCPWM\_EVT\_TZ0\_OST\_EN] 配置是否使能 OST0 事件生成。\\0：关闭 \\1：使能 \\(R/W)
\label{fielddesc:MCPWMEVTTZ1OSTEN}\item [MCPWM\_EVT\_TZ1\_OST\_EN] 配置是否使能 OST1 事件生成。\\0：关闭 \\1：使能 \\(R/W)
\item [见下页...]
\end{reglist}\end{regdesc}
\end{register}
\addtocounter{Regfloat}{-1}
\begin{register}{H}{MCPWM\_EVT\_EN\_REG}{0x0120}
\begin{regdesc}\begin{reglist}
\item [...接上页]
\label{fielddesc:MCPWMEVTTZ2OSTEN}\item [MCPWM\_EVT\_TZ2\_OST\_EN] 配置是否使能 OST2 事件生成。\\0：关闭 \\1：使能 \\(R/W)
\label{fielddesc:MCPWMEVTCAP0EN}\item [MCPWM\_EVT\_CAP0\_EN] 配置是否使能 capture0 事件生成。\\0：关闭 \\1：使能 \\(R/W)
\label{fielddesc:MCPWMEVTCAP1EN}\item [MCPWM\_EVT\_CAP1\_EN] 配置是否使能 capture1 事件生成。\\0：关闭 \\1：使能 \\(R/W)
\label{fielddesc:MCPWMEVTCAP2EN}\item [MCPWM\_EVT\_CAP2\_EN] 配置是否使能 capture2 事件生成。\\0：关闭 \\1：使能 \\(R/W)
\end{reglist}\end{regdesc}
\end{register}


\begin{register}{H}{MCPWM\_TASK\_EN\_REG}{0x0124}\label{regdesc:MCPWMTASKENREG}
\regfield{(reserved)}{10}{22}{0000000000}%
\regfield{MCPWM\_TASK\_CAP2\_EN}{1}{21}{{0}}%
\regfield{MCPWM\_TASK\_CAP1\_EN}{1}{20}{{0}}%
\regfield{MCPWM\_TASK\_CAP0\_EN}{1}{19}{{0}}%
\regfield{MCPWM\_TASK\_CLR2\_OST\_EN}{1}{18}{{0}}%
\regfield{MCPWM\_TASK\_CLR1\_OST\_EN}{1}{17}{{0}}%
\regfield{MCPWM\_TASK\_CLR0\_OST\_EN}{1}{16}{{0}}%
\regfield{MCPWM\_TASK\_TZ2\_OST\_EN}{1}{15}{{0}}%
\regfield{MCPWM\_TASK\_TZ1\_OST\_EN}{1}{14}{{0}}%
\regfield{MCPWM\_TASK\_TZ0\_OST\_EN}{1}{13}{{0}}%
\regfield{MCPWM\_TASK\_TIMER2\_PERIOD\_UP\_EN}{1}{12}{{0}}%
\regfield{MCPWM\_TASK\_TIMER1\_PERIOD\_UP\_EN}{1}{11}{{0}}%
\regfield{MCPWM\_TASK\_TIMER0\_PERIOD\_UP\_EN}{1}{10}{{0}}%
\regfield{MCPWM\_TASK\_TIMER2\_SYNC\_EN}{1}{9}{{0}}%
\regfield{MCPWM\_TASK\_TIMER1\_SYNC\_EN}{1}{8}{{0}}%
\regfield{MCPWM\_TASK\_TIMER0\_SYNC\_EN}{1}{7}{{0}}%
\regfield{MCPWM\_TASK\_GEN\_STOP\_EN}{1}{6}{{0}}%
\regfield{MCPWM\_TASK\_CMPR2\_B\_UP\_EN}{1}{5}{{0}}%
\regfield{MCPWM\_TASK\_CMPR1\_B\_UP\_EN}{1}{4}{{0}}%
\regfield{MCPWM\_TASK\_CMPR0\_B\_UP\_EN}{1}{3}{{0}}%
\regfield{MCPWM\_TASK\_CMPR2\_A\_UP\_EN}{1}{2}{{0}}%
\regfield{MCPWM\_TASK\_CMPR1\_A\_UP\_EN}{1}{1}{{0}}%
\regfield{MCPWM\_TASK\_CMPR0\_A\_UP\_EN}{1}{0}{{0}}%
\reglabel{Reset}\regnewline%
\begin{regdesc}\begin{reglist}
\label{fielddesc:MCPWMTASKCMPR0AUPEN}\item [MCPWM\_TASK\_CMPR0\_A\_UP\_EN] 配置是否接收 PWM 生成器 0 时间戳 A 的影子寄存器的更新任务。 \\0：无效\\1：接收\\ (R/W)
\label{fielddesc:MCPWMTASKCMPR1AUPEN}\item [MCPWM\_TASK\_CMPR1\_A\_UP\_EN] 配置是否接收 PWM 生成器 1 时间戳 A 的影子寄存器的更新任务。 \\0：无效\\1：接收\\ (R/W)
\label{fielddesc:MCPWMTASKCMPR2AUPEN}\item [MCPWM\_TASK\_CMPR2\_A\_UP\_EN] 配置是否接收 PWM 生成器 2 时间戳 A 的影子寄存器的更新任务。 \\0：无效\\1：接收\\ (R/W)
\label{fielddesc:MCPWMTASKCMPR0BUPEN}\item [MCPWM\_TASK\_CMPR0\_B\_UP\_EN] 配置是否接收 PWM 生成器 0 时间戳 B 的影子寄存器的更新任务。 \\0：无效\\1：接收\\ (R/W)
\label{fielddesc:MCPWMTASKCMPR1BUPEN}\item [MCPWM\_TASK\_CMPR1\_B\_UP\_EN] 配置是否接收 PWM 生成器 1 时间戳 B 的影子寄存器的更新任务。 \\0：无效\\1：接收\\ (R/W)
\label{fielddesc:MCPWMTASKCMPR2BUPEN}\item [MCPWM\_TASK\_CMPR2\_B\_UP\_EN] 配置是否接收 PWM 生成器 2 时间戳 B 的影子寄存器的更新任务。 \\0：无效\\1：接收\\ (R/W)
\item [见下页...]
\end{reglist}\end{regdesc}
\end{register}
\addtocounter{Regfloat}{-1}
\begin{register}{H}{MCPWM\_TASK\_EN\_REG}{0x0124}
\begin{regdesc}\begin{reglist}
\item [...接上页]
\label{fielddesc:MCPWMTASKGENSTOPEN}\item [MCPWM\_TASK\_GEN\_STOP\_EN] 配置是否接收所有 PWM 生成器的停止任务。 \\0：无效\\1：接收\\ (R/W)
\label{fielddesc:MCPWMTASKTIMER0SYNCEN}\item [MCPWM\_TASK\_TIMER0\_SYNC\_EN] 配置是否接收定时器 0 的同步任务。 \\0：无效\\1：接收\\ (R/W)
\label{fielddesc:MCPWMTASKTIMER1SYNCEN}\item [MCPWM\_TASK\_TIMER1\_SYNC\_EN] 配置是否接收定时器 1 的同步任务。 \\0：无效\\1：接收\\ (R/W)
\label{fielddesc:MCPWMTASKTIMER2SYNCEN}\item [MCPWM\_TASK\_TIMER2\_SYNC\_EN] 配置是否接收定时器 2 的同步任务。 \\0：无效\\1：接收\\ (R/W)
\label{fielddesc:MCPWMTASKTIMER0PERIODUPEN}\item [MCPWM\_TASK\_TIMER0\_PERIOD\_UP\_EN] 配置是否接收定时器 0 的周期更新任务。 \\0：无效\\1：接收\\ (R/W)
\label{fielddesc:MCPWMTASKTIMER1PERIODUPEN}\item [MCPWM\_TASK\_TIMER1\_PERIOD\_UP\_EN] 配置是否接收定时器 1 的周期更新任务。 \\0：无效\\1：接收\\ (R/W)
\label{fielddesc:MCPWMTASKTIMER2PERIODUPEN}\item [MCPWM\_TASK\_TIMER2\_PERIOD\_UP\_EN] 配置是否接收定时器 2 的周期更新任务。 \\0：无效\\1：接收\\ (R/W)
\label{fielddesc:MCPWMTASKTZ0OSTEN}\item [MCPWM\_TASK\_TZ0\_OST\_EN] 配置是否接收 OST0 任务。 \\0：无效\\1：接收\\ (R/W)
\label{fielddesc:MCPWMTASKTZ1OSTEN}\item [MCPWM\_TASK\_TZ1\_OST\_EN] 配置是否接收 OST1 任务。 \\0：无效\\1：接收\\ (R/W)
\item [见下页...]
\end{reglist}\end{regdesc}
\end{register}
\addtocounter{Regfloat}{-1}
\begin{register}{H}{MCPWM\_TASK\_EN\_REG}{0x0124}
\begin{regdesc}\begin{reglist}
\item [...接上页]
\label{fielddesc:MCPWMTASKTZ2OSTEN}\item [MCPWM\_TASK\_TZ2\_OST\_EN] 配置是否接收 OST2 任务。 \\0：无效\\1：接收\\ (R/W)
\label{fielddesc:MCPWMTASKCLR0OSTEN}\item [MCPWM\_TASK\_CLR0\_OST\_EN] 配置是否接收 OST0 清除任务。 \\0：无效\\1：接收\\ (R/W)
\label{fielddesc:MCPWMTASKCLR1OSTEN}\item [MCPWM\_TASK\_CLR1\_OST\_EN] 配置是否接收 OST1 清除任务。 \\0：无效\\1：接收\\ (R/W)
\label{fielddesc:MCPWMTASKCLR2OSTEN}\item [MCPWM\_TASK\_CLR2\_OST\_EN] 配置是否接收 OST2 清除任务。 \\0：无效\\1：接收\\ (R/W)
\label{fielddesc:MCPWMTASKCAP0EN}\item [MCPWM\_TASK\_CAP0\_EN] 配置是否接收 capture0 任务。 \\0：无效\\1：接收\\ (R/W)
\label{fielddesc:MCPWMTASKCAP1EN}\item [MCPWM\_TASK\_CAP1\_EN] 配置是否接收 capture1 任务。 \\0：无效\\1：接收\\ (R/W)
\label{fielddesc:MCPWMTASKCAP2EN}\item [MCPWM\_TASK\_CAP2\_EN] 配置是否接收 capture2 任务。 \\0：无效\\1：接收\\ (R/W)
\end{reglist}\end{regdesc}
\end{register}


\begin{register}{H}{MCPWM\_CLK\_REG}{0x0128}\label{regdesc:MCPWMCLKREG}
\regfield{(reserved)}{31}{1}{0000000000000000000000000000000}%
\regfield{MCPWM\_CLK\_EN}{1}{0}{{0}}%
\reglabel{Reset}\regnewline%
\begin{regdesc}\begin{reglist}
\label{fielddesc:MCPWMCLKEN}\item [MCPWM\_CLK\_EN] 配置是否强制使能时钟。\\
0：无效 \\
1：强制使能时钟 \\
(R/W)
\end{reglist}\end{regdesc}
\end{register}


\begin{register}{H}{MCPWM\_VERSION\_REG}{0x012C}\label{regdesc:MCPWMVERSIONREG}
\regfield{(reserved)}{4}{28}{0000}%
\regfield{MCPWM\_DATE}{28}{0}{{0x2107230}}%
\reglabel{Reset}\regnewline%
\begin{regdesc}\begin{reglist}
\label{fielddesc:MCPWMDATE}\item [MCPWM\_DATE] 版本控制寄存器。 (R/W)
\end{reglist}\end{regdesc}
\end{register}


